\documentclass[11pt]{article}
\usepackage[a4paper,margin=27mm]{geometry}
\usepackage[T1]{fontenc}
\usepackage{lmodern,microtype}
\usepackage{amsmath,amssymb,amsthm,mathtools}
\usepackage{enumitem,booktabs,array,longtable}
\usepackage[hidelinks,hypertexnames=false]{hyperref}
\hypersetup{pdftitle={The Contamination Triangle}}

\newtheorem{theorem}{Theorem}[section]
\newtheorem{lemma}[theorem]{Lemma}
\newtheorem{proposition}[theorem]{Proposition}
\newtheorem{corollary}[theorem]{Corollary}
\newtheorem{conjecture}[theorem]{Conjecture}
\theoremstyle{definition}
\newtheorem{definition}[theorem]{Definition}
\theoremstyle{remark}
\newtheorem{remark}[theorem]{Remark}

\newcommand{\R}{\mathbb R}
\newcommand{\one}{\mathbf 1}
\newcommand{\Di}{\Delta^{\circ}}
\newcommand{\M}{\mathcal M}
\newcommand{\C}{\mathcal C}
\newcommand{\cl}{\operatorname{cl}}
\newcommand{\co}{\operatorname{conv}}
\newcommand{\supp}{\operatorname{supp}}
\newcommand{\relint}{\operatorname{relint}}
\newcommand{\aff}{\operatorname{aff}}
\newcommand{\lin}{\operatorname{line}}
\newcommand{\sd}{\mathbin{\triangle}}
\newcommand{\lp}{\underline{P}}
\newcommand{\up}{\overline{P}}
\newcommand{\dH}{d_{\mathrm H}}
\newcommand{\src}[1]{\texttt{#1}}
\newcommand{\status}[1]{\par\smallskip\noindent\textbf{Status:} #1\par\smallskip}
\newcommand{\Ltwo}{\mathcal L_2^{\mathrm s}}
\newcommand{\Nb}{N_0^{\mathrm b}}

\title{The Contamination Triangle}
\author{Research note on coverage under decomposable domains}
\date{14 October 2026 (round 9)}

\begin{document}
\emergencystretch=2em
\maketitle

\section{Introduction}\label{sec:intro}

\paragraph{The coverage question.}
A dominance argument for coherence has two halves. The first shows that, for any ranking of models that satisfies certain axioms, every model is weakly dominated by an undominated coherent model. For finitely many worlds this is the KKM theorem of \src{Loss Dominance Generalized v2} (Theorem~\ref{thm:kkm} below). The second half, \emph{coverage}, asks that every coherent model be undominated for at least one ranking that satisfies the axioms. Without coverage the argument shows only that the permissible models lie among the coherent ones. It does not show that every coherent model is permissible.

In the setting adopted here, models are sets of desirable gambles. A ranking compares models by the advice they give on binary decision problems. At each world, good advice is better than silence, silence is better than bad advice, and bad advice is better than contradictory advice (guidance, (G)). Comparisons depend only on the problems on which two models differ (independence, (IB)). Rankings are preorders, which may be incomplete and need not be numerical. Each ranking evaluates the models in a \emph{decomposable domain}, its class of evaluable models. Restricting the domain restricts only which models are evaluated; the decision problems are never altered or removed. To stop a target being protected merely by excluding its borderline competitors, the domain must contain every model that differs from the target only near the boundary of its cone of desirable gambles (condition (LF)).

\paragraph{Why the contamination triangle.}
Coverage is known for precise models, for every coherent model when there are two worlds, for normalised boxes, and for a large class of interior polygons (Section~\ref{sec:positive}). In three worlds the simplest imprecise models not known to be covered are those whose credal set is a triangle in the interior of the simplex. The most natural of these is the $\varepsilon$-contamination of the uniform prior,
\[
T_1\coloneqq(1-\varepsilon)u+\varepsilon\Delta,\qquad u=(\tfrac13,\tfrac13,\tfrac13),\quad 0<\varepsilon<1 .
\]
Its cone of desirable gambles has three flat facets that meet along three corner rays. Additive rankings built by mixing gauges cover hexagons and 9-gons arbitrarily close to $T_1$, but provably no triangle. In earlier work on coherence-preserving maps the same target was shown to have no continuous realiser. So $T_1$ is where the structure of the axioms, and not the convenience of a construction, decides coverage.

\paragraph{What is shown.}
\begin{enumerate}[leftmargin=*]
\item \emph{Additive rankings} (Sections~\ref{sec:reduction}--\ref{sec:main}). If an additive ranking on the $U$-domain (the class of measurable subsets of a fixed open cone $U$ around $\cl D_{T_1}\setminus\{0\}$) covers $T_1$, the Bayes models of its interior weightings form a Hausdorff-continuous family of credal sets with half-plane charts that passes through $T_1$. In three worlds the converse holds too. The main theorem (Theorem~\ref{thm:main}, Corollary~\ref{cor:coverage}) is that on every decomposable domain whose pairs near $\cl D_{T_1}$ (those with an orientation in a cone $D_{K'}$, $K'\subseteq\relint T_1$) are non-adjacent, and in particular on the $U$-domain, no additive ranking covers $T_1$. This holds for every $\varepsilon\in(0,1)$, at every interior target weighting, and with no regularity assumption. The proof applies to every all-facing triangle in general position, hence to every interior $\varepsilon$-contamination.
\item \emph{General decomposable domains} (Section~\ref{sec:general}). Without non-adjacency, in every additive cover of $T_1$ the advice range of almost every mixed pair contains both acceptances, and no vertex is pinned near the target weighting; competition between the two acceptances can occur only in neighbourhoods of the vertices of $\Delta$ that shrink to the vertices. Additive coverage of $T_1$ on some (LF) decomposable domain is reduced to a purely geometric problem, CF$(T_1)$, about Hausdorff-continuous credal families. CF$(T_1)$ is open.
\item \emph{Coherent almost-preference rankings} (Section~\ref{sec:cp}). If each world's ranking is a coherent almost-preference in Walley's sense, coverage of $T_1$ on a non-adjacent domain is impossible when every world is complete (on the $U$-domain, or under a mild domain condition (UN)), and when strict comparisons are detected by countably additive previsions. In general, coverage is reduced to one property (U) of a single integral field.
\item \emph{Rankings that add up} (\S\S\ref{sec:M0}--\ref{sec:M0-lex}). If every world of an arbitrary ranking on the $U$-domain adds up on the relevant comparisons, a supporting additive assessment exists whose exceptional weightings form a nonempty countable union of chords of $\Delta$, each cutting off one vertex. At those weightings coherent and incoherent minimisers coexist and the assessment cannot tell them apart; a cover requires the original ranking to do so. A two-world construction shows that a ranking can resolve such a tie while keeping every axiom; in three worlds, for two-level lexicographic rankings, the tie cannot sit at the target weighting.
\item \emph{Weakened guidance} (\S\ref{sec:open-ties}). The theorem of Section~\ref{sec:main} survives permitted ties on problems with small stakes, for every band width in the sup norm and for Euclidean band widths below an explicit threshold.
\item Section~\ref{sec:open} states the two concrete open questions (\S\ref{sec:two-questions}) and lists what is excluded and what is open; Section~\ref{sec:positive} recalls the positive results.
\end{enumerate}

\paragraph{Status of the results.}
Every result below carries a status label right after its statement. ``Refereed'' or ``confirmed'' means that a reader who did not write the note re-derived each step and returned the verdict CORRECT, CONFIRMED, CORRECT WITH REPAIR or CONFIRMED WITH REPAIR. All requested repairs are incorporated in the statements given here; where a referee's repair conflicts with the note, the referee's version is used. A result first proved \emph{in} a referee report has been checked by nobody else unless a second reading is recorded. CHECKED means verified with a proof sketch; NUMERICS means supported only by computation. Notes are cited by their file names (for example \src{N15} for \src{N15\_notes.md}); the Sources section at the end lists them.

{\small
\begin{longtable}{>{\raggedright\arraybackslash}p{6.4cm}>{\raggedright\arraybackslash}p{2.1cm}>{\raggedright\arraybackslash}p{5.6cm}}
\toprule
Result & Where & Status\\
\midrule
\endhead
Additive cover $\Rightarrow$ unique coherent interior Bayes family through the target & Prop.~\ref{prop:necessary} & proved; refereed once (R6)\\
Reduction to a credal family with half-plane charts & Prop.~\ref{prop:reduction} & proved; refereed once (N14)\\
Closed choice from interior Bayes uniqueness ($n=3$) & Thm.~\ref{thm:ccibu} & author-checked; re-derived without error in N13; no formal referee report\\
\textbf{No additive cover of $T_1$ (all-facing triangles) on non-adjacent domains} & Thm.~\ref{thm:main}, Cor.~\ref{cor:coverage} & \textbf{proved; refereed independently twice (N18, N20)}; imported lemmas refereed once each (N14, R12, R18)\\
Same with bounded charge penalties, given interior minimisers & Cor.~\ref{cor:charges} & proved (N14 Thm C2, in a referee report; splice step refereed in N24); use refereed (N18, N20)\\
General decomposable domains: compact collar replaces the conic $N_0$; ray transfer of charts & \S\ref{sec:general} preamble, Lem.~\ref{lem:raytransfer} & proved (N27); confirmed by one referee (N27 referee)\\
General decomposable domains: (TS) $\Rightarrow$ structure; dichotomy & Thms.~\ref{thm:TS}, \ref{thm:dichotomy} & proved; refereed once (N23); branch (II) also uses Cor.~\ref{cor:R232}\\
Lemma R23.1, Cor.\ R23.2, Cor.\ R23.4 & Lem.~\ref{lem:R231}, Cors.~\ref{cor:R232}, \ref{cor:R234} & proved in a referee report (N23); second reading CORRECT (CHECK\_REPORT.md)\\
\textbf{(TS) holds in every cover; (E5a) is empty} & Thm.~\ref{thm:TSalways}, Cor.~\ref{cor:TSalways} & \textbf{proved (N28); confirmed by one referee (N28 referee)}; rests on once-refereed inputs\\
Cut apexes at the vertex; vertex-free interiors; non-precise credal sets on crossing chords & Thm.~\ref{thm:apex}, Prop.~\ref{prop:vertexfree} & proved (N28); confirmed by one referee\\
Global one-sidedness of chords & Lem.~\ref{lem:R1N28} & proved in a referee report (N28 referee), one reader\\
Reduction to CF$(T_1)$ & Prop.~\ref{prop:CF} & CHECKED (N28); confirmed with repair (N28 referee); converse at the status of Thm.~\ref{thm:ccibu}\\
Diagonal symmetries & Prop.~\ref{prop:diag} & proved (N28), confirmed; ranking-level form proved in a referee report (N28 referee), one reader\\
CF$(T_1)$ (residual of (E4$'$)) & \S\ref{sec:general-open} & open; a sector-cap example passes every derived constraint (NUMERICS)\\
(CP-a): full support of countably additive parts is automatic & Thm.~\ref{thm:fullsupport} & proved; refereed once (N24)\\
(CP-a): integral reduction under (CP-a) alone & Thm.~\ref{thm:integral} & proved; refereed once (N24)\\
\textbf{(CP-a) with complete worlds: no cover of $T_1$} & Thm.~\ref{thm:complete}, Cor.~\ref{cor:complete} & \textbf{proved; refereed once (N24)}\\
(CP-a) with countably additive detection: no cover & Cor.~\ref{cor:DCA} & proved; refereed once (N13, N24)\\
General (CP-a): cover forces a tie weight; (U) as continuity & Cor.~\ref{cor:residual}, Prop.~\ref{prop:Ucont} & residual proved, refereed once (N24); continuity form proved by the referee, one reader\\
(U) for every (CP-a) system & Conj.~\ref{conj:U} & open\\
M0: supporting assessment with properties ($\alpha$)--($\delta$) & Prop.~\ref{prop:M0} & proved; inputs refereed once or proved in referee reports (N13, N19, N24)\\
Structure of the exceptional set: countably many atom chords & Thm.~\ref{thm:Estructure} & proved (N25), confirmed by one referee; ``no fan points'' proved in a referee report (N25 referee), one reader\\
$p^*\in E(\psi)$ blocks the chain of \S\ref{sec:main}; genericity & Thm.~\ref{thm:pstarE}, Prop.~\ref{prop:generic} & proved (N25); confirmed with repair (N25 referee)\\
Consequences of (E$_\psi$): ideal tie-breaking, transport along chords & Prop.~\ref{prop:Epsi} & proved (N25); confirmed by one referee\\
Residuals R-a$'$, R-b$'$ of M0 & \S\ref{sec:M0-residual} & open\\
\textbf{Two-world ranking resolving an exceptional additive tie} & Thm.~\ref{thm:N26.1} & \textbf{proved (N26); confirmed by one referee (N26 referee)}\\
Two-level lexicographic class $\Ltwo$: no 2-dim.\ faces, rigid chords, $p^*\notin E$ & Thm.~\ref{thm:Ltwo} & proved (N26); confirmed with repairs (N26 referee); does not transfer to general M0 rankings\\
$\Ltwo$: the pivot theorem alone removes the chords & Prop.~\ref{prop:R1N26} & proved in a referee report (N26 referee), one reader\\
Strict-preference variants conflict with (S) & Prop.~\ref{prop:A} & proved in a referee report (N13); one reader\\
Walley-2000 variant with invisible desirability & \S\ref{sec:cp-b} & open\\
Forced small-stakes ties: no additive cover & \S\ref{sec:open-ties} & proved; refereed once (N4)\\
Permitted small-stakes ties: Thm.~\ref{thm:main} survives below $\theta_x(T)$ & Thm.~\ref{thm:permitted} & proved (N27); confirmed by one referee (N27 referee)\\
Permitted ties, non-adjacent (LF) domains & Cor.~\ref{cor:permitted-nonadj} & CHECKED (N27); confirmed at CHECKED level\\
Euclidean band $\theta\in[\theta_x(T_1),1)$ & \S\ref{sec:open-ties} & open\\
Other triangle classes, $n\ge4$, segments, boundary targets & \S\ref{sec:open-other} & open\\
\bottomrule
\end{longtable}}

\section{Setting}\label{sec:setting}

\subsection{Problems, models and advice}
Let $\Omega\coloneqq\{1,\dots,n\}$ be the worlds, $n\ge2$; from Section~\ref{sec:reduction} on, $n=3$. A gamble is $f\in\R^n$; it pays $f_i$ at world $i$. Write $\Delta$ for the probability simplex, $e_i$ for its vertices, $\one=(1,\dots,1)$, $\supp p\coloneqq\{i:p_i>0\}$, and $\Di$ for the interior of $\Delta$ relative to $\aff\Delta$. Fix the standard Gaussian measure $\mu$ on $\R^n$; it has the same null sets as Lebesgue measure.

\paragraph{What a gamble represents.} Any gamble $f$ represents the binary choice $\{f,0\}$, which is equivalent to the binary choice $\{g,h\}$ for any gambles $g,h$ with $f=g-h$. That is all that $f$ represents. The gambles $f$ and $-f$ represent the same binary choice, with the roles of the two options swapped; we call this choice the \emph{pair} $\{f,-f\}$. A \emph{model} is a Lebesgue-measurable set $D\subseteq\R^n$ of desirable (accepted) gambles, taken modulo $\mu$-null sets. On the choice represented by $f$ it gives one of four kinds of advice:
\begin{itemize}[leftmargin=*]
\item $f\in D$ and $-f\notin D$: you prefer $f=g-h$ to $0$ and do not prefer $-f=h-g$ to $0$, which is true iff you prefer $g$ to $h$ and do not prefer $h$ to $g$; so $D$ advises choosing $g$ in $\{g,h\}$, for every $g,h$ with $f=g-h$;
\item $f\notin D$ and $-f\in D$: $D$ advises choosing $h$ in $\{g,h\}$, for every $g,h$ with $f=g-h$;
\item $f\notin D$ and $-f\notin D$: no advice (also called \emph{silence});
\item $f\in D$ and $-f\in D$: contradictory advice.
\end{itemize}
As bookkeeping only, each pair is indexed once by its member in $H\coloneqq\{f:f_1>0\}$, and for $f\in H$ we write the advice $a_D(f)$ as $+$ (accept $f$ only), $-$ (accept $-f$ only), $0$ (no advice) or $\times$ (both); nothing depends on this convention. At world $i$, accepting $f$ (that is, choosing $g$ in $\{g,h\}$) is \emph{good} if $f_i>0$ and \emph{bad} if $f_i<0$, and the roles swap for $-f$. Coordinate hyperplanes are null and are ignored. The \emph{background order} at world $i$, better first, is
\[
\text{good}\;\prec\;\text{no advice}\;\prec\;\text{bad}\;\prec\;\text{contradictory}.
\]
From here on the letters $g,h,x,\dots$ denote gambles, each representing its binary choice in this way; the options of the choice no longer appear.
A gamble is \emph{mixed} if it has a positive and a negative coordinate.

For a credal set $K$ (a nonempty closed convex subset of $\Delta$) put
\[
P_K(g)\coloneqq\min_{m\in K}m\cdot g,\qquad D_K\coloneqq\{g:P_K(g)>0\}\cup\{g\ge0:g\ne0\}.
\]
Every coherent set of desirable gambles agrees a.e.\ with some $D_K$ (\src{Full Coverage by Homothetic Credal Families v3}, App.~A), and $D_K\approx D_{K'}$ modulo null sets only if $K=K'$. A model is \emph{coherent} if it agrees a.e.\ with some $D_K$.

\subsection{Decomposable domains, rankings and axioms}
The domain of a ranking is the class of models that it evaluates. Restricting the domain means only restricting which models are evaluated. Decision problems are never altered or removed, and they are never given restricted options: every gamble $f$ still represents the binary choice $\{f,0\}$, equivalently $\{g,h\}$ with $f=g-h$, and the world orders of a ranking still concern every binary problem. This is the only way to make sense of a domain restriction in the accuracy-ordering framework.

\begin{definition}\label{def:domain}
A \emph{decomposable domain} is a nonempty class $\M$ of models (the \emph{evaluable models}) that is closed under splicing, i.e.\ $(X\cap A)\cup(Y\setminus A)\in\M$ whenever $X,Y\in\M$ and $A=-A$ is measurable, and under convergence in measure, i.e.\ $X\in\M$ whenever $X_m\in\M$ and $\mu(X_m\sd X)\to0$.
\end{definition}

\begin{remark}\label{rem:advicerange}
The decomposable domains are exactly the classes of the form
\[
\M=\{D\text{ measurable}:\ a_D(f)\in A(f)\text{ for a.e.\ }f\in H\},
\]
where $f\mapsto A(f)\subseteq\{+,0,-,\times\}$ is measurable (each set $\{f:a\in A(f)\}$ is measurable) and nonempty for a.e.\ $f$; the map $A$ is determined by $\M$ up to a null set. We call $A(f)$ the \emph{advice range} of the domain at $f$: the kinds of advice that evaluable models give on the binary choice represented by $f$. So ``the advice range at $f$ is $\{+,0\}$'' says that every evaluable model either accepts $f$ (and not $-f$) or gives no advice on that choice, and that evaluable models of both kinds exist; and when the advice range at $f$ is a single kind of advice, every evaluable model gives that advice on the choice represented by $f$. That choice is still a decision problem; it simply does not distinguish evaluable models.
\end{remark}
\status{proved (\src{Domain Closure and Extendability}, Prop.~2, author-checked): take a countable dense family $(X_k)$ in $\M$ and put $A(f)\coloneqq\{a_{X_k}(f):k\ge1\}$; every evaluable model has its advice in $A$ a.e., and every measurable selection from $A$ is a limit in measure of finite splices of the $X_k$. Splicing alone is not enough (same source, first part of Prop.~2). This is a special case of the theorem of Hiai and Umegaki on closed decomposable subsets of $L^1$~\cite{HU1977}.}

\paragraph{Earlier terminology.} Versions of this note and of the working notes written before 14 October 2026 called decomposable domains ``menu domains'', called the advice range a ``menu'' of available actions, called the $U$-domain of \S\ref{sec:aos} the ``accept-or-silence domain'', and called problems on which every evaluable model gives no advice ``forced silent''. That language misdescribed a restriction of the class of evaluable models as a restriction of the decision problems, and it has been removed from all current files; files and folders whose names used it were renamed on the same day.

A \emph{ranking} $L$ consists of a decomposable domain $\M_L$, its class of evaluable models, and a preorder $\preceq_i$ on $\M_L$ at each world $i$ (smaller is better). The world orders compare evaluable models by the advice they give on all binary problems. Write $X\prec_iY$ for the strict part and $X\sim_iY$ for ties. $Y$ \emph{dominates} $X$ if $Y\preceq_iX$ at every world and $Y\prec_jX$ at some world. $U_L$ is the set of undominated members of $\M_L$, and $\C_L$ the set of its coherent members. A model $C$ is \emph{competitive on} $S\subseteq\Omega$ if every $X\in\M_L$ has some $i\in S$ with $C\preceq_iX$. The following requirements are imposed within $\M_L$ (\src{Full Coverage \dots v3}, \S1; \src{Loss Dominance Generalized v2}, \S\S1--2).
\begin{enumerate}[leftmargin=*,label=\textbf{(\arabic*)}]
\item[\textbf{(G)}] \emph{Guidance.} If $X$ gives advice at least as good as $Y$ at world $i$ for a.e.\ $g$, then $X\preceq_iY$; if moreover it is strictly better on a non-null set, then $X\prec_iY$.
\item[\textbf{(IB)}] \emph{Independence across binary problems.} For measurable $A=-A$ and measurable $F,G\subseteq A^c$ (with all four models in $\M_L$): $(D\cap A)\cup F\preceq_i(E\cap A)\cup F$ iff $(D\cap A)\cup G\preceq_i(E\cap A)\cup G$.
\item[\textbf{(LA)}] $\C_L\ne\varnothing$.
\item[\textbf{(S)}] \emph{Closed competitive choice.} There is a map $p\mapsto C_p\in\C_L$ on $\Delta$ such that $C_p$ is competitive on $\supp p$ and $\{p:C_p\preceq_iX\}$ is closed for every $i$ and every $X\in\M_L$.
\item[\textbf{(C)}] A model tied at every world with a member of $\C_L$ belongs to $\C_L$.
\end{enumerate}
Restricting the domain leaves (G) unchanged: it compares the advice that two evaluable models give on each problem. For a class of rankings one also asks (E): every model is evaluated by some ranking (supplied by an all-model Brier ranking), and (A): every coherent model is in some $\C_L$. Condition (A) is \emph{coverage}.

\begin{theorem}\label{thm:kkm}
Let $\Omega$ be finite and let $L$ satisfy (LA) and (S). Then every $X\in\M_L$ is weakly dominated by an undominated selected model $C_p$, and every undominated model is tied at every world with a selected model.
\end{theorem}
\status{proved (\src{Loss Dominance Generalized v2}, Thm.~thm:kkm, via the KKM lemma; valid for preorders). Cited as ``Fact 0'' in the notes.}

\begin{definition}\label{def:LF}
A ranking $L$ \emph{covers} a coherent model $D$ if $D\in\C_L$ and (LF) holds at $D$: there is an open $N\supseteq\partial D\setminus\{0\}$ such that every measurable $R$ with $R\sd D\subseteq N$ a.e.\ belongs to $\M_L$. For a decomposable domain this means that for a.e.\ pair with an orientation $h\in N$ the advice range contains both ``accept $h$'' and no advice.
\end{definition}
This is the open-set form of (LF): $N$ need not contain a conic neighbourhood of $\partial D\setminus\{0\}$ (for instance $N=\{g:\operatorname{dist}(g,\partial D)<e^{-\|g\|}\}$ contains none; \src{N27} \S1). On the $U$-domain (\S\ref{sec:aos}) the explicit open cone $U$ serves as $N$. For additive rankings whose interior Bayes models are coherent, charts are invariant along almost every ray (Lemma~\ref{lem:raytransfer}), so the conic form holds at the level of charts; no strengthening of (LF) is needed anywhere in this note (\src{N27} Repair 27.1, confirmed in \src{N27\_referee}).

\subsection{The target}
Put $c_0\coloneqq(1-\varepsilon)/3$ and
\[
\kappa\coloneqq\frac{1+2\varepsilon}{1-\varepsilon}=1+\frac{3\varepsilon}{1-\varepsilon}>1 .
\]
The vertices of $T_1$ are $v_k=c_0\one+\varepsilon e_k$, and
\[
D_{T_1}=\Bigl\{g:\ \tfrac{1-\varepsilon}3(g_1+g_2+g_3)+\varepsilon g_k>0\ \text{for }k=1,2,3\Bigr\}\cup\{g\ge0:g\ne0\}.
\]
The facets of $D_{T}$ for a triangle $T=\co\{v_1,v_2,v_3\}$ are $F_k\coloneqq\{h:h\cdot v_k=0,\ h\cdot v_j\ge0\ (j\ne k)\}$, two-dimensional closed cones in $v_k^\perp$. Two facets $F_k,F_j$ meet in a \emph{corner ray} $r_{kj}\propto v_k\times v_j$. Coordinate $i$ \emph{crosses} on $F_k$ if $h_i=0$ at some $h\in\relint F_k$. For $T_1$ the corner rays are
\[
r_{12}=(-1,-1,\kappa+1),\quad r_{23}=(\kappa+1,-1,-1),\quad r_{31}=(-1,\kappa+1,-1),
\]
and on $F_1$ the coordinates $2,3$ cross at $w_{21}=(-1,0,\kappa)$ and $w_{31}=(-1,\kappa,0)$ (cyclically on $F_2,F_3$). These follow from $c_0(\kappa-1)=\varepsilon$.

\begin{definition}\label{def:allfacing}
A triangle $T\subseteq\Di$ is in \emph{general position} if no line through two of its vertices contains a vertex of $\Delta$. It is \emph{all-facing} if $v_k\in\operatorname{int}\co\{e_k,v_i,v_j\}$ for each $k$, with $\{i,j,k\}=\{1,2,3\}$.
\end{definition}
Every interior $\varepsilon$-contamination $T=(1-\varepsilon)p_0+\varepsilon\Delta$ ($p_0\in\Di$) is all-facing and in general position: its edges are parallel to those of $\Delta$ at positive distance, and $v_k$ lies on the open segment from $p_0\in\operatorname{int}T$ to $e_k$ (\src{R15} \S6, \src{R18} \S6).

\subsection{Additive rankings and charts}
An \emph{additive ranking} has measurable penalties $\ell_i(g,a)$, $a\in A(g)$, with an integrable envelope, and ranks world $i$ by $L_i(X)\coloneqq\int_H\ell_i(g,a_X(g))\,d\mu(g)$. Only the advice that evaluable models actually give enters: the penalty is defined for the kinds of advice in the advice range $A(g)$. On a problem on which all evaluable models give the same advice, the penalty is the same for every evaluable model and drops out of every comparison; it is still part of the ranking. On decomposable domains, (G) holds iff $\ell_i(g,\cdot)$ is strictly increasing in world-$i$ rank on $A(g)$ for a.e.\ $g$ (\src{R1} Lemma 1.1); such penalties are \emph{strictly guiding}. A model is $p$-Bayes if it minimises $p\cdot L$ on $\M$; equivalently its advice minimises $p\cdot\ell(g,\cdot)$ over the advice range $A(g)$ for a.e.\ $g$ (\src{R1} Lemma 1.2).

For a pair with $A(g)\subseteq\{+,0,-\}$ put $c\coloneqq\ell(g,0)-\ell(g,+)$, $c'\coloneqq\ell(g,0)-\ell(g,-)$ and $d\coloneqq c-c'$ (vectors over worlds). By (G), $c$ and $d$ have the strict sign pattern of $g$ and $c'$ that of $-g$. The \emph{chart} $R(h)$ of an orientation $h$ is the set of priors at which accepting $h$ is the unique minimiser. For advice range $\{+,0\}$, $R(g)=\{p:c\cdot p>0\}$. For $\{+,0,-\}$, $R(g)=\{c\cdot p>0,\,d\cdot p>0\}$, $R(-g)=\{c'\cdot p>0,\,d\cdot p<0\}$, and no advice is the minimiser on $\{c\cdot p\le0,\,c'\cdot p\le0\}$ (\src{R1} Prop.~2.1). Such a pair is \emph{adjacent} if its d-face $F(g)\coloneqq\{p\in\Di:c\cdot p=c'\cdot p>0\}$ is nonempty; otherwise it is \emph{split}. In words: a pair is adjacent when, at some interior weighting, accepting $g$ and accepting $-g$ are tied as the best advice and each wins on one side, so that the two acceptances compete directly across a boundary between interior weightings. The advice range of a split pair may still contain both acceptances, provided the region where no advice is best separates their charts. Adjacency therefore depends on the penalties, not only on the advice range. Mixed pairs with advice range $\{+,-\}$ (on which every evaluable model accepts one orientation and not the other) are always adjacent (\src{R4} Thm.~A).

\subsection{The \texorpdfstring{$U$}{U}-domain}\label{sec:aos}
Fix a credal set $K'$ in the relative interior of the target $T$ (for $T_1$, $K'\coloneqq u+s(T_1-u)$ with $0<s<1$) and put $U\coloneqq D_{K'}$, an open convex cone containing $\cl D_T\setminus\{0\}$. The \emph{$U$-domain} $\M_U$ is the class of measurable subsets $D$ of $U$ (modulo null sets); it is closed under splicing and under convergence in measure. This is a restriction on the class of evaluable models, and only that; its consequences, for every evaluated model $D\in\M_U$, are the following.
\begin{itemize}[leftmargin=*]
\item $U$ never contains both $f$ and $-f$, because $P_{K'}(f)>0$ and $P_{K'}(-f)>0$ cannot both hold (and $-f\notin U$ for $f\ge0$, $f\ne0$, since $K'\subseteq\Di$). So no evaluated model gives contradictory advice.
\item For $f\in U$ and any $g,h$ with $f=g-h$, $D$ advises choosing $g$ in $\{g,h\}$ iff $f\in D$, and gives no advice iff $f\notin D$.
\item For $f$ with $-f\in U$ the same holds with the roles of $g$ and $h$ swapped: $D$ advises choosing $h$ iff $-f\in D$, and otherwise gives no advice.
\item For $f$ with $f\notin U$ and $-f\notin U$, every evaluated $D$ has $f\notin D$ and $-f\notin D$, so $D$ gives no advice in $\{g,h\}$.
\end{itemize}
Nothing more follows from the restriction. The problems of the last kind (pairs $\{f,-f\}$ with $f\notin U$ and $-f\notin U$) remain decision problems, and the world orders still concern them; every evaluable model gives no advice on them, so they do not distinguish evaluable models. It is not enough that $f\notin U$: if $-f\in U$, evaluable models differ on $f$. (LF) holds at $D_T$ with any open $N$ such that $\partial D_T\setminus0\subseteq N\subseteq U$. On $\M_U$ we index each pair that meets $U$ by its orientation $g\in U$. Only the vector $c(g)=\ell(g,0)-\ell(g,+)$, $g\in U$, enters comparisons between evaluable models (on every pair $\{f,-f\}$ with $f\notin U$ and $-f\notin U$ all evaluable models give no advice, so its penalty drops out), and the $q$-Bayes model is
\[
B_q=\{g\in U:\ q\cdot c(g)>0\}.
\]
An additive ranking on a general decomposable domain is \emph{non-adjacent near $T$} if, for some credal set $K'\subseteq\relint T$, a.e.\ pair with an orientation in $D_{K'}$ is non-adjacent. The advice range of such a pair may contain both acceptances; what is excluded is that they compete directly as best advice across a boundary between interior weightings. This holds on $\M_U$, whenever the advice ranges of these pairs are contained in $\{+,0\}$ (a sufficient special case that concerns the domain alone), and whenever every interior Bayes credal set is two-dimensional (\src{R4} Thm.~B).

\subsection{Coherent almost-preferences}\label{sec:CPa-def}
Let $L^\infty(H,\mu)$ be the bounded measurable functions on $H$ modulo $\mu$-null sets. A \emph{state} is a positive normalised linear functional on $L^\infty(H,\mu)$, that is, a finitely additive probability on the measurable sets that vanishes on $\mu$-null sets.

\begin{definition}\label{def:CPa}
World $i$ of a ranking on a decomposable domain $\M$ is a \emph{coherent almost-preference over losses}, condition (CP-a), if there are
\begin{enumerate}[leftmargin=*,label=(\roman*)]
\item a bounded measurable penalty $\lambda_i(g,a)$ with $\lambda_i(g,a)<\lambda_i(g,b)$ whenever $a$ is strictly better than $b$ in world-$i$ rank at $g$ (for a.e.\ $g$), and
\item a coherent lower prevision $\lp_i$ on $L^\infty(H,\mu)$ in the sense of Walley~\cite{W1991}, with ``$f\ge0$'' read as $f\ge0$ a.e.,
\end{enumerate}
such that, writing $\Lambda_iX\coloneqq\lambda_i(\cdot,a_X(\cdot))$,
\[
X\preceq_iY\iff\lp_i(\Lambda_iY-\Lambda_iX)\ge0 .
\]
\end{definition}
Equivalently (\cite[Thm.~3.8.1]{W1991}, together with the lower envelope theorem \cite[Thm.~3.3.3]{W1991} and weak-* compactness \cite[\S3.6]{W1991}): $\{f:\lp_i(f)\ge0\}$ is a coherent class of almost-desirable gambles (axioms D0--D4); or $\preceq_i$ is the pull-back along $X\mapsto-\Lambda_iX$ of a coherent almost-preference (axioms R0--R5); or there is a nonempty weak-* closed convex set $\M_i$ of states with
\[
X\preceq_iY\iff P(\Lambda_iY-\Lambda_iX)\ge0\ \text{for every }P\in\M_i .
\]
The null-set reading is a restriction in Walley's own sense: (CP-a) holds iff there is a Walley-coherent lower prevision on bounded measurable gambles that is \emph{null-invariant} ($\lp_i(f)=\lp_i(f')$ whenever $f=f'$ a.e.), equivalently every member of its credal set vanishes on null sets (\src{N24} \S5). Not assumed: countable additivity, separability or non-atomicity of $\M_i$; completeness; conglomerability. World $i$ is \emph{complete} if $\M_i$ is a singleton. Instances include additive rankings with integral penalties ($\M_i=\{\rho\mu\}$), additive rankings with bounded charge penalties ($\M_i=\{P\}$), Pareto unanimity of additive rankings that share the encoding, and credal-set rankings. On $\M_U$ the canonical encoding is the error indicator $\Lambda_iX=1_{E_i(X)}$ with $E_i(X)\coloneqq X\sd(U\cap\{h_i>0\})$ (a function on $U$, carried to the pairs by the measure-preserving bijection $g\mapsto\operatorname{sign}(g_1)\,g$ from $U$ onto the representatives of the pairs meeting $U$; it vanishes on the pairs with $f\notin U$ and $-f\notin U$).

\subsection{Adding up and supporting additive assessments}\label{sec:addup-def}
This subsection fixes the terms used for rankings that are not assumed to be additive (Glossary of \src{START HERE}; \src{N17}, \src{N19}). Work on $\M_U$ with a ranking $L$ satisfying strict (G), (IB), (LA), (S), (C). The \emph{error set} $E_i(X)$ of a model at world $i$ is the set of $h\in U$ with $h_i>0$ that $X$ does not accept, together with those with $h_i<0$ that it accepts, as above. A model is \emph{regular} if it gives the correct advice on unanimous problems, i.e.\ accepts a.e.\ nonnegative gamble of $U$.

\begin{definition}\label{def:addup}
\begin{enumerate}[leftmargin=*,label=(\alph*)]
\item The \emph{relevant comparisons} at world $i$ are the pairs $B\preceq_iX$ with $B\in\C_L$ (coherent and undominated) and $X$ regular.
\item World $i$ \emph{adds up} if some finitely additive weight $\nu_i$ on $U$, zero on null sets and positive on non-null sets, endorses every relevant comparison: $\nu_i(E_i(B))\le\nu_i(E_i(X))$.
\item If $\nu_i$ is such a weight, let $\rho_i$ be the density of its countably additive part; then $\rho_i>0$ a.e.\ on $U$ (\src{N19} Lemma R1). The integral loss $\psi_i(X)\coloneqq\int_{E_i(X)}\rho_i\,d\mu$ is a \emph{supporting additive assessment}. It is an additive ranking on the same domain. It endorses the relevant comparisons when $\nu_i$ is countably additive, but it is not claimed to be numerically or topologically close to $\preceq_i$, and it can lose distinctions that $\preceq_i$ makes.
\end{enumerate}
\end{definition}

Restricting to regular $X$ is essential. Ranking first by the mass of unanimous problems of $U$ on which the model gives no advice, and only then by $\preceq_i$, preserves every axiom, the selection, $\C_L$ and the covered target, yet it destroys adding up for unrestricted comparisons (\src{N19} \S3.3, Prop.~R3, proved in that referee report). Such a priority manufactures apparent non-additivity without helping coverage, so only the relevant comparisons are tested.

Exactly one of the following holds.
\begin{itemize}[leftmargin=*]
\item[\textbf{(M0)}] Every world adds up.
\item[\textbf{(M1)}] At some world a finite family of relevant comparisons is an \emph{obstruction} to strictly positive error weights: some positive combination of their error-difference vectors $1_{E_i(X)}-1_{E_i(B)}$ is $\le0$ a.e.\ and not a.e.\ zero. The comparisons need not form an executable sequence of trades.
\item[\textbf{(M2)}] At some world there is no finite obstruction, but still no strictly positive weight endorses all relevant comparisons.
\end{itemize}
The division is logical bookkeeping. The substantive result is the reduction behind (M0) (Proposition~\ref{prop:M0}): supporting weights, if they exist, yield an additive system that is coherent everywhere except for ties at exceptional weightings. The restrictions known for (M1) (at least four comparisons, overlapping gains, crossing error sets; \src{N17} N17.3, refereed in \src{N19} \S4) concern (M1) certificates in this formulation only; they do not constrain every failure of adding up.

\section{Additive coverage as a realisation problem}\label{sec:reduction}

From now on $n=3$. This section reduces additive coverage of an interior target to a statement about a family of credal sets.

\begin{proposition}\label{prop:necessary}
Let $L$ be a strictly guiding additive ranking on a decomposable domain $\M$ satisfying (LA), (S), (C). Then:
\begin{enumerate}[leftmargin=*,label=(\alph*)]
\item for every $p\in\Di$ the $p$-Bayes model is a.e.\ unique and equals $D_{K_p}$ a.e.\ for a credal set $K_p$;
\item at each vertex $e_i$ the Bayes model is the world-$i$-best evaluable model, and it is coherent;
\item at every $p\in\partial\Delta$ some $p$-Bayes model is coherent;
\item if $L$ covers $T=D_K$ with $K\subseteq\Di$, then $T$ is the $p^*$-Bayes model for some $p^*\in\Di$.
\end{enumerate}
\end{proposition}
\status{proved; refereed once (R6). Sources: \src{R1} Thm.~2.2(N), Lemmas 1.4--1.5, Prop.~2.6; \src{R2} Lemma U. Uniqueness here is uniqueness of the Bayes minimiser at each interior weight, not of the selection, which may be reindexed (Consolidated Statements \S0.3).}

\begin{definition}\label{def:realisation}
Let $T\subseteq\Di$ be a credal set and $U=D_{K'}$ as above. A \emph{realisation} of $T$ on $\M_U$ at $p^*\in\Di$ is a measurable $c:U\to\R^3$ such that $c_i(g)$ has the strict sign of $g_i$ for a.e.\ $g$ and:
\begin{enumerate}[leftmargin=*]
\item[\textbf{(BU)}] for \emph{every} $q\in\Di$ the tie set $\{g\in U:q\cdot c(g)=0\}$ is null and $B_q=D_{K_q}$ a.e.\ for a credal set $K_q$;
\item[\textbf{(TP)}] $K_{p^*}=T$.
\end{enumerate}
\end{definition}
(These are conditions (R1), (R2) of \src{N12}.) By Proposition~\ref{prop:necessary}, a strictly guiding additive ranking on $\M_U$ that covers $T$ yields a realisation; Theorem~\ref{thm:ccibu} gives the converse. Write $P(q,g)\coloneqq P_{K_q}(g)$ and, for $g\in U$, $H(g)\coloneqq\{q\in\aff\Delta:q\cdot c(g)>0\}$, an open half-plane.

\begin{proposition}\label{prop:reduction}
Let $c$ be a realisation of $T$ at $p^*$. Then:
\begin{enumerate}[leftmargin=*,label=(\alph*)]
\item $K_q\supseteq K'$ for every $q\in\Di$, and $K_q$ is unique;
\item $q\mapsto K_q$ is Hausdorff-continuous on $\Di$, and $P$ is jointly continuous with $|P(q,g)-P(q',g')|\le\|g-g'\|_\infty+\|g\|\,\dH(K_q,K_{q'})$;
\item the vertex models $U\cap\{g_i>0\}=D_{\co(K'\cup\{e_i\})}$ are coherent;
\item $\tilde c\coloneqq p^*\odot c$ is a realisation of $T$ at $u$;
\item \emph{(half-plane charts)} for a.e.\ $g\in U$, $\{q\in\Di:P(q,g)>0\}$ equals $H(g)\cap\Di$ up to a null set, and $e_i\in H(g)$ if $g_i>0$, $e_i\notin\cl H(g)$ if $g_i<0$.
\end{enumerate}
Conversely, a Hausdorff-continuous family $q\mapsto K_q\supseteq K'$ with $K_{p^*}=T$ and vertex-respecting half-plane charts as in (e) comes from a realisation, and its tie sets are automatically null.
\end{proposition}
\begin{proof}[Proof sketch]
(a) $D_{K_q}$ is open, convex and a.e.\ inside the open convex $U$, so $D_{K_q}\subseteq U$, i.e.\ $K_q\supseteq K'$; injectivity of $K\mapsto D_K$ gives uniqueness. (b) If $q_m\to q$ then $1\{q_m\cdot c>0\}\to1\{q\cdot c>0\}$ off the null tie set of $q$, so $B_{q_m}\to B_q$ in measure; for cones with interior this is Hausdorff convergence of credal sets (\src{R1} Prop.~3.1). (c) $D_{K_1}\cap D_{K_2}=D_{\co(K_1\cup K_2)}$. (d) $\tilde q\cdot\tilde c(g)=(\tilde q\odot p^*)\cdot c(g)$, and normalising $\tilde q\odot p^*$ is a projective bijection of $\Delta$ fixing the vertices. (e) Fubini on $\{(q,g):1[q\cdot c(g)>0]\ne1[P(q,g)>0]\}$. Converse: \src{N12} Prop.~1.4 with the simplification of \src{N14} \S2.1.
\end{proof}
\status{proved; refereed once (N14: CORRECT, minor repair incorporated). Source: \src{N12} \S1 (Lemma 1.1, Props 1.2--1.4). The irregularity of $c$ has no effect beyond the family $q\mapsto K_q$.}

\begin{remark}\label{rem:N20}
Condition (BU) must hold at \emph{every} interior $q$, or Hausdorff continuity must be assumed. If (BU) held only for a.e.\ $q$, then $K_q$ could be altered on a null set of priors that contains $p^*$, and (TP) would carry no information (\src{N20} \S8). For a covering ranking this is automatic, by Proposition~\ref{prop:necessary}(a).
\end{remark}

\begin{theorem}\label{thm:ccibu}
Let $n=3$ and let $L$ be an additive ranking on a decomposable domain with integrable envelopes and penalties strictly increasing in world rank a.e. Suppose that for every $w\in\Di$ the minimiser of $w\cdot L$ on $\M$ is unique modulo null sets and coherent. Then $L$ satisfies (LA), (S), (C), and every undominated model is coherent and the only model with its loss profile.
\end{theorem}
\begin{proof}[Proof sketch]
The attainable set $\mathcal R=\{L(X):X\in\M\}$ is compact and convex (Lyapunov). Every Pareto-undominated profile has a supporting weight $w$. If $w$ is interior, use the hypothesis. If $w=e_i$, strict guidance gives a unique minimiser, and it is the limit of interior Bayes models. If $w=(w_1,w_2,0)$, a finite subgradient of $t\mapsto\min\{y_3:y\in\operatorname{argmin}w\cdot y,\ y_1=t\}$ turns the profile into a directional minimiser on $\operatorname{argmin}_aw\cdot\ell(g,a)$; directional minimisers are unique (a tie would persist at nearby interior weights) and are limits of interior Bayes models, hence coherent. Then select $x(p)$ by minimising $Q_p(x)=\sum_ip_i(x_i-a_i)^2$ over $\mathcal R$ (with $a_i<\min_{\mathcal R}x_i$), followed at edge indices by a secondary minimisation of the omitted coordinate. Each $x(p)$ is undominated. Coordinates of $x(p)$ are lower semicontinuous in $p$, which makes the sections closed, and strict monotonicity of $Q_p$ gives competitiveness on $\supp p$.
\end{proof}
\status{author-checked (\src{CCIBU.tex}; also checked by the coordinator, \src{BRIEF2}); Lemmas 2--4 re-derived without error in \src{N13} \S2; no formal referee report. The argument is specific to three worlds.}

\begin{corollary}\label{cor:realisation}
Let $T=D_K$ with $K\subseteq\Di$, and let $\M$ be a decomposable domain with (LF) at $T$. Some strictly guiding additive ranking with integrable envelopes on $\M$ that satisfies (LA), (S), (C) covers $T$ iff some strictly guiding additive ranking with integrable envelopes on $\M$ has all interior Bayes models a.e.\ unique and coherent, with $T$ among them. On $\M_U$ this holds iff $T$ has a realisation (Definition~\ref{def:realisation}).
\end{corollary}
\status{proved, from Proposition~\ref{prop:necessary} and Theorem~\ref{thm:ccibu} (so its status is that of Theorem~\ref{thm:ccibu} in the ``if'' direction). For the realisation form, penalties are recovered from $c$ by \src{R1} Prop.~2.1 (converse) and a positive rescaling for integrability, which does not change charts.}

\section{No additive cover of the contamination triangle}\label{sec:main}

\begin{theorem}\label{thm:main}
Let $T\subseteq\Di$ be an all-facing triangle in general position, and $p^*\in\Di$. Then $T$ has no realisation at $p^*$ on any $U$-domain $\M_U$. In particular, for every $\varepsilon\in(0,1)$, $T_1$ has no realisation at any interior target weighting.
\end{theorem}
\status{proved; refereed independently twice (\src{N18}, \src{N20}: both CORRECT, with the local repairs incorporated below). Source: \src{N15} (main chain), assembled with \src{N12} Lemma 2.1, Lemma 4.3, Thm.~2.4, Thm.~3.1 (refereed once, \src{N14}), \src{R10b} Lemma UVN and Cor.~PL2 (refereed once, \src{R12}), and \src{R15} Thm.~6.2 algebra (refereed once, \src{R18}). No hypothesis of differentiability, nondegeneracy or regularity of $q\mapsto K_q$ is used beyond the Hausdorff continuity that (BU) provides (Remark~\ref{rem:N20}).}

\begin{corollary}\label{cor:coverage}
Let $L$ be a strictly guiding additive ranking with integrable envelopes on a decomposable domain $\M$, satisfying (G), (IB), (LA), (S), (C). If $L$ is non-adjacent near $T$ (for $T$ as in Theorem~\ref{thm:main}), then $L$ does not cover $T$. In particular $T_1$ is not covered by any additive ranking on the $U$-domain, on any decomposable domain on which the pairs with an orientation in some $D_{K'}$, $K'\subseteq\relint T_1$, have advice ranges contained in $\{+,0\}$, or whenever all interior Bayes credal sets are two-dimensional.
\end{corollary}
\status{proved; refereed twice (\src{N18} \S1.10, \src{N20} \S7: Cor.~7.2 of \src{N15} CORRECT). It rests on the reduction of \src{N14} \S3(iv), which was proved in that referee report: under non-adjacency, $\hat K_q\coloneqq\co(K_q\cup K')$ satisfies Proposition~\ref{prop:reduction} with $U=D_{K'}$, so by its converse it is a realisation on $\M_{D_{K'}}$. The sources also list domains whose advice ranges near $\partial D_T$ have two elements (``two-action'' domains there) as an instance; the hypothesis actually used, through \src{N14} \S3(iv), is non-adjacency of all pairs meeting $D_{K'}$, which is how it is stated here.}

\begin{corollary}\label{cor:charges}
The conclusion of Corollary~\ref{cor:coverage} holds when the penalties are bounded charges vanishing on $\mu$-null sets, provided the interior Bayes minimisers exist.
\end{corollary}
\status{proved; \src{N14} Thm.~C2 (proved in that referee report; its key splice step is re-proved and refereed as \src{N22} Lemma 3.0 / \src{N24}), used in \src{N15} Cor.~7.3, refereed CORRECT conditional on existence in \src{N18} and \src{N20}. The existence proviso is discharged for complete rankings in Theorem~\ref{thm:complete}.}

The rest of this section proves Theorem~\ref{thm:main}. Fix a realisation $c$ of an all-facing triangle $T$ in general position at $p^*$. By Proposition~\ref{prop:reduction}(d) we may and do assume $p^*=u$; nothing below uses this before \S\ref{sec:contradiction}, and the earlier steps hold verbatim at any interior $p^*$. Put $r_0\coloneqq\operatorname{dist}(p^*,\partial\Delta)$ in $\aff\Delta$, so $B(p^*,r_0)\subseteq\Di$.

\subsection{Chords and straight levels}

For each facet $F_k$ with corner rays $r,r'$ put $h(\tau)\coloneqq(1-\tau)r+\tau r'$ and $\Psi(q,\tau)\coloneqq P(q,h(\tau))$. For each $q$, $\tau\mapsto\Psi(q,\tau)$ is concave on all of $\R$, because $P(q,\cdot)$ is concave and $h$ is affine. Since $K_{p^*}=T$, $\Psi(p^*,\tau)=0$ for $\tau\in[0,1]$. A gamble $g\in U$ is \emph{good} if $\{q:P(q,g)=0\}$ is null and the half-plane chart of Proposition~\ref{prop:reduction}(e) holds for it; a.e.\ $g$ is good. A \emph{cluster line} of a mixed $g\in U$ is a limit, with orientation, of the lines $\partial H(g_m)$ for good $g_m\to g$ (normalise $c(g_m)$ to unit length; a limit unit vector defines an affine function on $\aff\Delta$ and its zero line). Cluster lines exist by compactness.

\begin{lemma}\label{lem:signs}
On each facet $F_k\setminus\{0\}$, $h_k<0$. After relabelling $i,j$ so that $r=r_{ki}$, the coordinates $h_i,h_j$ each cross zero exactly once on $(0,1)$, at $\tau_i<\tau_j$, with sign patterns $(-,+)$ on $[0,\tau_i)$, $(+,+)$ on $(\tau_i,\tau_j)$ and $(+,-)$ on $(\tau_j,1]$.
\end{lemma}
\begin{proof}
Write $v_k=ae_k+bv_i+cv_j$ with $a,b,c>0$ (all-facing). For $0\ne h\in F_k$, $0=h\cdot v_k=ah_k+b\,h\cdot v_i+c\,h\cdot v_j$, and $h\cdot v_i,h\cdot v_j\ge0$ are not both zero (the $v$'s span $\R^3$), so $h_k<0$. At $r_{ki}$, $r_k<0$ and $r_i<0$ (the same fact on $F_i$), and $r_j>0$, since $r\le0$ would give $r\cdot v_j<0$. Likewise at $r_{kj}$: $r_k,r_j<0<r_i$. Along the facet $h_i,h_j$ are affine, so each crosses once; they cannot both be $\le0$ at one point, since then $h\le0$, $h\ne0$ and $h\cdot v_k<0$. Hence at the crossing of $i$ we have $h_j>0$, and $\tau_i<\tau_j$.
\end{proof}
\status{proved; \src{N15} \S1 for $T_1$, general proof supplied independently by both referees of \src{N15} (\src{N18} (R-c), \src{N20} \S4.0).}

\begin{lemma}\label{lem:chords}
\begin{enumerate}[leftmargin=*,label=(\alph*)]
\item For good $g$: $P(\cdot,g)\ge0$ on $\cl H(g)\cap\Di$, $P(\cdot,g)\le0$ on $\Di\setminus H(g)$, and $P(\cdot,g)=0$ on $\partial H(g)\cap\Di$.
\item For every mixed $g\in U$ and every cluster line $\ell$ with closed sides $\Pi^\pm$: $P(\cdot,g)\ge0$ on $\Pi^+\cap\Di$, $P(\cdot,g)\le0$ on $\Pi^-\cap\Di$, $P(\cdot,g)=0$ on $\ell\cap\Di$; and $e_i\in\Pi^+$ if $g_i>0$, $e_i\in\Pi^-$ if $g_i<0$.
\end{enumerate}
\end{lemma}
\begin{proof}
(a) The open set $\{P(\cdot,g)>0\}$ is a.e.\ inside the open half-plane $H(g)$, hence inside it; similarly $\{P<0\}\subseteq\Di\setminus\cl H(g)$; continuity in $q$ gives the closed statements. (b) Pass to the limit using joint continuity (Proposition~\ref{prop:reduction}(b)); the vertex statements are the limits of the strict signs of $c$.
\end{proof}
\status{proved; refereed once (\src{N14}: CORRECT). Source: \src{N12} Lemma 2.1.}

The level shift $P(q,g-t\one)=P(q,g)-t$ (since $m\cdot\one=1$) turns the chord statement at level $0$ into one at every level. For $g\in\R^3$ let $I_g\coloneqq\{t:\ g-t\one\in U\text{ and is mixed}\}=(\min_ig_i,\,P_{K'}(g))$.

\begin{lemma}\label{lem:levels}
Let $f\coloneqq P(\cdot,g)$ on $\Di$ and $t\in I_g$. Then $\{f>t\}$ and $\{f<t\}$ are each the intersection of $\Di$ with an open half-plane, or $\varnothing$, or $\Di$. Consequently, when both half-planes are proper, $\{f=t\}$ is the closed region between two lines that do not meet in $\Di$.
\end{lemma}
\begin{proof}
For $t'\in I_g$, Lemma~\ref{lem:chords}(b) applied to $g-t'\one$ gives a line with closed sides $\Pi^\pm_{t'}$ such that $f\ge t'$ on $\Pi^+_{t'}\cap\Di$ and $f\le t'$ on $\Pi^-_{t'}\cap\Di$. \emph{$\{f>t\}$ is convex:} if $f(q_1),f(q_2)>t$, choose $t'\in(t,\min f(q_i))$, which lies in $I_g$ because $f\le P_{K'}(g)$; then $q_1,q_2\in\operatorname{int}\Pi^+_{t'}$, and the segment lies in $\operatorname{int}\Pi^+_{t'}\cap\Di\subseteq\{f\ge t'\}$. \emph{$\{f\le t\}$ is convex:} it equals $\bigcap_{t'\in(t,t+\eta)}\Pi^-_{t'}\cap\Di$ for small $\eta$. Two complementary convex subsets of the convex set $\Di$, one of them open, are separated by a line, which gives the half-plane form. The case $\{f<t\}$ is symmetric. If the two boundary lines crossed at $x\in\Di$, the two disjoint open half-planes would share a wedge at $x$.
\end{proof}
\status{proved; refereed twice (\src{N18} \S1.1, \src{N20} \S1, wording repair: ``region between two lines'' rather than ``strip''). Source: \src{N15} Lemma 2.1.}

\begin{lemma}\label{lem:sandwich}
Let $f=P(\cdot,g)$, $f(p^*)=0$, and suppose $\{f>0\}=\Di\cap\{n\cdot(q-p^*)>0\}$ for a unit $n\in\one^\perp$ and $f\le0$ on $\lambda\cap B(p^*,r_0)$, where $\lambda\coloneqq p^*+\R u_\perp$ with $u_\perp\perp n$ a unit vector. Put $D^+(s)\coloneqq f(p^*+sn)$ and $D^-(s)\coloneqq-f(p^*-sn)$. Then $D^\pm$ are nondecreasing near $0$, and for $q=p^*+yu_\perp+xn$ with $|y|\le r_0/2$ and $x>0$ small,
\[
D^+\Bigl(\frac{x\,r_0}{r_0+|y|}\Bigr)\ \le\ f(q)\ \le\ D^+\Bigl(\frac{x\,r_0}{r_0-|y|}\Bigr),
\]
with the mirror statement for $x<0$.
\end{lemma}
\begin{proof}
Monotonicity: for $t>0$ in $I_g$, $\{f>t\}$ is a half-plane inside $\{f>0\}$ missing $p^*$, so it meets the ray $p^*+\R_+n$ in a terminal sub-ray. Upper bound: if $D^+(s_1)<f(q)$ with $s_1>xr_0/(r_0-|y|)$, choose a level $t$ strictly between; the convex set $\{f<t\}$ contains the open chord $\{p^*+y'u_\perp:|y'|<r_0\}$ and $p^*+s_1n$, hence the open triangle they span, whose height over $y$ is $s_1(1-|y|/r_0)>x$; so $f(q)<t$, a contradiction. The lower bound is the same argument with the convex set $\{f\le t\}$ and apex $q$.
\end{proof}
\status{proved; refereed twice (\src{N18} \S1.3, \src{N20} \S1, with the convex-hull proof of the upper bound and the weakened hypothesis $f\le0$ on $\lambda\cap B(p^*,r_0)$). The constants $r_0/(r_0\pm|y|)$ are sharp. Source: \src{N15} Lemmas 2.2--2.3.}

\begin{lemma}\label{lem:caps}
Let $\gamma:\R\to\R$ be concave.
\begin{enumerate}[leftmargin=*,label=(\alph*)]
\item If $\{\gamma>0\}\ne\varnothing$, then $\{\gamma=0\}=\partial\{\gamma>0\}$ (at most two points) and $\gamma<0$ off $\cl\{\gamma>0\}$; if $\gamma$ vanishes on a nondegenerate interval, $\gamma\le0$ everywhere.
\item If $\gamma(\tau_0)=0$ and $\gamma>0$ on $(\tau_0,\tau_0+\eta)$, then $\gamma(\tau)/(\tau-\tau_0)$ is nonincreasing on $\R\setminus\{\tau_0\}$, and $0<\gamma'_+(\tau_0)\le\gamma'_-(\tau_0)<\infty$. The mirror statement holds on the left.
\item For a jointly continuous $\Gamma(q,\tau)$ concave in $\tau$, $\partial^-_\tau\Gamma=\inf_{s>0}(\Gamma(q,\tau)-\Gamma(q,\tau-s))/s$ is upper semicontinuous and $\partial^+_\tau\Gamma=\sup_{s>0}(\Gamma(q,\tau+s)-\Gamma(q,\tau))/s$ is lower semicontinuous.
\end{enumerate}
\end{lemma}
\status{proved (standard convex analysis); refereed twice (\src{N18} \S1.4, \src{N20} \S2, endpoint slip repaired). Source: \src{N15} Lemma 3.1.}

\subsection{Vertex neighbourhoods and single-vertex pinning}

\begin{lemma}\label{lem:UVN}
Let $x\in U$ be mixed and $l$ a world. If $x_l>0$ there are a ball $B_1\ni x$ and a relative neighbourhood $V_l$ of $e_l$ in $\Delta$ with $P(q,h)>0$ for all $q\in V_l\cap\Di$, $h\in B_1$, and $H(h)\supseteq V_l\cap\Di$ for a.e.\ $h\in B_1$. If $x_l<0$ the same holds with $P(q,h)<0$ and $H(h)\cap V_l=\varnothing$.
\end{lemma}
\status{proved; refereed once (\src{R12}: CORRECT). Source: \src{R10b} Lemma UVN. It uses the strict minimiser of the vertex objective and dominated convergence as $q\to e_l$, together with upward (downward) closure of coherent cones.}

Fix a vertex $v=v_k$ of $T$ with a crossing $w\in\relint F_k$, $w_c=0$, and a small ball $\mathcal U$ around $w$ inside $U$, meeting no other facet, and inside the balls of Lemma~\ref{lem:UVN} at $w$ and at points $w_\pm\in F_k$ near $w$ with $(w_\pm)_c\gtrless0$. Put $\mathcal U^+=\mathcal U\cap D_T$, $\mathcal U^-=\mathcal U\setminus\cl D_T$ and
\begin{gather*}
\Gamma_v\coloneqq\{q\in\Di:D_{K_q}\supseteq\mathcal U^+\},\qquad
\beta_v\coloneqq\Gamma_v\setminus\alpha_v,\\
\alpha_v\coloneqq\{q\in\Gamma_v:v\in K_q\}=\{q\in\Gamma_v:D_{K_q}\cap\mathcal U^-=\varnothing\}.
\end{gather*}
$\beta_v$ is relatively open in $\Gamma_v$ (\src{R10b} Lemma D, as repaired in \src{R12}). For $q\in\beta_v$, $D_{K_q}$ contains an exterior collar of $F_k\cap\mathcal U$; for $q\in\alpha_v$, $P(q,h)<0$ for all $h\in\mathcal U^-$.

\begin{lemma}\label{lem:PL2}
Let $G\subseteq\Gamma_v$ be open and connected, and suppose a.e.\ $h\in\mathcal U^-$ has no d-face meeting $G$ (automatic on $\M_U$, where no advice range contains both acceptances). Then $G\subseteq\alpha_v$ or $G\subseteq\beta_v$.
\end{lemma}
\begin{proof}[Proof sketch]
It suffices to show that no $Q\in G$ lies in $\cl\beta_v\setminus\beta_v$. Suppose $Q$ does; then $Q\in\alpha_v$, so $P(Q,h)<0$ on $\mathcal U^-$, and $Q\notin\cl H(h)$ for good $h\in\mathcal U^-$. Choose $s_m\in\beta_v$ with $s_m\to Q$, and good $h^\pm_m\in\mathcal U^-$ converging to $w_\pm$, so close to the facet that $P(\cdot,h^\pm_m)>0$ near $s_m$ (collar). Then the chart lines of $h^\pm_m$ pass within $2|s_m-Q|$ of $Q$; normalised, $c(h^\pm_m)\to a_\pm$ with $a_\pm\cdot Q=0$. Every point of $\{a_+\cdot q>0\}$ near $Q$ has a neighbourhood a.e.\ accepting some $h^+_m\in\mathcal U^-$, so it lies in $\cl\beta_v$; every point of $\beta_v$ near $Q$ has $P(\cdot,h^-_m)>0$ for large $m$, so $a_-\cdot q\ge0$ there. Hence $\{a_+>0\}\subseteq\cl\{a_->0\}$ near $Q$, and since both lines pass through $Q$, $a_+=\lambda a_-$ with $\lambda>0$. But $(a_-)_c\le0$, because $c_c(h^-_m)<0$, while Lemma~\ref{lem:UVN} at $w_+$ for world $c$ gives $\{a_+\cdot q>0\}\supseteq V_c\cap\Di$, hence $(a_+)_c>0$ (a line through $e_c$ and the interior point $Q$ would cut $V_c$). This is a contradiction. So $\beta_v$ is relatively clopen in the connected set $G$.
\end{proof}
\status{proved; refereed once (\src{R12}: CORRECT). Source: \src{R10b} Lemma M (Steps 2--4) and Cor.~PL2. Its validity for sets $G$ that do not contain $p^*$, which is how it is used below, was re-checked by both referees of \src{N15} (\src{N18} \S1.5, \src{N20} \S3).}

\begin{lemma}\label{lem:nopin}
Let $\tau_c\in(\sigma_1,\sigma_2)\subseteq[0,1]$ be a crossing parameter of $F_k$, and let $G\subseteq\Di$ be nonempty, open and connected, with $\Psi(\cdot,\sigma_1)\ge0$ and $\Psi(\cdot,\sigma_2)\ge0$ on $G$. Then $\Psi>0$ on $G\times(\sigma_1,\sigma_2)$.
\end{lemma}
\begin{proof}
By concavity in $\tau$, $\Psi\ge0$ on $G\times[\sigma_1,\sigma_2]$. Choose the ball $\mathcal U$ around $w=h(\tau_c)$ so small that $F_k\cap\cl\tilde{\mathcal U}\subseteq\R_+h([\sigma_1,\sigma_2])$ for a slightly larger concentric ball $\tilde{\mathcal U}$. By homogeneity $K_q$ lies in the wedge $W\coloneqq\{m\in\aff\Delta:m\cdot h(\sigma_1)\ge0,\ m\cdot h(\sigma_2)\ge0\}$, with apex $v_k$, and then $q\in\Gamma_{v_k}$ (\src{R10b} \S1: write $h\in\mathcal U^+$ as $w''+(v_k\cdot h)\one$ with $w''\in F_k\cap\tilde{\mathcal U}$). Pick $j$ with $w_j>0$; by Lemma~\ref{lem:UVN}, $V_j\cap\Di\subseteq\Gamma\cap\beta$. Every good $h\in\mathcal U^+$ has $H(h)\supseteq G\cup(V_j\cap\Di)$, so $\operatorname{int}\co(G\cup V_j)\subseteq\Gamma$ (a Fubini step, checked in \src{N20} \S3). By Lemma~\ref{lem:PL2} this connected open set lies in $\alpha$ or in $\beta$; it meets $\beta$ at $V_j$, so it lies in $\beta$, and $G\subseteq\beta$, i.e.\ $v_k\notin K_q$ for $q\in G$. Since $h(\tau)$ is a strictly positive combination of $h(\sigma_1),h(\sigma_2)$ for $\tau\in(\sigma_1,\sigma_2)$, and these are independent modulo $\one$, $m\cdot h(\tau)>0$ on $W\setminus\{v_k\}$; compactness of $K_q$ gives $\Psi(q,\tau)>0$.
\end{proof}
\status{proved; refereed twice (\src{N18} \S1.5, \src{N20} \S3). Source: \src{N15} Lemma 4.1. This is the step where non-adjacency enters (through Lemma~\ref{lem:PL2} on a set that leaves every ball around $p^*$).}

\subsection{Pivot chords and corners}

Write $A_\tau\coloneqq\{\Psi(\cdot,\tau)>0\}$, $B_\tau\coloneqq\{\Psi(\cdot,\tau)<0\}$, $Z_\tau\coloneqq\{\Psi(\cdot,\tau)=0\}$, $P_\tau\coloneqq\Di\setminus B_\tau$, all subsets of $\Di$.

\begin{theorem}\label{thm:pivot}
For every facet $F_k$ and $\tau\in(0,1)$ there is a line $\lambda_\tau\ni p^*$ with $Z_\tau=\lambda_\tau\cap\Di$; $A_\tau$ and $B_\tau$ are the two open sides of $\lambda_\tau$ in $\Di$. If $n_\tau$ denotes the unit normal of $\lambda_\tau$ in $\one^\perp$ pointing into $A_\tau$, then $\tau\mapsto n_\tau$ is continuous and strictly monotone in angle, $\lambda_{\tau_i}=\lin(p^*,e_i)$ at each crossing, and $n_\tau\cdot(e_l-p^*)>0$ whenever $h(\tau)_l>0$. Every cluster line of $h(\tau)$ equals $\lambda_\tau$, with the same orientation.
\end{theorem}
\begin{proof}[Proof (steps in the order required by \src{N18} (R-a) and \src{N20} \S4)]
\emph{Step 0.} By Lemma~\ref{lem:UVN}, $A_\tau\supseteq V_l\cap\Di$ when $h(\tau)_l>0$ and $B_\tau\supseteq V_k\cap\Di$; by Lemma~\ref{lem:levels} at level $0\in I_{h(\tau)}$, $A_\tau$ and $B_\tau$ are proper open half-planes intersected with $\Di$, with boundary lines $\lambda^\pm_\tau$, and $Z_\tau$ lies between them.

\emph{Step 1 ($p^*\in\lambda^+_\tau$).} If $\tau<\tau_j$ take $\sigma_1\in(0,\min(\tau,\tau_i))$, $\sigma_2\in(\max(\tau,\tau_i),\tau_j)$; if $\tau\ge\tau_j$ take $\sigma_1\in(\tau_i,\min(\tau,\tau_j))$, $\sigma_2\in(\max(\tau,\tau_j),1)$. In both cases $[\sigma_1,\sigma_2]$ contains a crossing, and $h(\sigma_1),h(\sigma_2)$ share a positive coordinate $l$ (Lemma~\ref{lem:signs}). Then $G\coloneqq\operatorname{int}(P_{\sigma_1}\cap P_{\sigma_2})\supseteq V_l(\sigma_1)\cap V_l(\sigma_2)\cap\Di$ is nonempty, open and convex, and $p^*\in\cl G$. Lemma~\ref{lem:nopin} gives $G\subseteq A_\tau$, so $p^*\in\cl A_\tau\setminus A_\tau=\lambda^+_\tau\cap\Di$. A line through $p^*$ and $e_l$ with $h(\tau)_l>0$ would cut $V_l\cap\Di\subseteq A_\tau$, so $n_\tau\cdot(e_l-p^*)>0$.

\emph{Step 2 (monotone pieces).} If $\sigma<\tau<\sigma'$ and $A_\sigma\cap A_{\sigma'}\ne\varnothing$, concavity in $\tau$ gives $A_\sigma\cap A_{\sigma'}\subseteq A_\tau$; near $p^*$ this is an inclusion of wedges, so $n_\tau\in\operatorname{cone}(n_\sigma,n_{\sigma'})$. On $(0,\tau_j)$ all $A_\tau$ contain points near $e_j$, so the $n_\tau$ lie in one open half-circle and betweenness makes the angle monotone there; likewise on $(\tau_i,1)$.

\emph{Step 3 (continuity).} If the one-sided limits $n(\tau_0\pm)$ differed, then on the open wedge $W\coloneqq\{n(\tau_0+)\cdot\delta>0>n(\tau_0-)\cdot\delta\}$ and on $-W$ we would have $\Psi(p^*+\delta,\tau_0)=0$ for small $\delta$. The zero set is convex (Lemma~\ref{lem:levels}), so it would contain a neighbourhood of $p^*$, contradicting Step 1.

\emph{Step 4 (no constancy; crossing values).} If $n_\tau$ were constant on some $[a,b]$, with line $\lambda$, then $\Psi(q,\cdot)=0$ on $[a,b]$ for $q\in\lambda\cap\Di$, so $\Psi(q,\cdot)\le0$ on $\R$ (Lemma~\ref{lem:caps}(a)). Hence $\lambda\cap\Di\subseteq\{\Psi(\cdot,\tau)\le0\}$ for every $\tau$, which forces $\lambda^+_\tau=\lambda$ for all $\tau\in(0,1)$; by continuity $n_\tau$ would then be constant on $(0,1)$. But $n_\tau\cdot(e_i-p^*)$ is $>0$ where $h_i>0$ (Step 1) and $<0$ where $h_i<0$ (there $V_i\cap\Di\subseteq B_\tau$, which lies in the open negative side of $\lambda^+_\tau$). So the angle is non-constant on every interval, the two monotone pieces glue to one strictly monotone continuous angle on $(0,1)$, and continuity of $n$ and the sign change give $\lambda_{\tau_i}=\lin(p^*,e_i)$.

\emph{Step 5 (thin zero set).} If $\lambda^-_{\tau_0}\ne\lambda^+_{\tau_0}$, then $\lambda^-_{\tau_0}\not\ni p^*$, and a half-disc $D^-$ at $p^*$ on the negative side of $\lambda^+_{\tau_0}$ lies in $Z_{\tau_0}$. Pick $\tau_1\ne\tau_0$ near $\tau_0$ and $q\in A_{\tau_1}\cap D^-$. Then $\Psi(q,\tau_1)>0=\Psi(q,\tau_0)$, so by Lemma~\ref{lem:caps}(a) $q\in A_\tau$ for $\tau$ tending to $\tau_0$ from one side, and continuity of $n$ gives $n_{\tau_0}\cdot(q-p^*)\ge0$, contradicting $q\in D^-$.

\emph{Cluster lines.} A cluster line $\ell$ of $h(\tau)$ must meet $\Di$, since otherwise $\Psi(\cdot,\tau)$ would be one-signed on $\Di$. By Lemma~\ref{lem:chords}(b) its closed positive side misses $B_\tau$ and its closed negative side misses $A_\tau$, so $\ell\cap\Di\subseteq Z_\tau$ and $\ell=\lambda_\tau$ with the same orientation.
\end{proof}
\status{proved; refereed twice (\src{N18} \S1.5, \src{N20} \S4; CORRECT WITH REPAIR, repairs (R-a), (R-b), (R-c) incorporated). Source: \src{N15} Thm.~4.2. It excludes pinning of a vertex on any arc, non-unique clusters, chord jumps and fat zero sets.}

\begin{theorem}\label{thm:nosliding}
Let $r$ be the common corner ray of $F_k$ and $F_j$, parametrised so that $r=h(0)$ on both facets. Then the one-sided limits $n^k(0+)$ and $n^j(0+)$ of the chord normals coincide.
\end{theorem}
\begin{proof}
If not, $P(\cdot,r)$ is $\ge0$ and $\le0$ on two opposite open wedges at $p^*$, and by convexity of its zero set (Lemma~\ref{lem:levels}; $r\in U$ is mixed) it vanishes on a ball $B$ around $p^*$. By strict monotonicity on $F_k$ there is $q\in B$ with $n^k(0+)\cdot(q-p^*)<0$ and $q\in A^k_\tau$ for some $\tau\in(0,1)$. By continuity at $0+$, $q\notin A^k_\sigma$ for $\sigma\in(0,\eta)$, some $\eta>0$. So $\{\tau:\Psi^k(q,\tau)>0\}$ is an interval that contains $\tau$ and misses $(0,\eta)$; Lemma~\ref{lem:caps}(a) then gives $P(q,r)=\Psi^k(q,0)<0$, a contradiction.
\end{proof}
\status{proved; refereed twice (\src{N18} \S1.6, \src{N20} \S4). Source: \src{N15} Thm.~4.3. This excludes ``sliding edges''.}

\subsection{Magnitudes from rotation}

Fix a facet and orient $\tau$ so that the angle $\theta_\tau$ of $n_\tau$ is continuous and strictly increasing. Put $u_\tau\coloneqq$ the rotation of $n_\tau$ by $+\pi/2$ in $\one^\perp$, $D^+_\tau(x)\coloneqq\Psi(p^*+xn_\tau,\tau)$, $D^-_\tau(x)\coloneqq-\Psi(p^*-xn_\tau,\tau)$, and let $L_*(\tau_1)$, $L^*(\tau_1)$ be the lower and upper left Dini derivatives of $\theta$ at $\tau_1$. For $0<\xi<r_0/2$ the \emph{anchor} $q_\tau(\xi)\coloneqq p^*+\xi u_\tau$ lies on $\lambda_\tau$, in $A_\sigma$ for $\sigma>\tau$ near $\tau$ and in $B_\sigma$ for $\sigma<\tau$.

\begin{lemma}\label{lem:rotation}
Let $\tau_1\in(0,1)$, $0<\xi<r_0/2$, $S^\pm\coloneqq\partial^\pm_\tau\Psi(q_{\tau_1}(\xi),\tau_1)$. Then $0<S^+\le S^-<\infty$ and
\[
\frac{S^+}{L^*}\cdot\frac{r_0-\xi}{r_0\xi}\ \le\ \liminf_{x\downarrow0}\frac{D^+_{\tau_1}(x)}x\ \le\ \limsup_{x\downarrow0}\frac{D^+_{\tau_1}(x)}x\ \le\ \frac{S^-}{L_*}\cdot\frac{r_0+\xi}{r_0\xi},
\]
with $1/0=\infty$, $1/\infty=0$. The same holds for $D^-_{\tau_1}$ with $|\partial^-_\tau\Psi(q',\tau_1)|$ on the left and $|\partial^+_\tau\Psi(q',\tau_1)|$ on the right, $q'=p^*-\xi u_{\tau_1}$.
\end{lemma}
\begin{proof}[Proof sketch]
For $\tau_0<\tau_1$ near $\tau_1$ let $\Delta\theta\coloneqq\theta_{\tau_1}-\theta_{\tau_0}>0$ and $q\coloneqq q_{\tau_0}(\xi)$. Relative to $\lambda_{\tau_1}$, $q$ has transverse coordinate $\xi\sin\Delta\theta$ and tangential coordinate $\xi\cos\Delta\theta$, so Lemma~\ref{lem:sandwich} gives $D^+(\xi\sin\Delta\theta\,\tfrac{r_0}{r_0+\xi})\le\Psi(q,\tau_1)\le D^+(\xi\sin\Delta\theta\,\tfrac{r_0}{r_0-\xi})$. Since $\Psi(q,\cdot)$ vanishes at $\tau_0$ and is positive just to its right, Lemma~\ref{lem:caps}(b) gives $(\tau_1-\tau_0)\Psi(q,\tau_2)/(\tau_2-\tau_0)\le\Psi(q,\tau_1)\le(\tau_1-\tau_0)\,\partial^+_\tau\Psi(q,\tau_0)$ for $\tau_2>\tau_1$. As $\tau_0\uparrow\tau_1$, upper semicontinuity of $\partial^-$ (Lemma~\ref{lem:caps}(c)) bounds $\limsup\partial^+\Psi(q_{\tau_0},\tau_0)$ by $S^-$, and joint continuity followed by a supremum over $\tau_2\downarrow\tau_1$ gives $S^+$ on the left; $(\tau_1-\tau_0)/\sin\Delta\theta$ has limsup $1/L_*$ and liminf $1/L^*$, and every small $x$ is of the form $\xi\sin\Delta\theta\,r_0/(r_0\pm\xi)$. Positivity and finiteness of $S^\pm$ come from Lemma~\ref{lem:caps}(b), since $\Psi(q_{\tau_1},\cdot)$ changes sign strictly at $\tau_1$ (Theorem~\ref{thm:pivot}).
\end{proof}
\status{proved; refereed twice (\src{N18} \S1.7, \src{N20} \S5; both verified the semicontinuity directions analytically, noting that gauge-mixture numerics cannot test them). Source: \src{N15} Lemma 5.1.}

\begin{lemma}\label{lem:risingsun}
If $\theta$ is continuous and strictly increasing on $[a,b]$, then $L_*(\tau)>0$ for some $\tau\in(a,b]$.
\end{lemma}
\status{proved (elementary); refereed twice. Source: \src{N15} Lemma 5.2.}

\begin{proposition}\label{prop:regularity}
\begin{enumerate}[leftmargin=*,label=(\alph*)]
\item For every compact $J\subset(0,1)$ there are $C,t_0$ with $|\Psi(p^*+\delta,\tau)|\le C|\delta|$ for $\tau\in J$, $|\delta|\le t_0$.
\item $\theta'(\tau)$ exists and lies in $(0,\infty)$ for a.e.\ $\tau$; call this full-measure set $E'$.
\item For $\tau\in E'$ there is $c_\tau>0$ with $c_\tau x\le D^\pm_\tau(x)\le Cx$ for small $x$.
\item If $\lim_{x\downarrow0}D^\pm_\tau(x)/x$ exist, then $\lim_{t\downarrow0}\Psi(p^*+t\delta,\tau)/t$ exists for every $\delta$ and equals $\alpha[n_\tau\cdot\delta]_+-\beta[n_\tau\cdot\delta]_-$.
\end{enumerate}
\end{proposition}
\begin{proof}[Proof sketch]
(a) Lemma~\ref{lem:risingsun} gives points $a<b$ left of $J$ and $c$ right of $J$ with $L_*>0$; Lemma~\ref{lem:rotation} bounds $D^\pm$ linearly there, Lemma~\ref{lem:sandwich} extends the bound to $|\Psi(p^*+\delta,\sigma)|\le2C'|\delta|$ for $\sigma\in\{a,b,c\}$, and concavity in $\tau$ (extrapolation from the left, minimum of endpoints) transfers it to $J$. (b) Lebesgue's theorem for monotone $\theta$; $\theta'(\tau)=0$ would make the lower bound of Lemma~\ref{lem:rotation} infinite, contradicting (a). (c) Lemma~\ref{lem:rotation} with $L_*=L^*=\theta'$. (d) Lemma~\ref{lem:sandwich} at scale $t$.
\end{proof}
\status{proved; refereed twice (\src{N18} \S1.8, \src{N20} \S5.3). Source: \src{N15} Prop.~5.3. Part (b) also excludes singular (Cantor-type) chord maps. Oscillation of $D^\pm_\tau(x)/x$ inside $[c_\tau,C]$ is \emph{not} excluded locally; \S\ref{sec:proj} shows it is irrelevant.}

\subsection{Projectivity}\label{sec:proj}

\begin{lemma}\label{lem:twoslope}
Let $I$ be an open interval, $n:I\to\one^\perp$ continuous with $|n|=1$ and not constant on $I$, and $\alpha,\beta:I\to(0,\infty)$ continuous. Suppose that for every $\delta\in\one^\perp$ the function $f_\delta(\tau)\coloneqq\alpha(\tau)[n(\tau)\cdot\delta]_+-\beta(\tau)[n(\tau)\cdot\delta]_-$ is concave on $I$. Then $a\coloneqq\alpha n$ is affine on $I$ and $\beta/\alpha$ is constant.
\end{lemma}
\begin{proof}[Proof sketch]
On a small interval $J$ with $n(J)$ in an arc of length $<\pi/2$, let $C_J\coloneqq\{\delta:n(\tau)\cdot\delta>0\ \forall\tau\in J\}$. For $\delta\in C_J$, $f_\delta=a\cdot\delta$ is concave; for $\delta\in-C_J$, $f_\delta=b\cdot\delta$ with $b\coloneqq\beta n$ is concave. So $a,b$ are locally Lipschitz, and their distributional second derivatives are vector measures with $a''=-n\,m_a$, $b''=n\,m_b$, $m_a,m_b\ge0$ (using continuity of $n$). With $\times$ the determinant on $\one^\perp$, $d(a\times a')=a\times a''=0$ because $a\parallel n$; so $a\times a'\equiv L_a$ and likewise $b\times b'\equiv L_b$ are constant on $I$. Now $a\times a'=\alpha^2(n\times n')$ and $b\times b'=\beta^2(n\times n')$. If $L_a=0$ then $n'=0$ a.e.\ and $n$ is constant, which is excluded. So $(\beta/\alpha)^2=L_b/L_a$ is constant, $b=\rho a$, and $(m_b+\rho m_a)n=b''-\rho a''=0$ forces $m_a=0$: $a$ is affine.
\end{proof}
\status{proved; refereed once (\src{N14} \S2.6: CORRECT; regularity is derived from concavity and atoms of $a''$ are handled); its use in \src{N15} checked by both \src{N18} and \src{N20}. Source: \src{N12} Lemma 4.3.}

\begin{proposition}\label{prop:blowup}
Fix a facet and a compact interval $J_1\subset(0,1)$. There are $t_m\downarrow0$ and continuous $\alpha,\beta:J_1\to(0,\infty)$ such that, with $\Psi_t(\delta,\tau)\coloneqq\Psi(p^*+t\delta,\tau)/t$, the limits
\[
G(\delta,\tau)\coloneqq\lim_m\Psi_{t_m}(\delta,\tau),\qquad G_2(\delta,\tau)\coloneqq\lim_{\epsilon\downarrow0}G(\epsilon\delta,\tau)/\epsilon
\]
exist for all $(\delta,\tau)\in\one^\perp\times J_1$, $G_2(\delta,\tau)=\alpha_\tau[n_\tau\cdot\delta]_+-\beta_\tau[n_\tau\cdot\delta]_-$, and $\tau\mapsto G_2(\delta,\tau)$ is concave for every $\delta$.
\end{proposition}
\begin{proof}[Proof sketch]
Take $J_1\Subset J\Subset J'\Subset(0,1)$. By Proposition~\ref{prop:regularity}(a) and concavity, $\Psi_t(\delta,\cdot)$ is $L|\delta|$-Lipschitz on $J$ uniformly in $t$; and (Lemma~\ref{lem:sandwich} at scale $t$) $\Psi_t(\delta,\tau)\le\Psi_t(\delta',\tau)$ when $n_\tau\cdot\delta<n_\tau\cdot\delta'$ and $t$ is small (property (M)).

\emph{First blow-up.} Arzel\`a--Ascoli on a countable dense set of $\delta$ and a diagonal argument give $t_m$ and a limit $G$, concave and Lipschitz in $\tau$, that inherits (M).

\emph{Sweep.} Let $\tau_0\in E'\cap\operatorname{int}J$, $\vartheta\coloneqq\theta'(\tau_0)\in(0,\infty)$. Suppose the profile $x\mapsto G(\cdot,\tau_0)$ along $n_{\tau_0}$ had a jump $\mathcal J>0$ at $x_0$. Take a far point $\delta\approx Yu_{\tau_0}+x_0n_{\tau_0}$. Rotating the chord through an angle $\pm3\epsilon/Y$, i.e.\ moving $\tau$ by about $\pm3\epsilon/(Y\vartheta)$, moves $n_\sigma\cdot\delta$ across the jump. The Lipschitz bound in $\tau$ then gives $\mathcal J-o(1)\le L|\delta|\cdot6\epsilon/(Y\vartheta)\approx6L\epsilon/\vartheta$, so $\mathcal J=0$. Hence $G(\delta,\tau_0)=\varphi_{\tau_0}(n_{\tau_0}\cdot\delta)$ with $\varphi_{\tau_0}$ continuous and nondecreasing; (M) and equi-Lipschitz convergence on the dense set $E'\cap J$ extend the convergence to all $(\delta,\tau)$.

\emph{Window.} With $\delta_Y\coloneqq Yu_{\tau_0}+x_0n_{\tau_0}$, the concave functions $s\mapsto G(\delta_Y,\tau_0+s/(Y\vartheta))$, $|s|\le r$, converge uniformly to $s\mapsto\varphi_{\tau_0}(x_0+s)$: a single far point sees every transverse position as the chord turns through an angle of order $1/Y$. So $\varphi_{\tau_0}$ is concave on $\R$.

\emph{Second zoom.} Concavity and $\varphi(0)=0$ give $\alpha_{\tau_0}=\varphi'(0+)$ and $\beta_{\tau_0}=\varphi'(0-)$, with $c_{\tau_0}\le\alpha\le\beta\le C$ (Proposition~\ref{prop:regularity}(c)); equi-Lipschitz convergence of $G(\epsilon\delta,\cdot)/\epsilon$ on $E'\cap J_1$ extends the two-slope form to all of $J_1$ with $\alpha,\beta$ continuous, and $\alpha>0$ everywhere because a concave function $\ge0$ vanishing at an interior point vanishes identically.
\end{proof}
\status{proved; refereed twice (\src{N18} \S1.9, \src{N20} \S6; cosmetic repair $\tau_0\in\operatorname{int}J$). Source: \src{N15} \S6 Steps 1--4.}

Say that \textbf{(PC)} holds if for each facet $F_k$ there is a linear map $\Gamma_k:v_k^\perp\to\one^\perp$ with $\Gamma_kh\ne0$ on $\relint F_k$, such that every cluster line of $h\in\relint F_k$ is $\{q:\Gamma_kh\cdot(q-p^*)=0\}$ with positive side $\{\Gamma_kh\cdot(q-p^*)>0\}$. Say that \textbf{(NS)} holds if $\Gamma_kr\propto_+\Gamma_jr$ at each corner ray $r=F_k\cap F_j$.

\begin{corollary}\label{cor:PCNS}
(PC) and (NS) hold, and each $\Gamma_k$ has rank $2$.
\end{corollary}
\begin{proof}
By Proposition~\ref{prop:blowup} and Lemma~\ref{lem:twoslope} ($n$ is nowhere locally constant by Theorem~\ref{thm:pivot}), $a=\alpha n$ is affine on each compact $J_1\subset(0,1)$. Positively proportional affine maps with non-constant direction differ by a constant factor (\src{N20} \S6.5), so the normalisations patch and $a$ extends affinely to $[0,1]$; define $\Gamma_kh(\tau)\coloneqq a(\tau)$ and extend linearly to $v_k^\perp=\operatorname{span}\{r,r'\}$. The direction field $n_\tau$ is the zero-chord field and does not depend on the subsequence. With Theorem~\ref{thm:pivot} this is (PC). If $a(0)$ and $a(1)$ were parallel the direction would be constant, and if antiparallel $a$ would vanish inside, so the rank is $2$ and $\Gamma_kr=a(0)\ne0$. Finally $\Gamma_kr\propto_+n^k(0+)=n^j(0+)\propto_+\Gamma_jr$ by Theorem~\ref{thm:nosliding}, which is (NS).
\end{proof}
\status{proved; refereed twice (\src{N18} \S1.9--1.10, \src{N20} \S\S6.5, 7). Source: \src{N15} \S6 Step 5 and Thm.~7.1.}

\subsection{The contradiction}\label{sec:contradiction}

\begin{lemma}\label{lem:edgebox}
Let $a\ne b$. For a.e.\ pair $(A,B)\in U^2$ with $A_a<0<A_b$, $B_b<0<B_a$ and $A_bB_a<|A_a|\,|B_b|$, the closed edge $[e_a,e_b]$ does not meet $H(A)\cap H(B)$.
\end{lemma}
\begin{proof}
The ratio condition gives an open interval of $\lambda>0$ on which $(A+\lambda B)_a<0$ and $(A+\lambda B)_b<0$; by Fubini, for a.e.\ $(A,B)$ some such $A+\lambda B$ is good. Suppose $q_0\in[e_a,e_b]\cap H(A)\cap H(B)$. The open set $O\coloneqq H(A)\cap H(B)\cap\Di$ accumulates at $q_0$, and on $O$, $P(\cdot,A+\lambda B)\ge P(\cdot,A)+\lambda P(\cdot,B)>0$ a.e. If $A+\lambda B\notin U$ this contradicts $P\le P_{K'}\le0$; otherwise $O\subseteq H(A+\lambda B)$ and $q_0\in\cl H(A+\lambda B)$. But $c(A+\lambda B)$ is negative in coordinates $a$ and $b$, hence negative on the whole closed edge.
\end{proof}
\status{proved; refereed once (\src{N14} \S2.3: CORRECT, with an explicit sanity test). Source: \src{N12} Thm.~2.4. It needs no regularity and no coherence at edge priors.}

Under (PC) put $c_k(h)\in\R^3$ for the vector with $c_k(h)\cdot p^*=0$ whose restriction to $\one^\perp$ is $\Gamma_kh$; then the cluster line of $h$ is $\{c_k(h)\cdot q=0\}$, a line in all of $\aff\Delta$.

\begin{proposition}\label{prop:boxbridge}
Assume (PC), (NS) and the conclusions of Theorem~\ref{thm:pivot}.
\begin{enumerate}[leftmargin=*,label=(\alph*)]
\item \emph{Crossing rows.} For each crossing coordinate $i$ of $F_k$ there is $\lambda_{ik}>0$ with $c_k(h)_i=\lambda_{ik}h_i$ for all $h\in v_k^\perp$.
\item \emph{Corners.} $c_k(r)\propto_+c_j(r)$ at each corner ray $r=F_k\cap F_j$.
\item \emph{Box.} Let $A\in\relint F_k$, $B\in\relint F_{k'}$ and $q_0\in[e_a,e_b]$ satisfy $A_a<0<A_b$, $B_b<0<B_a$, $A_bB_a<|A_a||B_b|$. Then $c_k(A)\cdot q_0\le0$ or $c_{k'}(B)\cdot q_0\le0$.
\end{enumerate}
\end{proposition}
\begin{proof}
(a) At the crossing $w_{ik}$ the unique cluster line is $\lin(p^*,e_i)$, so $c_k(w_{ik})\cdot e_i=0$; the functional $h\mapsto c_k(h)_i$ on the plane $v_k^\perp$ vanishes on $\R w_{ik}=v_k^\perp\cap e_i^\perp$, hence equals $\lambda_{ik}h_i$, and $\lambda_{ik}>0$ because $c_k(h)\cdot e_i=\Gamma_kh\cdot(e_i-p^*)>0$ when $h_i>0$ (Theorem~\ref{thm:pivot}). (b) is (NS). (c) Suppose $c_k(A)\cdot q_0>0$ and $c_{k'}(B)\cdot q_0>0$. By uniqueness of cluster lines, the normalised vectors $c(A')$ of good $A'\to A$ converge to the normalised $c_k(A)$ (an affine function on $\aff\Delta$ is the restriction of a unique linear functional on $\R^3$), so $q_0\in H(A')$ for a.e.\ $A'$ in a small ball around $A$; likewise for $B$. The sign and ratio conditions are open. This produces a set of pairs of positive measure contradicting Lemma~\ref{lem:edgebox}.
\end{proof}
\status{proved; refereed once (\src{N14} \S2.5: CORRECT; nondegeneracy at the $b$-crossings from \src{R18} \S5). Source: \src{N12} Thm.~3.1 items 3--6. Its use with the conclusions of \S4.3 was checked in \src{N18} and \src{N20}.}

\paragraph{The contamination triangle.}
Take $T=T_1$ and $p^*=u$. On $F_k$ the crossing coordinates are the two others, and $c_k(h)\cdot u=0$ gives $c_k(h)_k=-\sum_{i\ne k}\lambda_{ik}h_i$. Normalise $\lambda_{31}=\lambda_{12}=\lambda_{23}=1$ and put $x_1\coloneqq\lambda_{21}$, $x_2\coloneqq\lambda_{32}$, $x_3\coloneqq\lambda_{13}$. Then
\begin{gather*}
c_1(h)\propto(-(x_1h_2+h_3),\ x_1h_2,\ h_3),\qquad
c_2(h)\propto(h_1,\ -(h_1+x_2h_3),\ x_2h_3),\\
c_3(h)\propto(x_3h_1,\ h_2,\ -(x_3h_1+h_2)).
\end{gather*}

\begin{proposition}\label{prop:T1}
For $T_1$, (PC), (NS) and the box of Proposition~\ref{prop:boxbridge}(c) are inconsistent, for every $\varepsilon\in(0,1)$.
\end{proposition}
\begin{proof}
\emph{Corner equations.} At $r_{12}=(-1,-1,\kappa+1)$: $c_1(r)=(x_1-\kappa-1,\,-x_1,\,\kappa+1)$ and $c_2(r)=(-1,\,1-x_2(\kappa+1),\,x_2(\kappa+1))$. These are positively proportional iff $c_1(r)=x_2^{-1}c_2(r)$, i.e.\ iff $x_1+1/x_2=\kappa+1$. Cyclically,
\begin{equation}\label{eq:corner}
x_1+\frac1{x_2}=x_2+\frac1{x_3}=x_3+\frac1{x_1}=\kappa+1 .
\end{equation}
With $\phi(t)\coloneqq\kappa+1-1/t$ this says $x_1=\phi(x_2)$, $x_2=\phi(x_3)$, $x_3=\phi(x_1)$, so $x_1$ is a fixed point of $\phi^3$. The M\"obius map $\phi$ has matrix $\bigl(\begin{smallmatrix}\kappa+1&-1\\1&0\end{smallmatrix}\bigr)$ with determinant $1$ and trace $\kappa+1>2$, so it is hyperbolic and $\phi^3$ has the same two fixed points $x_\pm$, the roots of $x^2-(\kappa+1)x+1=0$. Hence
\[
x_1=x_2=x_3=x\in\{x_+,x_-\},\qquad x+\frac1x=\kappa+1,\qquad x_+>1,\ x_-=1/x_+ .
\]

\emph{Box.} Let $\beta\coloneqq1/\kappa$, $q_0\coloneqq(s,1-s,0)$, and for small $\eta,\zeta>0$
\[
A=(-\beta(1+\eta),\,\eta,\,1)\in\relint F_1,\qquad B=(1,\,-\zeta,\,-\beta(1-\zeta))\in\relint F_3 ,
\]
near the two crossings of coordinate $2$ (membership uses $c_0(1-\beta)=\varepsilon\beta$). Then $c_1(A)\cdot q_0=x\eta(1-2s)-s$ and $c_3(B)\cdot q_0=sx-(1-s)\zeta$, and the ratio condition for $(a,b)=(1,2)$ is $\eta<\beta\zeta(1+\eta)$. Take $\zeta\coloneqq\theta sx/(1-s)$ and $\eta\coloneqq s/(\theta x(1-2s))$ with $\theta<1$. Both chart values are positive, and the ratio condition reads $(1-s)/((1-2s)(1+\eta))<\theta^2\beta x^2$. If $x^2>\kappa$, this holds for $\theta$ close to $1$ and $s$ small, contradicting Proposition~\ref{prop:boxbridge}(c). So $x^2\le\kappa$.

The transposition $\sigma$ of coordinates $1,2$ maps $T_1$ to itself and fixes $u$; $c^\sigma(g)\coloneqq\sigma c(\sigma g)$ is again a realisation, so the whole chain applies to it, and its facet data are $x^\sigma=1/x$. Hence also $x^{-2}\le\kappa$, i.e.
\[
\frac1\kappa\le x^2\le\kappa .
\]

\emph{Contradiction.} Since $t\mapsto t+1/t$ is increasing on $[1,\infty)$ and invariant under $t\mapsto1/t$, the box gives $x+1/x\le\sqrt\kappa+1/\sqrt\kappa$. With $y=\sqrt\kappa>1$,
\[
(\kappa+1)-\Bigl(y+\frac1y\Bigr)=\frac{(y-1)(y^2+1)}{y}>0,
\]
which contradicts $x+1/x=\kappa+1$. Concretely, $x_+^2>\kappa$ violates the upper bound and $x_-^2<1/\kappa$ the lower one; at $\varepsilon=0.3$, $x_+^2/\kappa\approx3.80$.
\end{proof}
\status{proved. The corner system is \src{R15} Prop.~5.2 and \src{R4} Prop.~G1 (refereed in \src{R18} and \src{R8}); the box value $1/\kappa\le x^2\le\kappa$ and the certificate are \src{R15} \S6 (refereed in \src{R18} \S6, which recomputed them by hand). \src{N20} \S7 confirms that this contradiction holds for every $\kappa>1$, i.e.\ every $\varepsilon\in(0,1)$. The transposition step replaces the explicit mirror certificate near $e_1$ used in \src{R15}; this is a change of presentation only. The corner equations alone are consistent (two sign-admissible solutions $x_\pm$), so the box is indispensable. The same equations, with the same contradiction, appear for continuous maps in \src{CONTAMINATION\_NOGO\_SCOPE\_AND\_GEOMETRY.md} \S2.}

\paragraph{General all-facing triangles.}
For an all-facing triangle in general position the same three ingredients give a contradiction by the algebra of \src{R15} Thm.~6.2. Work at $p^*=u$ with $\mu_{ik}=\lambda_{ik}/3$, write $v_{kl}$ for the $l$-th coordinate of $v_k$ and $\rho_k(i{:}j)\coloneqq\mu_{ik}/\mu_{jk}$. At the corner $r=r_{ac}$ with third index $b$, Lemma~\ref{lem:signs} gives $r_a,r_c<0<r_b$; put $A_r\coloneqq-r_c/r_b$, $B_r\coloneqq-r_a/r_b$, $s\coloneqq\rho_a(b{:}c)$, $t\coloneqq\rho_c(a{:}b)$. Then (NS) reads $A_r/s+B_rt=1$, and the boxes for the pairs $(a,b)$ and $(c,b)$ read $v_{ca}/v_{cc}\le st\le v_{aa}/v_{ac}$. Together with $r\perp v_a,v_c$ these give
\[
\frac{v_{cb}}{v_{cc}}\le\rho_a(b{:}c)\le\frac{v_{ab}}{v_{ac}},\qquad\frac{v_{ca}}{v_{cb}}\le\rho_c(a{:}b)\le\frac{v_{aa}}{v_{ab}} .
\]
Applying the second inequality at the other corner of facet $a$ pinches $\rho_a(b{:}c)=v_{ab}/v_{ac}$ on every facet. The corner equation then gives $t=v_{aa}/v_{ab}$, while pinching on facet $c$ gives $t=v_{ca}/v_{cb}$. So $v_{aa}v_{cb}=v_{ab}v_{ca}$, i.e.\ $\det(v_a,v_c,e_c)=0$, and $e_c$ lies on the line through $v_a,v_c$, contradicting general position.
\status{proved; refereed once (\src{R18} \S6: every step recomputed). In \src{R15} the box came from a forbidden-support lemma under extra hypotheses; here it comes from Proposition~\ref{prop:boxbridge}(c), which uses only chord lines (\src{N12} Thm.~3.1, \src{N14} \S2.5).}

\begin{proof}[Proof of Theorem~\ref{thm:main}]
Reduce to $p^*=u$ (Proposition~\ref{prop:reduction}(d)). Theorem~\ref{thm:pivot}, Theorem~\ref{thm:nosliding} and Corollary~\ref{cor:PCNS} give unique pivot chords, (PC) and (NS). Proposition~\ref{prop:boxbridge} turns these into crossing rows, corner equations and the box. Proposition~\ref{prop:T1} (for $T_1$) and the general pinching argument give a contradiction.
\end{proof}

\section{General decomposable domains}\label{sec:general}

Call a decomposable domain \emph{two-sided} if, near the target, evaluable models may accept $f$, accept $-f$, or give no advice: a.e.\ pair with an orientation near $\partial D_T$ has an advice range containing $\{+,0,-\}$. Without non-adjacency, pairs near $\partial D_T$ whose advice range is $\{+,0,-\}$ may have wedge charts $\{c\cdot q>0,\,d\cdot q>0\}$. Non-adjacency is used in exactly two places in \S\ref{sec:main} (\src{N21} \S1, confirmed in \src{N23} \S1): (U-glob) global straight levels and Lemma~\ref{lem:PL2} on a set that leaves every ball around $p^*$, inside Lemma~\ref{lem:nopin}; and (U-box) the box step, where the witnesses' actual charts must contain the edge prior. Everything else is local at $p^*$ or adjacency-tolerant. Throughout this section the ranking is strictly guiding additive on a decomposable domain with (LF) at $T$, and (BU) holds at every interior prior with $K_{p^*}=T$ (R4 notation: for a pair with advice range $\{+,0,-\}$, $R(g)=\{\alpha>0,d>0\}$, $R(-g)=\{\beta>0,d<0\}$ with $\alpha\coloneqq c\cdot q$, $\beta\coloneqq c'\cdot q$).

\paragraph{The neighbourhood $N_0$ (compact collar).} (LF) supplies only an open $N$, which need not contain a conic neighbourhood of $\partial D_T\setminus0$. Every use of a neighbourhood in this section concerns gambles near one fixed compact set of boundary gambles at bounded norm (\src{N27} \S1.1, inventory re-derived independently in \src{N27\_referee} \S1). Let $\Sigma\subseteq\partial D_T\setminus0$ be the compact set formed by the three facet segments $h([0,1])$ between the normalised corner rays, the crossings in every normalisation used, and the limits of the box witnesses of Proposition~\ref{prop:boxbridge}(c) in the normalisation of Proposition~\ref{prop:T1} (for $T_1$: $(-\beta,0,1)$ and $(1,0,-\beta)$) and of the test gambles of Theorem~\ref{thm:cuts}. Fix a compact $K''\subseteq\operatorname{int}T$ and $\eta>0$ such that the compact collar
\[
\Nb\coloneqq\{g:\operatorname{dist}(g,\Sigma)\le\eta\}
\]
lies in $N$, has $P_{K''}>0$ on it, misses $-\cl D_T$ and contains no $h\le0$; this is possible because $\Sigma$ is compact, $N$ is open and $P_{K''}>0$ on $\Sigma$. Below, $N_0$ means $\Nb$. Two numerical conditions are imposed:
\begin{enumerate}[leftmargin=*,label=(N\arabic*)]
\item level shifts are restricted to $|c|\le\bar c$ with $\bar c<\eta/\sqrt3$, so that $h(\tau)-c\one\in\operatorname{int}\Nb$ for $\tau\in[0,1]$;
\item the radius $r_1$ of Lemma~\ref{lem:B1} is chosen after $\bar c$, so small that $|\Psi_k(q,\tau)|<\bar c$ on $B(p^*,r_1)\times[0,1]$ for each facet (possible by joint continuity, since $\Psi_k(p^*,\cdot)=0$ on $[0,1]$).
\end{enumerate}
With these choices every statement of this section holds with the same proof (\src{N27} Repair 27.1; \src{N27\_referee} \S1.3: CONFIRMED, constants re-derived; \src{N28} Lemma 0 gives the same repair and was CONFIRMED WITH REPAIR, the repair being exactly (N1)--(N2)). The balls around crossings, the witness balls and the test-gamble neighbourhoods are taken inside $\operatorname{int}\Nb$, which is possible because their centres lie in, or converge to, $\Sigma$. For a.e.\ gamble $h$ of the open cone $\R_+N$, accepting $h$ lies in the advice range (\src{N28} Lemma 0(b), confirmed). The $U$-domain and Corollary~\ref{cor:coverage} are unaffected: they use the explicit cones $U$ and $D_{K'}$, never $N$ (\src{N27} \S1.5, confirmed).

\begin{lemma}\label{lem:raytransfer}
Let an additive ranking (strictly guiding, or with permitted ties as in \S\ref{sec:open-ties}) on a decomposable domain have, for every $q\in\Di$, accepted set $B_q\coloneqq\{h:q\in R(h)\}$ equal a.e.\ to $D_{K_q}\cup\{h\ge0\}$. Then:
\begin{enumerate}[leftmargin=*,label=(\alph*)]
\item for a.e.\ $(g,t)\in\R^3\times(0,\infty)$, $R(tg)\cap\Di=R(g)\cap\Di$;
\item under strict guidance, for a.e.\ mixed $g$ and a.e.\ $t>0$, the advice ranges at $g$ and $tg$ contain the same acceptances; they can differ only in whether they contain $\times$, and whether they contain no advice ($0$) at pairs with advice range $\{+,0,-\}$ whose no-advice cell $\Di\setminus(R(g)\cup R(-g))$ has empty interior;
\item for a.e.\ $g$ in the cone over $\operatorname{int}\Nb$ and every open convex $C\subseteq\{q:K_q\cap K''\ne\varnothing\}$: $R(-g)\cap C=\varnothing$, no d-face of $g$ meets $C$, and $R(g)\cap C=\{c(g)\cdot q>0\}\cap C$.
\end{enumerate}
\end{lemma}
\begin{proof}[Proof sketch]
(a) For fixed $(q,t)$ the $g$-section of $\{(q,g,t):1[q\in R(g)]\ne1[q\in R(tg)]\}$ lies in $(B_q\sd D')\cup t^{-1}(B_q\sd D')$ with $D'\coloneqq D_{K_q}\cup\{h\ge0\}$ a cone, hence is null; Fubini, and two relatively open convex sets that agree a.e.\ are equal. (b) An acceptance in the advice range of a mixed pair has a nonempty chart, an acceptance outside it an empty chart. (c) Apply (a) to $g/t\in\Nb$ and Lemma~\ref{lem:B1}; a d-face point would have nearby points of $C$ in $R(-g)$.
\end{proof}
\status{proved (\src{N27} Lemma 27.2, Cor.~27.3, Cor.~27.3$'$); confirmed by one referee (\src{N27\_referee} \S2, which also supplied the full $c$-form in (c)). No guidance is used in (a). Only the charts transfer along rays, not the penalty vectors.}

\begin{lemma}\label{lem:quasiconcave}
For every $g\in\R^3$, $q\mapsto P(q,g)$ is quasi-concave on $\Di$: the sets $\{P(\cdot,g)>c\}$ and $\{P(\cdot,g)\ge c\}$ are convex for every $c\in\R$. Consequently, if $A_\tau\coloneqq\{\Psi(\cdot,\tau)>0\}$ satisfies $A_\tau\cap B=H^+\cap B$ for a ball $B$ at $p^*$ and an open half-plane $H^+$ whose boundary passes through $p^*$, then $A_\tau\subseteq H^+$.
\end{lemma}
\status{proved; refereed once (\src{N23} \S2). Source: \src{N21} Lemma N21.1 and Cor.~N21.1$'$. It uses only convexity of every chart in the table of charts for each advice range (\src{R1}) and the chord property of good gambles, and holds for every $g$, including nonnegative gambles and pairs on which every evaluable model gives the same advice (\src{N28\_referee} \S2). Non-adjacency is sufficient, not necessary, for quasi-convexity (\src{N23} counterexample: the precise identity system with adjacent data).}

\begin{lemma}\label{lem:B1}
There is $r_1>0$, satisfying (N2), with $K_q\supseteq K''$ for $q\in B_1\coloneqq B(p^*,r_1)$. For a.e.\ $g\in N_0$ and every $q$ with $K_q\cap K''\ne\varnothing$: $q\notin R(-g)$, $q\notin F(g)$, and $q\in R(g)$ iff $c(g)\cdot q>0$. So near-facet gambles have half-plane charts on $B_1$.
\end{lemma}
\status{proved; refereed once (\src{N23} \S3, which strengthened it to a pointwise statement). Source: \src{N21} Lemma N21.2.}

Say that \textbf{(TS)} holds if for every facet $F_k$ and every $\tau\in(0,1)$, $p^*\in\cl A_\tau\cap\cl B_\tau$.

\begin{theorem}\label{thm:TS}
On every decomposable domain with (LF), under (BU), (TP) and (TS), the conclusions of Theorems~\ref{thm:pivot} and~\ref{thm:nosliding}, Lemma~\ref{lem:rotation}, Propositions~\ref{prop:regularity} and~\ref{prop:blowup} hold with $\Di$ replaced by $B_1$ in every local statement and $r_0$ by $r_1$. In particular (NS) holds, and (PC) holds for $h$ near $\Sigma$, extended to the cone over $\Sigma$ by linearity of $\Gamma_k$.
\end{theorem}
\status{proved; refereed once (\src{N23} \S4: CORRECT, two wording repairs). Source: \src{N21} Thm.~N21.4. The one new argument is the negative vertex sign at crossings via clusters of $c$-vectors on $B_1$. The wording of (PC) follows \src{N27\_referee} \S1.4: the cluster clause is proved only near $\Sigma$, which is all that is used downstream; the full conic clause also follows from Lemma~\ref{lem:raytransfer}.}

Fix a pair $(a,b)$ with a failing box, i.e.\ with $\Lambda_{ab}>1$ for the facet data of \S\ref{sec:contradiction}; here $\Lambda_{ab}$ is \src{R15}'s vertex-limit box ratio, and witnesses near the two $b$-crossings contradict the box of Proposition~\ref{prop:boxbridge}(c) on $[e_a,e_b]$ near $e_b$ when $\Lambda_{ab}>1$. For $T_1$ every world $b$ has exactly one such $a$, since $\Lambda_{ab}\Lambda_{cb}=\kappa^{-2}<1$ and $\Lambda\in\{x_+^2/\kappa,\,1/(\kappa x_+^2)\}$ (exact; \src{N23} \S6).

\begin{theorem}\label{thm:cuts}
Assume (PC), (NS) and a realisation of $T_1$ on a decomposable domain with (LF). For every $b$ there are good gambles $g_m$ with advice range $\{+,0,-\}$ near one of the two $b$-crossings $w^*\in F_{k^*}$, with d-faces $F(g_m)\cap\Di=(a_m,z_m)$, such that $a_m\to a^*\in[e_b,p^*]\setminus B_1$, $z_m\to z^*\in[e_b,e_a)$, $K_q\subseteq\lin(e_b,v_{k^*})$ for $q\in[a^*,z^*]\cap\Di$, and some interior $q_m\to z^*$ have $K_{q_m}\to\{e_b\}$ in the Hausdorff metric.
\end{theorem}
\status{proved; refereed once (\src{N23} \S6). Source: \src{N21} Thm.~N21.6. Theorem~\ref{thm:apex} below sharpens the apex limit to $a^*=e_b$.}

\begin{corollary}\label{cor:availability}
Under the hypotheses of Theorem~\ref{thm:cuts}: accepting $h$ lies in the advice range for a.e.\ $h$ with $h_b>0$ that is not $\ge0$; the advice range of a.e.\ mixed pair contains both acceptances; $K_{e_b}=\{e_b\}$ for every $b$; and for every edge prior $q_0$ with support $S$, every limit of $K_q$ as $q\to q_0$ through $\Di$ lies in $\Delta_S$.
\end{corollary}
\status{proved; refereed once (\src{N23} \S6, wording repair). Source: \src{N21} Cor.~N21.7.}

\begin{lemma}\label{lem:R231}
Let $F_k$ have a crossing $\tau_c\in(\sigma_1,\sigma_2)$ and $0<\rho\le r_1$. If $\Psi(\cdot,\sigma_1)\ge0$ and $\Psi(\cdot,\sigma_2)\ge0$ on $B(p^*,\rho)$, then $v_k\in K_q$, $\Psi(q,\cdot)\equiv0$ on $[\sigma_1,\sigma_2]$ and $\Psi(q,\cdot)\le0$ on $[0,1]$ for every $q\in B(p^*,\rho)$.
\end{lemma}
\begin{proof}[Proof sketch]
As in Lemma~\ref{lem:nopin}, $B(p^*,\rho)\subseteq\Gamma_{v_k}$; Lemma~\ref{lem:B1} rules out d-faces there; Lemma~\ref{lem:PL2} applies to the ball; and now the $\alpha$-point is $p^*$ itself, since $K_{p^*}=T\ni v_k$. So the whole ball lies in $\alpha_{v_k}$.
\end{proof}
\status{proved in a referee report (\src{N23} \S7.1, Lemma R23.1). Second reading: CORRECT (recorded in \src{CHECK\_REPORT.md}); every step re-derived, including the use of Lemma~\ref{lem:PL2} on a ball that contains $p^*$ and the choice $\mathcal U\subseteq\operatorname{int}N_0$ needed for Lemma~\ref{lem:B1}.}

\begin{corollary}\label{cor:R232}
Failure of the $B$-half of (TS) alone is impossible: on each facet, if $p^*\in\cl A_\tau$ for every $\tau\in(0,1)$, then $p^*\in\cl B_\tau$ for every $\tau$. Hence (TS) is equivalent to its $A$-half (TS$_A$).
\end{corollary}
\status{proved in a referee report (\src{N23} \S7.2, Cor.~R23.2). Second reading: CORRECT (recorded in \src{CHECK\_REPORT.md}); in particular items 1--4 of \src{N21} Thm.~N21.4 use only the $A$-half of (TS). It refutes the residual case (E5b) left open in \src{N21}.}

\begin{theorem}\label{thm:dichotomy}
Let a strictly guiding additive ranking on a decomposable domain with (LF) cover $T_1$ (integral penalties, or bounded charges given interior minimisers). Then either
\begin{enumerate}[leftmargin=*,label=(\Roman*)]
\item (TS) holds, and then (PC), (NS) hold, the advice range of a.e.\ mixed pair contains both acceptances, all vertex models are precise, interior limits at edge priors lie in the edge face, and every vertex carries the cut structure of Theorem~\ref{thm:cuts}; or
\item (TS) fails, and then some facet $F_k$ is pinned: there are an open convex $G$ with $p^*\in\cl G$, on which $v_k\in K_q$ and $\Psi(q,\cdot)\equiv0$ on an interval around a crossing; every segment from $G$ towards a vertex neighbourhood $V_l$ contains an interior $Q$ with $K_Q=\{v_k\}$; and accepting $h$ lies in the advice range for a.e.\ $h$ with $v_k\cdot h>0$.
\end{enumerate}
\end{theorem}
\status{proved; refereed once (\src{N23} \S8, assembling \src{N21} Thm.~N21.8 and Prop.~N21.9, both refereed). Branch (II) as stated also needs Corollary~\ref{cor:R232}, which excludes a failure of the $B$-half alone; that corollary has now had a second reading (CORRECT). The further refinement $A_{\tau_c}\cap B_1=\varnothing$ at the crossing (\src{N23} Prop.~R23.3) is not used in the statement and remains a one-reader result. Branch (II) never occurs (Theorem~\ref{thm:TSalways}).}

\begin{corollary}\label{cor:R234}
In every covering of $T_1$ as in Theorem~\ref{thm:dichotomy}, the advice range of a.e.\ mixed pair contains at least one acceptance. Hence no such covering exists on a decomposable domain with (LF) on which, for a non-null set of mixed pairs, the advice range contains neither acceptance. If (TS) holds and, for a non-null set of mixed pairs, the advice range lacks one acceptance, there is no covering either.
\end{corollary}
\status{the first two sentences: proved in a referee report (\src{N23} \S7.5, Cor.~R23.4); second reading (\src{CHECK\_REPORT.md}): CORRECT, from Corollary~\ref{cor:availability} if (TS) holds and from Corollary~\ref{cor:R232} with \src{N21} Prop.~N21.9(ii)--(iii) otherwise. The last sentence: \src{N21} Thm.~N21.8(b), refereed once (\src{N23}); by Corollary~\ref{cor:TSalways}(b) it no longer needs the hypothesis (TS).}

\begin{proposition}\label{prop:noWLOG}
Restricting the domain gives no reduction to the non-adjacent case: removing ``$-$'' from the advice range of a non-null set of pairs (that is, removing from the class every model that accepts $-g$ there) removes the vertex selections from the class of evaluable models, so \src{R11} Lemma RS does not apply, and removing it while keeping the penalties changes the Bayes family and can destroy coherence (an explicit two-world example exists).
\end{proposition}
\status{proved; refereed once (\src{N23} \S5, scope repair: what is shown is that the restriction route fails, not that no reduction exists). Source: \src{N21} Prop.~N21.5.}

\subsection{Competing acceptances: no pinning, and the cuts reach the vertices}\label{sec:general-N28}

In this subsection a \emph{cover} means a strictly guiding additive ranking (integral penalties, or bounded charges vanishing on null sets with interior minimisers) on a decomposable domain with (LF) at $D_T$, satisfying (LA), (S), (C) and covering $T$; $K_q$ are its interior Bayes credal sets and $K_{p^*}=T$. For a facet $F_k$ write $A^k_\tau,B^k_\tau$ for the sets $A_\tau,B_\tau$ of \S\ref{sec:main}. Define the \emph{wall} of a probability $m_0$ and the set of priors \emph{avoiding} the vertex $e_b$ by
\[
\mathcal W_{m_0}\coloneqq\{q\in\Di:K_q=\{m_0\}\},\qquad \mathcal A_b\coloneqq\{q\in\Di:e_b\notin K_q\}.
\]

\begin{lemma}\label{lem:walls}
In every cover:
\begin{enumerate}[leftmargin=*,label=(\alph*)]
\item each wall $\mathcal W_{m_0}$ is convex and relatively closed in $\Di$;
\item each $\mathcal A_b$ is open and convex;
\item for a.e.\ pair $g$ and every $m_0\in\Delta$ with $m_0\cdot g\ne0$, $\cl F(g)\cap\Di$ contains no $q$ with $m_0\in K_q$;
\item for a vertex $v_m$ of $T$, $i\ne m$ and $\tau\in(0,1)$: $\mathcal W_{v_m}\subseteq A^i_\tau$, which is convex and does not contain $p^*$.
\end{enumerate}
\end{lemma}
\begin{proof}[Proof sketch]
(a) $\mathcal W_{m_0}=\bigcap_{h:\,m_0\cdot h>0}\{P(\cdot,h)>0\}$, an intersection of convex sets (Lemma~\ref{lem:quasiconcave}); for $m\in K_q\setminus\{m_0\}$ the gamble $h=(m_0-m)-((m_0-m)\cdot m)\one$ separates. (b) $\mathcal A_b=\{P(\cdot,\one-e_b)>0\}$. (c) \src{R4} Thm.~B: $K_q\subseteq g^\perp$ on $\cl F(g)\cap\Di$ for a.e.\ $g$. (d) On $\mathcal W_{v_m}$, $P(q,h^i(\tau))=v_m\cdot h^i(\tau)>0$, while $P(p^*,h^i(\tau))=0$.
\end{proof}
\status{proved (\src{N28} Lemmas 2.1--2.4); confirmed by one referee (\src{N28\_referee} \S2).}

\begin{theorem}\label{thm:TSalways}
Let $T$ be an all-facing triangle in general position (in particular $T_1$, for every $\varepsilon\in(0,1)$). In every cover of $T$, (TS) holds. Hence branch (II) of Theorem~\ref{thm:dichotomy} never occurs: there is no single-vertex pinning near $p^*$, by a ball or by a thin set.
\end{theorem}
\begin{proof}[Proof sketch]
Suppose (TS) fails. By Corollary~\ref{cor:R232} the $A$-half fails on some facet $F_m$, and branch (II) gives a \emph{wall of $F_m$ toward $V_l$} for some $l\ne m$: an open set $O$ with $p^*\in\cl O$ and $v_m\in K_q$ on $O$, such that every segment from $O$ to $V_l\cap\Di$ meets $\mathcal W_{v_m}$. (i) Such a wall forbids (TS) at every $\tau$ on $F_l$: since $h^l(\tau)_l<0$, the $c$-cluster sign puts $e_l$ on the closed negative side of the chord, so some point $v\in V_l\cap\Di$ lies strictly beyond it, while the wall points lie in $A^l_\tau$ (Lemma~\ref{lem:walls}(d)) inside the positive half-plane (Lemma~\ref{lem:quasiconcave}); for $o$ close to $p^*$ the wall point on $[o,v]$ would be closer to $p^*$ than $\operatorname{dist}(p^*,\mathcal W_{v_m})>0$. (ii) By Lemma~\ref{lem:R231} and \src{R10b} Thm.~PL, $B$-failures on $F_l$ cannot straddle a crossing next to the pinned set, so $F_l$ has walls toward both other vertices. (iii) Two relatively closed walls that meet every segment from sets accumulating at $p^*$ toward all three vertex neighbourhoods have $p^*$ in the convex hull of their union. (iv) Both walls lie in the convex set $A^n_\tau$ ($n$ the third index), which misses $p^*$.
\end{proof}
\status{proved (\src{N28} Prop.~3.1, Lemmas 3.2--3.4, Thm.~3.5); confirmed by one referee (\src{N28\_referee} \S3, every step re-derived, including that the walls of branch (II) meet \emph{every} segment from the pinned set to $V_l\cap\Di$, as the surround step needs). It rests on once-refereed inputs: \src{N21} Prop.~N21.9 and Thm.~N21.4 (refereed in \src{N23}); \src{R10b} Thm.~PL, Lemma UVN and Lemma M (refereed in \src{R12}); Lemma~\ref{lem:R231} and Corollary~\ref{cor:R232} (\src{N23}, with a second reading).}

\begin{corollary}\label{cor:TSalways}
\begin{enumerate}[leftmargin=*,label=(\alph*)]
\item Every cover of an all-facing $T$ in general position is in branch (I) of Theorem~\ref{thm:dichotomy}: (TS), (PC), (NS) hold, and Theorem~\ref{thm:cuts} and Corollary~\ref{cor:availability} hold at every world with a failing box.
\item $T_1$ has no additive cover on any decomposable domain with (LF) on which, for a non-null set of mixed pairs, the advice range lacks one of the two acceptances: for instance, advice ranges contained in $\{+,0\}$ or in $\{0,-\}$ there, or a non-null set of mixed pairs on which every evaluable model gives no advice.
\item In every cover of $T_1$, $K_q\to\{e_b\}$ as $q\to e_b$ for every $b$, and every limit of $K_q$ at an edge prior with support $S$ lies in $\Delta_S$.
\end{enumerate}
\end{corollary}
\status{proved (\src{N28} Cor.~3.6); confirmed by one referee. (b), (c) are \src{N21} Cor.~N21.7 and Thm.~N21.8(b), whose only extra hypothesis was (TS). So in every cover the advice range of almost every mixed pair contains both acceptances, and competition between them is not merely permitted but required.}

For the rest of this subsection fix a world $b$ of $T_1$, the index $a$ with a failing box $\Lambda_{ab}>1$, and the third index $c$. The two $b$-crossings are $w'\in\relint F_a$ and $w\in\relint F_c$ (for $(a,b,c)=(1,2,3)$, $w'\propto(-1,0,\kappa)$ and $w\propto(\kappa,0,-1)$). Under (TS) both crossing chords are the line $\lambda_b\coloneqq\lin(p^*,e_b)$. The tangent lines from $e_b$ to $T_1$ are $\ell'\coloneqq\lin(e_b,v_a)$ and $\ell\coloneqq\lin(e_b,v_c)$. One has $\{m\in\Delta:m\cdot w'\le0,\ m\cdot w\le0\}=\{e_b\}$ (\src{N28} Fact 4.0, confirmed). With the data $x_+$ the failing pairs $(a,b)$ are $(1,2),(2,3),(3,1)$; with $x_-$ they are $(2,1),(3,2),(1,3)$ (\src{N28\_referee}, NUMERICS from Proposition~\ref{prop:boxbridge}(c)).

\begin{theorem}\label{thm:apex}
In Theorem~\ref{thm:cuts}, the apex limit is $a^*=e_b$. Competition between accepting $g$ and accepting $-g$ at a positive distance from $\partial\Delta$ is impossible. The same holds for every all-facing triangle in general position at every world with a failing box.
\end{theorem}
\begin{proof}[Proof sketch]
If $a^*\in\Di$, let $w^{**}$ be the other $b$-crossing. Along $[a^*,z^*]$, $K_q\subseteq\ell_{w^*}$ gives $P(q,w^{**})\ge0$. If $z^*\ne e_b$, convexity of $\{P(\cdot,w')\ge0\}$ and $P(\cdot,w')\le0$ off the positive side give a triangle in $B_1$ on which $P(\cdot,w')=0$, contradicting thin zero sets. If $z^*=e_b$, then $P(a^*,w^{**})=0$ forces $e_b\in K_{a^*}$; separating $a^*$ from the open convex $\mathcal A_b\ni p^*$ puts a ball around a point of $(a^*,e_b)$ outside $\mathcal A_b$, and d-face points of $g_m$ in that ball contradict Lemma~\ref{lem:walls}(c).
\end{proof}
\status{proved (\src{N28} Thm.~4.1); confirmed by one referee (\src{N28\_referee} \S4.1). This settles the case $a^*\in\Di$ that \src{N23} \S7.6 singled out.}

\begin{lemma}\label{lem:R1N28}
In every cover of an all-facing $T$ in general position, for every facet $F_k$ and $\tau\in(0,1)$, with $H^\pm_\tau$ the open sides of $\lambda_\tau$ ($H^+_\tau$ containing $A^k_\tau$): $P(\cdot,h(\tau))>0$ only on $H^+_\tau$, $P(\cdot,h(\tau))<0$ on all of $\Di\cap H^-_\tau$, and $P(\cdot,h(\tau))\le0$ on $\lambda_\tau\cap\Di$. So the chord line separates $\Di$ globally.
\end{lemma}
\begin{proof}
By Theorem~\ref{thm:TSalways}, $Z_\tau\cap B_1=\lambda_\tau\cap B_1$. If $q_1\in\Di\cap H^-_\tau$ had $P(q_1,h(\tau))\ge0$, the convex set $\{P(\cdot,h(\tau))\ge0\}$ (Lemma~\ref{lem:quasiconcave}) would contain the triangle spanned by $q_1$ and $\lambda_\tau\cap B(p^*,r_1/2)$, on which also $P\le0$; so $P=0$ on a triangle meeting $B_1\setminus\lambda_\tau$, a contradiction. The positive statement is Lemma~\ref{lem:quasiconcave}.
\end{proof}
\status{proved in a referee report (\src{N28\_referee} \S4.2, Lemma R1); one reader only. It supersedes the CHECKED remark after \src{N28} Prop.~4.3.}

\begin{proposition}\label{prop:vertexfree}
In every cover of $T_1$:
\begin{enumerate}[leftmargin=*,label=(\roman*)]
\item no interior credal set $K_q$, $q\in\Di$, contains a vertex of $\Delta$;
\item for every facet, $\tau\in(0,1)$ and $q\in\Di\setminus H^+_\tau$, $K_q$ meets $\{m:m\cdot h(\tau)\le0\}$, and it meets $\{m:m\cdot h(\tau)<0\}$ if $q\in H^-_\tau$;
\item for $q\in\lambda_b\cap\Di$, $K_q$ contains a point with $m\cdot w'\le0$ and a different point with $m\cdot w\le0$, so $K_q$ is not a singleton; on the whole open $e_a$-side of $\lambda_b$, $K_q$ meets $\{m\cdot w'<0\}$, and on the whole open $e_c$-side it meets $\{m\cdot w<0\}$;
\item no precise prior lies on a crossing chord. In particular a precise identity cap $K_q=\{q\}$ on a full neighbourhood of $e_b$ (\src{R4} \S9), and every diagonal-projective precise region meeting $\lambda_b$, are incompatible with every cover.
\end{enumerate}
\end{proposition}
\begin{proof}[Proof sketch]
(i) If $e_b\in K_y$, separate $y$ from $\mathcal A_b$, which contains points near $e_a$ and $e_c$ (Corollary~\ref{cor:TSalways}(c)); the separating functional is negative near $e_b$, and by Theorem~\ref{thm:apex} some apex $a_m\in\cl F(g_m)$ lies there, against Lemma~\ref{lem:walls}(c). (ii) is Lemma~\ref{lem:R1N28}. (iii) is (ii) for $w'$ and $w$ at the crossing parameter; a common point would be $e_b$, excluded by (i). (iv) A singleton on $\lambda_b$ would be $\{e_b\}$.
\end{proof}
\status{proved (\src{N28} Prop.~4.3); confirmed by one referee (\src{N28\_referee} \S4.3). The strict forms in (ii) and (iii) use Lemma~\ref{lem:R1N28} (one reader). The arithmetic for the identity cap was checked numerically in \src{N28} \texttt{checks.py} (3).}

\subsection{Reduction to a problem about credal families}\label{sec:general-CF}

\begin{definition}\label{def:CF}
Let $T\subseteq\Di$ be a credal set. The \emph{concurrent-chart family problem} CF$(T)$ asks for a Hausdorff-continuous map $q\mapsto K_q$ from $\Di$ to credal sets, with $K_{p^*}=T$ for some $p^*\in\Di$, such that for a.e.\ mixed $g\in\R^3$ there are $c,c'\in\R^3$ with $\operatorname{sign}c_i=\operatorname{sign}g_i=-\operatorname{sign}c'_i$ ($i=1,2,3$) and, for \emph{every} $q\in\Di$,
\[
P_{K_q}(g)>0\iff c\cdot q>\max(0,c'\cdot q),\qquad P_{K_q}(-g)>0\iff c'\cdot q>\max(0,c\cdot q).
\]
\end{definition}
In words: for almost every pair, the partition of the prior triangle into ``accept $g$'', ``accept $-g$'' and ``neither'' is the cell decomposition of $\max\{0,c\cdot q,c'\cdot q\}$, three lines through one point, with the vertex signs of \src{R1}.

\begin{proposition}\label{prop:CF}
If an additive ranking with integral penalties covers $T_1$ on some decomposable domain with (LF), then CF$(T_1)$ has a solution. Conversely, every solution of CF$(T_1)$ yields such a cover, on the all-model domain, provided Theorem~\ref{thm:ccibu} holds. The prior $p^*$ may be chosen arbitrarily in $\Di$.
\end{proposition}
\begin{proof}[Proof sketch]
($\Rightarrow$) Propositions~\ref{prop:necessary} and~\ref{prop:reduction} give the family. By Corollary~\ref{cor:TSalways}(b) the advice range of a.e.\ mixed pair contains both acceptances, so advice ranges are $\{+,0,-\}$ or $\{+,-\}$; a $\{+,-\}$ chart is the cell decomposition of $(\lambda d+e_0\one,(\lambda-1)d+e_0\one)$ for $0<\lambda<1$ and small $e_0>0$. ($\Leftarrow$) Take advice ranges $\{+,0,-\}$, penalties from \src{R1} Prop.~2.1 (converse), rescaled per pair for integrability; a measurable choice of $(c,c')$ exists because the admissible set is Borel (test $C\subseteq R(c,c')\Leftrightarrow\min_CP(\cdot,g)>0$ over closed rational polygons $C$) and the Jankov--von Neumann theorem applies. Tie sets lie in $\pm\partial D_{K_q}$, so Bayes models are unique and coherent at every interior $q$, and Theorem~\ref{thm:ccibu} gives (LA), (S), (C). The last sentence is Proposition~\ref{prop:diag}(a).
\end{proof}
\status{CHECKED in \src{N28} (Prop.~5.1); confirmed with repair by one referee (\src{N28\_referee} \S5.1). Repairs incorporated: the measurable-selection test above replaces the note's ``test on a countable dense set of $q$'', which presumed convexity of $\{P(\cdot,g)>0\}$ for an arbitrary candidate; ($\Rightarrow$) needs both acceptances in the advice range a.e., so it is specific to $T_1$ (and to targets whose covers have both acceptances in the advice range a.e.); ($\Leftarrow$) has the status of Theorem~\ref{thm:ccibu} (author-checked, no formal referee report).}

\begin{proposition}\label{prop:diag}
Let $D,D'$ be positive diagonal matrices and $\nu$ the normalisation $x\mapsto x/(\one\cdot x)$.
\begin{enumerate}[leftmargin=*,label=(\alph*)]
\item If $K$ solves CF$(T)$ at $p^*$, then $\tilde q\mapsto K_{\nu(D'^{-1}\tilde q)}$ solves CF$(T)$ at $\nu(D'p^*)$, with chart data $D'^{-1}c,D'^{-1}c'$.
\item If $K$ solves CF$(T)$ at $p^*$, then $q\mapsto\nu(D^{-1}K_q)$ solves CF$(\nu(D^{-1}T))$ at $p^*$, with data $(\tilde c,\tilde c')(Dg)\coloneqq(c,c')(g)$.
\item The map $X\mapsto DX$ on models, with $A(g)\mapsto A(D^{-1}g)$ on advice ranges, preserves decomposable domains, splicing, convergence in measure, null sets, guidance ranks, (IB), coherence ($D\cdot D_K=D_{\nu(D^{-1}K)}$) and (LF). Hence a ranking covers $D_T$ iff its transport covers $D_{\nu(D^{-1}T)}$; for additive rankings, multiply the penalties by the Radon--Nikodym factor of the transported Gaussian measure.
\end{enumerate}
\end{proposition}
\status{(a), (b): proved (\src{N28} Prop.~5.2); confirmed by one referee. (a) is Proposition~\ref{prop:reduction}(d). (c): proved in a referee report (\src{N28\_referee} \S5.2, Prop.~R3), one reader only. By (c), for every $T$ and every class of rankings closed under this transport, coverability of $T$ depends only on the orbit of $T$ under $T\mapsto\nu(D^{-1}T)$; diagonal gamble rescaling is a symmetry of the whole axiom system, and no appeal to Proposition~\ref{prop:CF} or Theorem~\ref{thm:ccibu} is needed.}

By \S\ref{sec:general-N28}, every solution of CF$(T_1)$ coming from a cover satisfies (TS), (PC), (NS), $K_q\to\{e_b\}$ at the vertices with edge confinement, cut apexes at the vertices, vertex-free interiors, global one-sidedness of chords, non-precise credal sets on the crossing chords, and the wall and d-face constraints of Lemma~\ref{lem:walls}. By Proposition~\ref{prop:CF}($\Leftarrow$) every solution comes from a cover (at the status of Theorem~\ref{thm:ccibu}), so the residual called (E4$''$) in \src{N28} is CF$(T_1)$ itself, together with these necessary conditions; the list is not exhaustive (\src{N28\_referee} \S6).

\subsection{What remains open on general decomposable domains}\label{sec:general-open}
\begin{itemize}[leftmargin=*]
\item \textbf{(E5a)} \emph{Single-vertex pinning near $p^*$.} Closed: it never occurs (Theorem~\ref{thm:TSalways}).
\item \textbf{(E4$'$)} \emph{The near-all-orientation cut regime.} Reduced to CF$(T_1)$ (Proposition~\ref{prop:CF}). Competition at a finite distance from $\partial\Delta$ is excluded (Theorem~\ref{thm:apex}), and a precise cap on a full neighbourhood of a vertex is excluded (Proposition~\ref{prop:vertexfree}(iv)).
\item \emph{Vertex-scale consistency is not excluded.} \src{N28} inferred that, with the identity cap of \src{R4} \S9 excluded, the claim that (E4$'$) is consistent at the scale of the vertices had lost its support. That inference is wrong (\src{N28\_referee} \S4.3). The cuts of \src{R4} \S9 run from the apex to the failing edge and never reach $\lambda_b$. The identity $K_q=\{q\}$ restricted to a \emph{sector} $S_b$ at $e_b$ adjacent to the failing edge passes every constraint derived in \S\ref{sec:general-N28}, Lemma~\ref{lem:R1N28} included, and carries the cuts (angular extent $\varphi_{\max}\approx0.367,0.196,0.127,0.026$ at $\varepsilon=0.05,0.2,0.3,0.6$, in an angular parameter measured from the failing edge with $\varphi=\tfrac12$ on $\lambda_b$; status NUMERICS, \src{N28\_referee} counterexample C1). A first-order blow-up further suggests that precise priors near $e_b$ occur only strictly on the failing side of $\lambda_b$ (NUMERICS N1; conjectured by the referee).
\item \emph{Where a solution of CF$(T_1)$ would have to live.} Precise sector caps adjacent to the failing edge at each vertex carry the cuts; a non-precise layer along $\lambda_b$ (and, numerically, on the whole non-failing half-neighbourhood of $e_b$) straddles the tangent wedge at $e_b$ without touching $e_b$, poking strictly beyond $\ell$ on the open $e_c$-side and beyond both $\ell$ and $\ell'$ on the chord; the sector caps and the non-precise layer must be glued to the core along the chords under every concurrent chart. The gluing across the sector boundary is the natural next test.
\end{itemize}
\begin{conjecture}\label{conj:N21}
CF$(T_1)$ has no solution.
\end{conjecture}
\status{open (\src{N28} Conj.~N28, restated in \src{N28\_referee} \S6; it replaces \src{N21} Conj.~N21). By Proposition~\ref{prop:CF}($\Rightarrow$) it implies that no additive ranking with integral penalties covers $T_1$ on any (LF) decomposable domain; the converse implication has the status of Theorem~\ref{thm:ccibu}. No construction is known, and the sector-cap example shows that the constraints derived so far do not exclude one.}

\section{Coherent almost-preferences and rankings that add up}\label{sec:cp}

Standing hypotheses \textbf{(H)}: $n=3$; a decomposable domain $\M$; the ranking satisfies strict (G), (IB), (LA), (S), (C), and covers $T_1$ with (LF). Let $N_i$ be the set of problems whose advice range has at least two elements, i.e.\ on which evaluable models can give different advice (on $\M_U$: the pairs $\{f,-f\}$ with $f\in U$ or $-f\in U$, not only those whose representative in $H$ lies in $U$). Under (CP-a) each state splits as $P=P^c+P^s$ (Yosida--Hewitt) with $P^c=\rho_P\mu$, $\rho_P\in L^1_+$, and $P^s$ purely finitely additive; any finitely many purely finitely additive parts have a common carrier of arbitrarily small $\mu$-measure (\src{N14} Lemma Y; \src{N11} Fact 1.2).

\subsection{Normal form and full support}

\begin{proposition}\label{prop:normalform}
Under (CP-a), $\preceq_i$ is a preorder, (IB) and weak (G) hold automatically, and every $P\in\M_i$ weakly respects the ranking: $X\preceq_iY\Rightarrow P(\Lambda_iX)\le P(\Lambda_iY)$ (property (W)). Strict (G) at $i$ holds iff $\up_i(1_A)\coloneqq\max_{P\in\M_i}P(1_A)>0$ for every non-null $A\subseteq N_i$.
\end{proposition}
\status{proved; refereed once each (\src{N13} for \src{N11} Prop.~1.1; \src{N24} for \src{N22} Lemma 1.0).}

\begin{theorem}\label{thm:fullsupport}
Under (CP-a) and strict (G) at world $i$, the set $\mathcal P^+_i\coloneqq\{P\in\M_i:\rho_P>0\text{ a.e.\ on }N_i\}$ is nonempty, and $(1-t)P+tQ\in\mathcal P^+_i$ whenever $P\in\mathcal P^+_i$, $Q\in\M_i$, $t\in[0,1)$.
\end{theorem}
\begin{proof}[Proof sketch]
Key lemma: if a weak-* compact convex set $\M$ of states has $P^c(A)=0$ for every $P\in\M$ and some non-null $A$, there is a non-null $N\subseteq A$ with $P(N)=0$ for all $P\in\M$. For one step, Kneser's minimax theorem applied to $(P,f)\mapsto P(f)$ on $\M\times\{0\le f\le1_{A'},\int f=(1-\eta^2)\mu(A')\}$ gives a large $B\subseteq A'$ of uniformly small $\M$-mass, because each single member has small carriers. Iteration and countable additivity of $\mu$ then give $N$. With the countable chain condition, the essential union of the supports of the $\rho_P$ must cover $N_i$, by Proposition~\ref{prop:normalform}; a countable convex combination of members (normalised partial sums, norm-closed) gives an element of $\mathcal P^+_i$.
\end{proof}
\status{proved; refereed once (\src{N24} \S1: CORRECT WITH REPAIR, cosmetic). Source: \src{N22} Lemma 1.1, Thm.~1.2. Of independent interest: a weak-* compact convex set of states on $L^\infty(\mu)$ that charges every non-null set contains a state whose normal part is equivalent to $\mu$.}

\subsection{The integral reduction and the residual (U)}

For $P=(P_1,P_2,P_3)\in\prod_i\mathcal P^+_i$ define the integral field $\psi^P_i(X)\coloneqq\int\lambda_i(g,a_X(g))\,\rho_{P_i}(g)\,d\mu(g)$, let $Z_p(P)$ be the set of $g$ at which $p\cdot\psi^P(g,\cdot)$ has two or more minimisers on $A(g)$, and put $E(P)\coloneqq\{p\in\Di:\mu(Z_p(P))>0\}$.

\begin{theorem}\label{thm:integral}
Assume (H) and (CP-a). For every $P\in\prod_i\mathcal P^+_i$, $\psi^P$ is a strictly guiding integral field on the same domain $\M$ (so (LF) at $T_1$ is kept), and:
\begin{enumerate}[leftmargin=*,label=($\alph*$)]
\item $T_1$ is $\psi^P$-Bayes at some $p_T\in\Di$;
\item for every $p\in\Delta$ some $\psi^P$-$p$-Bayes model is coherent;
\item $E(P)$ is Lebesgue-null, and at every $p\in\Di\setminus E(P)$ the $\psi^P$-$p$-Bayes model is a.e.\ unique and coherent;
\item at each vertex the Bayes model is the world-best evaluable model, which is coherent.
\end{enumerate}
\end{theorem}
\begin{proof}[Proof sketch]
($a$) $T_1$ is competitive on $\Omega$ (\src{R11} Lemma IC, refereed in \src{R13}). Splicing off a common small carrier of the three singular parts shows that no model is better than $T_1$ for all three integral parts; Lyapunov and separation give a weight, which the Lemma IC construction pushes into $\Di$. ($b$) Approximate the infimum of the full charge objective by models equal to a pointwise integral minimiser off the carriers; by Theorem~\ref{thm:kkm} select weak dominators; by (W) for the full members they are near-minimisers of the integral objective; compactness in measure of coherent models gives a coherent minimiser. ($c$) Tonelli on $H\times\Delta$; off $E(P)$ the minimiser is pointwise unique, so it is the coherent one. ($d$) Strict guidance and the vertex lemma.
\end{proof}
\status{proved; refereed once (\src{N24} \S2). Source: \src{N22} Thm.~2.1, building on \src{N11} Cor.~3.2 and \src{N13} Prop.~B (refereed in \src{N13}, \src{N24}). No detection, attainment or non-atomicity hypothesis is used; atoms enter only through finitely many small carriers.}

\begin{corollary}\label{cor:residual}
Assume (H) and (CP-a) on a decomposable domain non-adjacent near $T_1$. Then $E(P)\ne\varnothing$ for every $P\in\prod_i\mathcal P^+_i$. Equivalently, $T_1$ cannot be covered once the following holds:
\begin{quote}
\textbf{(U)} some $P\in\prod_i\mathcal P^+_i$ has no non-null pointwise Bayes tie at any interior weight.
\end{quote}
\end{corollary}
\begin{proof}
If $E(P)=\varnothing$, then (on $\M_U$; on other non-adjacent domains after the reduction used in Corollary~\ref{cor:coverage}) the vector $c_i(g)\coloneqq\rho_{P_i}(\bar g)\bigl(\lambda_i(\bar g,0)-\lambda_i(\bar g,a_g)\bigr)$, $g\in U$, where $\bar g$ is the representative in $H$ of the pair of $g$ and $a_g\in\{+,-\}$ is the advice ``accept $g$'', satisfies (BU) at \emph{every} interior weight and (TP) at $p_T$, which Corollary~\ref{cor:coverage} excludes. Uniqueness off a null set of weights would not suffice (Remark~\ref{rem:N20}).
\end{proof}
\status{proved; refereed once (\src{N24} \S2.2, which checked that the every-weight form of (BU) is what is supplied). Source: \src{N22} Cor.~2.3.}

\begin{theorem}\label{thm:complete}
Assume (H), and suppose every world is complete: $\M_i=\{P_i\}$ for states $P_i$ with arbitrary purely finitely additive and atomic parts. Assume \textbf{(UN)}: for each $i$ there is a non-null $U_+\subseteq N_i$ on which every coherent model of $\M$ gives the advice in the advice range that is best at world $i$ (true on $\M_U$ with $U_+=U\cap\{g\gg0\}$, and on the all-model domain). Write $\varphi_i\coloneqq P_i\circ\Lambda_i$. Then:
\begin{enumerate}[leftmargin=*,label=(\alph*)]
\item for every $p\in\Delta$, $\inf_\M p\cdot\varphi$ is attained by a selected model;
\item for every $p\in\Di$ the minimiser is a.e.\ unique, coherent and selected;
\item $P\in\prod_i\mathcal P^+_i$ and $E(P)=\varnothing$.
\end{enumerate}
\end{theorem}
\begin{proof}[Proof sketch]
(a) Take a minimising sequence, selected weak dominators $C_{q_m}$ (Theorem~\ref{thm:kkm}), and subsequences with $q_m\to q^*$, $C_{q_m}\to C_\infty$ coherent and $\varphi(C_{q_m})\to z$. Worsening $C_m$ off a small carrier shows that the attainable world-$i$ values contain intervals of length bounded below, which uses (UN); so there are $X$ with $\varphi_i(X)$ just above $z_i$. Completeness turns $\varphi_i(C_{q_m})<\varphi_i(X)$ into $C_{q_m}\preceq_iX$, and the closed sections of (S) give $C_{q^*}\preceq_iX$. Hence $\varphi(C_{q^*})\le z$ and $C_{q^*}$ attains. (b) A selected weak dominator of a minimiser has equal profile; completeness makes this a tie at every world, and (C) gives coherence. Uniqueness follows from complementary splices and a bounded non-null window. (c) Ties of $\varphi$ are ties of $L$, so detection is automatic.
\end{proof}
\status{proved; refereed once (\src{N24} \S3: CORRECT; window-radius repair incorporated). Source: \src{N22} Thm.~3.1.}

\begin{corollary}\label{cor:complete}
On the $U$-domain, and on every decomposable domain non-adjacent near $T_1$ that satisfies (UN), no ranking satisfying strict (G), (IB), (LA), (S), (C) whose worlds are complete coherent almost-preferences over losses covers $T_1$. Equivalently, no additive ranking with bounded charge penalties vanishing on null sets covers $T_1$ there; the existence proviso of Corollary~\ref{cor:charges} is discharged by (S).
\end{corollary}
\status{proved; refereed once (\src{N24} \S3). Source: \src{N22} Cor.~3.2, using Theorem~\ref{thm:main}.}

\begin{corollary}\label{cor:DCA}
Assume (H), (CP-a) and countably additive detection: whenever $X\prec_iY$, some countably additive $Q\in\M_i$ has $Q(\Lambda_iY-\Lambda_iX)>0$ (for example, every member of $\M_i$ is countably additive, or $\M_i$ is the weak-* closure of its countably additive members). Then the countable aggregate of a norm-dense sequence of countably additive members is an integral field with a.e.-unique coherent Bayes models at every interior weight and $T_1$ among them. Consequently $T_1$ is not covered on decomposable domains non-adjacent near $T_1$.
\end{corollary}
\status{proved; refereed once (\src{N13} for \src{N11} Thm.~2.1, which uses \src{R11} Thm.~W with \src{R13} Prop.~S; \src{N24} \S3.3 for the application of Corollary~\ref{cor:residual}).}

\begin{proposition}\label{prop:Ucont}
On $\M_U$, fix $P\in\prod_i\mathcal P^+_i$ and let $K_q$ be the credal sets of the unique coherent $\psi^P$-Bayes models for $q\in\Di\setminus E(P)$. Then $q\mapsto K_q$ is Hausdorff-continuous on $\Di\setminus E(P)$, and $p_0\notin E(P)$ iff $K$ has a limit at $p_0$ along $\Di\setminus E(P)$. So (U) holds iff some member choice has a coherent Bayes family that extends continuously to all of $\Di$.
\end{proposition}
\status{proved in a referee report (\src{N24} \S6.1, Prop.~R1), one reader only; the non-adjacent version is stated there at the level of an expectation and was not re-checked.}

\begin{remark}\label{rem:EPchords}
On $\M_U$ the field $\psi^P$ is strictly signed and integrable, and its Bayes models are coherent at a.e.\ interior weight (Theorem~\ref{thm:integral}($c$)). So Theorem~\ref{thm:Estructure} below applies to it: $E(P)$ is a countable union of open chords of $\Delta$, each the common chart line of a non-null set of gambles and each cutting off one vertex, with no isolated exceptional points. Hence on $\M_U$, (U) fails exactly when every member choice $P$ has such a chord. (This application is ours; the sources state Theorem~\ref{thm:Estructure} for the M0 setting, but its proof uses only the properties just listed.)
\end{remark}

\begin{proposition}\label{prop:CD}
If every $\M_i$ has a countable subfamily detecting all strict comparisons (for example, $\M_i$ weak-* separable: finitely generated, contained in the countably additive states, or the closed convex hull of a norm-separable set), there is $\hat P\in\prod_i\mathcal P^+_i$ with $E(\hat P)$ contained in the set of interior weights at which $\inf_\M p\cdot(\hat P\circ\Lambda)$ is not attained. Hence countable detection together with attainment of one charge field at every interior weight excludes a covering on non-adjacent domains.
\end{proposition}
\status{proved in a referee report (\src{N24} \S6.2, Prop.~R2), one reader only. Attainment at every interior weight is open even for finitely generated incomplete $\M_i$ (\src{N22} Remarks 3.3).}

At a tie weight a hypothetical cover has a specific anatomy: if $P$ consists of countably additive members, the face of Bayes profiles at $p_0\in E(P)$ has dimension $\ge1$, is realised by selected models, and carries a strict comparison that $P$ does not see (\src{N22} Prop.~4.1, refereed with repair in \src{N24} \S4).

\begin{conjecture}\label{conj:U}
(U) holds for every (CP-a) system satisfying (H). On non-adjacent domains this is equivalent to: no (CP-a) ranking covers $T_1$.
\end{conjecture}
\status{open (\src{N22} Conj.~N22.U). It subsumes the earlier detection, atom and existence questions (\src{N11} Conj.~A, B; \src{N14} existence proviso).}

\subsection{Strict preference}
Walley's strict-preference model \cite[\S\S3.7.7--3.7.9]{W1991}, read in a preorder framework as $\prec\cup{=}$, gives condition (CP-s): $X\preceq_iY$ iff $X=Y$, or $\lp_i(\Lambda_iY-\Lambda_iX)>0$, or $\Lambda_iY-\Lambda_iX\ge0$ a.e.\ and $\ne0$. Walley 2000 \cite[\S6]{W2000} gives (CP-b): $X\preceq_iY$ iff $\Lambda_iY-\Lambda_iX\in\mathcal D_i\cup\{0\}$ for a coherent set $\mathcal D_i$ of desirable gambles (D1--D4); (CP-s) is the case $\mathcal D_i=\{\lp_i>0\}\cup\{f\ge0,f\ne0\}$.

\begin{proposition}\label{prop:A}
Let two worlds $\omega\ne j$ satisfy (CP-s) with weak-* compact sets of countably additive states with a.e.-positive densities, on $\M_U$ (or any decomposable domain whose advice ranges contain, for every selected model, advice worse at every world than the model's own on a non-null set of unanimous problems, and advice better at $\omega$ than the model's own wherever the model errs at $\omega$), and suppose the world-ideal models $W_\omega\ne W_j$ on a non-null set. Then no selection $u:\Delta\to\C$ satisfies (S), even without (LA), (C) or a target.
\end{proposition}
\begin{proof}[Proof sketch]
Under strict positivity, $Y\prec Z$ iff $\lp(\Lambda Z-\Lambda Y)>0$ at these worlds, so upper sets $\{Y:C\prec Y\}$ are open (Walley's openness axiom). For a selected $C^*=C_{p_0}$ build an incoherent $X$ with $\lp_\omega(\Lambda_\omega X-\Lambda_\omega C^*)=0$; then $p_0\notin\{p:C_p\preceq_\omega X\}$, which is closed, while every $p$ with $C_p\prec_\omega C^*$ lies in it. So near $p_0$ on the edge $[e_\omega,e_j]$ no selected model is strictly better than $C^*$ at $\omega$ or $j$; competitiveness forces $C_p=C^*$ there. The selection is locally constant on the connected edge, so $W_\omega=C_{e_\omega}=C_{e_j}=W_j$.
\end{proof}
\status{proved in a referee report (\src{N13} \S3, Prop.~A), one reader only. Consequence: additive rankings are \emph{not} instances of (CP-s) or of (CP-b) with $\mathcal D=\mathcal D^+$ that are compatible with (S), so for these variants ``impossibility contains the additive impossibility'' fails; (CP-a) is the only Walley variant for which the reduction is proved and known to be non-vacuous. For general $\M_i$ the inconsistency is \src{N13} Conj.~N13.A (open).}

\subsection{The Walley-2000 variant}\label{sec:cp-b}
Under (CP-b) put $\M_i\coloneqq\{P:P\ge0\text{ on }\mathcal D_i\}$; then (IB) and strict (G) are automatic and every $P\in\M_i$ satisfies (W) (\src{N11} Prop.~6.3, refereed in \src{N13}). If countably many countably additive members strictly separate $\mathcal D_i$, the conclusions of Corollary~\ref{cor:DCA} hold (\src{N11} Thm.~6.4, refereed in \src{N13}); with strongly non-atomic members instead, \src{N11} Thm.~6.4 gives the finitely additive reduction of \src{N11} Thm.~5.5. If $\mathcal D_i$ contains \emph{invisible} desirable gambles ($P(f)=0$ for all $P\in\M_i$), the ranking is genuinely non-Archimedean and no real additive field obeying (W) detects it. Under a lexicographic representation with countably additive levels the interior lexicographic Bayes models are unique, coherent and selected (\src{N11} Prop.~6.5), but whether $T_1$ is lexicographically Bayes at some non-standard interior prior is the open lexicographic separation problem (\src{R11} \S5, \src{N10} on \src{N9.9}).
\status{the variant with invisible desirability is open.}

\subsection{Rankings that add up: the supporting assessment and its exceptional set}\label{sec:M0}

This subsection and the next two concern an \emph{arbitrary} ranking that adds up, in the sense of Definition~\ref{def:addup}; the ranking need not be additive, complete or a coherent almost-preference. Standing hypotheses \textbf{(H$_{\mathrm{M0}}$)}: $n=3$; the $U$-domain $\M_U$ for $T_1$; a ranking with strict (G), (IB), (LA), (S), (C) and $T_1\in\C_L$ ((LF) is automatic on $\M_U$); a selection $p\mapsto C_p$ as in (S); and (M0). Let $\psi$ be a supporting additive assessment, put $c_i(g)\coloneqq\rho_i(g)\operatorname{sign}(g_i)$, and write
\[
B_q\coloneqq\{g\in U:q\cdot c(g)>0\},\quad Z_q\coloneqq\{g\in U:q\cdot c(g)=0\},\quad E(\psi)\coloneqq\{q\in\Di:\mu(Z_q)>0\}.
\]
$B_q$ is the strict $\psi$-$q$-Bayes set, $Z_q$ the tie set, and $E(\psi)$ the set of \emph{exceptional weightings}. The $\psi$-$q$-Bayes models are the models $B_q\cup S$ with $S\subseteq Z_q$ measurable. Let $\mathcal R\coloneqq\{\psi(X):X\in\M_U\}$ (compact and convex by Lyapunov), $V_\psi(t)\coloneqq\min_{y\in\mathcal R}t\cdot y$, and let $\mathfrak F_t\coloneqq\operatorname{argmin}_{\mathcal R}t\cdot y$ be the \emph{face} at $t$.

\begin{proposition}\label{prop:M0}
Under (H$_{\mathrm{M0}}$), $\psi$ is a strictly guiding additive ranking on the same domain $\M_U$, and:
\begin{enumerate}[leftmargin=*,label=($\alph*$)]
\item[($\alpha$)] $T_1$ is $\psi$-Bayes at some interior weighting $p^*$;
\item[($\beta$)] at \emph{every} weighting $p\in\Delta$ at least one $\psi$-$p$-Bayes model is coherent;
\item[($\gamma$)] $E(\psi)$ is Lebesgue-null, and at every $q\in\Di\setminus E(\psi)$ the $\psi$-$q$-Bayes model is a.e.\ unique and coherent, $B_q=D_{K_q}$;
\item[($\delta$)] the vertex Bayes models are coherent.
\end{enumerate}
If $E(\psi)=\varnothing$, then $c$ is a realisation of $T_1$ at $p^*$, which Theorem~\ref{thm:main} excludes. So a cover forces $E(\psi)\ne\varnothing$.
\end{proposition}
\status{proved. For countably additive weights: \src{N17} Prop.~2.1, refereed CORRECT in \src{N19}; for general weights: via \src{N13} Prop.~B (proved in that referee report; its ($\beta$) re-derived in \src{N24} \S2.1); the restriction to relevant comparisons is \src{N19} Prop.~R3(a) (proved in that referee report, one reader). The citation follows \src{N25\_referee} \S3, which notes that \src{N22} Thm.~2.1 is stated for (CP-a) and that the M0 form is \src{N19} \S2.2 with \src{N13} Prop.~B. The last sentence is the argument of \src{N22} Cor.~2.3 (refereed in \src{N24} \S2.2).}

So the unresolved issue in (M0) is not that coherence disappears at exceptional weightings. At $t\in E(\psi)$ coherent and incoherent minimisers \emph{coexist} (Theorem~\ref{thm:Estructure}(d)); $\psi$ assigns them equal profiles and cannot tell them apart; by Theorem~\ref{thm:kkm} and (C) every undominated model is coherent, so the original ranking must tell them apart, by strict comparisons that $\psi$ does not register.

\begin{theorem}\label{thm:Estructure}
Let $c$ be a strictly signed integrable field on $\M_U$ whose Bayes models are coherent at a.e.\ interior weighting (property ($\gamma$); for instance a supporting assessment under (H$_{\mathrm{M0}}$)). Then:
\begin{enumerate}[leftmargin=*,label=(\alph*)]
\item $V_\psi$ is concave on $\Delta$; at $t\in\Di$ it is differentiable iff $t\notin E(\psi)$, and its superdifferential at $t$ is $\mathfrak F_t$.
\item $E(\psi)=\bigcup_k(\ell_k\cap\Di)$, a countable (possibly empty) union of \emph{atom chords}. For each $k$ there is an \emph{oriented atom} $G_k$: a non-null set of mixed gambles $g\in U$ whose vectors $c(g)$ are all positive multiples of one vector $b_k$; then $\ell_k=\{q:q\cdot b_k=0\}$ is their common chart line. Each $G_k$ lies in one open orthant, $b_k$ has no zero coordinate, and $\ell_k$ misses the vertices of $\Delta$ and cuts off exactly one of them, that of the \emph{odd world} of $b_k$ (the world whose sign differs from the other two). There are no \emph{fan points}: no interior weighting at which a non-null set of gambles whose chart directions are spread without atoms has its chart lines through that weighting.
\item $E(\psi)$ meets every line other than the $\ell_k$ in a countable set; if it contains a segment, it contains the whole open chord.
\item For every $q\in\Di$, $B_q$ is coherent, and it is the coherent $q$-Bayes model with the largest credal set, $B_q=D_{K^\flat_q}$. The map $q\mapsto K^\flat_q$ is upper semicontinuous on $\Di$, equals $K_q$ off $E(\psi)$, and is continuous exactly at the points of $\Di\setminus E(\psi)$. Every extreme point of $\mathfrak F_q$ has a coherent realiser, while models $B_q\cup S$ with $S\subseteq Z_q$ bounded and non-null are incoherent.
\item For an oriented atom $G$ with line $\ell$: at every $q\notin E(\psi)$ on the positive side of $\ell$ the Bayes model contains $G$, and on the negative side it misses $G$. At all but countably many points $t$ of the chord, $B_t$ and $B_t\cup G$ are coherent, so $G$ is a coherent conic shell between two coherent cones.
\end{enumerate}
\end{theorem}
\begin{proof}[Proof sketch]
(a) $V_\psi(t)=\int_U\sum_{i:g_i>0}t_i\rho_i\,d\mu-\int_U[t\cdot c]_+\,d\mu$, and $\Phi'(t;\delta)+\Phi'(t;-\delta)=\int_{Z_t}|\delta\cdot c|\,d\mu$ for the convex part $\Phi$; $c\propto\one$ is impossible on $Z_t$. (b) With $\gamma(g)\coloneqq c(g)/|c(g)|$ and $m\coloneqq\gamma_*\mu$ on the sphere, $\mu(Z_q)=m(q^\perp\cap S^2)$; the atoms of $m$ give the chords and the non-atomic part gives the fan points, of which there are countably many since two great circles meet in two points. \emph{No fans:} along each line $g+\R\one$ the tied good gambles share one oriented chart line through the fan point, and $P(q,\cdot)$ is $1$-Lipschitz in the sup norm; two nearby fibres with different lines then give $|J(g)|+|J(g')|\le2\operatorname{dist}$ for the lengths of their tied intervals, which fails at a density point. (d) Directional limits $B_q\cup(Z_q\cap\{\delta\cdot c>0\})$ are limits in measure of coherent Bayes models, hence coherent; $B_q$ is a countable intersection of them; Straszewicz's theorem for extreme points. (e) For $g\in G$ the sign of $q\cdot c(g)$ is that of $q\cdot b$.
\end{proof}
\status{(a), (c), (d), (e) and the chord description in (b): proved (\src{N25} Thm.~2.1, Cor.~2.2, Thm.~3.3, Thm.~4.1, the claim referred to as N25 Claim 1); confirmed by one referee (\src{N25\_referee} \S\S1--2). ``No fan points'': proved in a referee report (\src{N25\_referee} \S2, Prop.~F), checked by nobody else. ``Cuts off one vertex'' is the referee's observation that it follows from \src{N25} Thm.~4.1(a), since $b_k$ has mixed signs and no zero coordinate. Incoherence of bounded non-null splices in (d) is the splice fact of \src{R2} Lemma U (two cones cannot differ by a bounded non-null set). Prop.~F uses only ($\gamma$): no (E$_\psi$), no (S), no lexicographic structure. The \src{N26\_referee} (\S3.3) records fan points as open for general M0 rankings and exhibits a strictly signed field with a fan point; that field does not satisfy ($\gamma$), so it does not conflict with Prop.~F, which the \src{N26} referee did not have. Nothing bounds the number of chords, and nothing prevents $E(\psi)$ from being dense. Crossings of two atom chords (with a two-dimensional face) are not excluded in general.}

\begin{theorem}\label{thm:pstarE}
Assume ($\alpha$) and ($\gamma$), and suppose $p^*\in E(\psi)$. Then there are a direction $\delta$, a sequence $s_m\downarrow0$ with $p^*+s_m\delta\notin E(\psi)$, and a relatively open set of gambles $h\in\partial D_{T_1}\setminus0$ with $P_{T_1}(h)=0$ but $P(p^*+s_m\delta,h)\to P_{K_\delta}(h)\ne0$, where $D_{K_\delta}$ is the directional limit $B_{p^*}\cup(Z_{p^*}\cap\{\delta\cdot c>0\})$. So $P(\cdot,h)$ does not converge to $0$ at $p^*$, and the chain of \S\ref{sec:main} cannot start at $p^*$: Lemma~\ref{lem:levels} at level $0$ through $p^*$ and Step~1 of Theorem~\ref{thm:pivot} both need that convergence. If $p^*$ is a generic point of an atom chord of sign pattern $\pm\sigma$, the affected $h$ have pattern $\sigma$ or $-\sigma$: they lie on one middle arc of a facet and near the opposite corner ray.
\end{theorem}
\status{proved (\src{N25} Thm.~9.1); confirmed with repair (\src{N25\_referee} \S4). Repair: the note's ``$q\mapsto P(q,h)$ has no limit at $p^*$'' is not proved; what is proved is non-convergence to $P_{T_1}(h)=0$, which is all the blocking conclusion needs. Addition (CHECKED, \src{N25\_referee} \S4, transferring \src{N26} Prop.~4.3): if the tie set at $p^*$ is a single oriented atom $G$ and no other chord passes through $p^*$, then $T_1=B_{p^*}\cup G$ is the accepting extreme of the face and $\sigma$ is the sign pattern of a corner ray.}

Let $\mathcal N$ be the convex cone of supporting assessments (the countably additive parts of endorsing weights), and let the \emph{atom set} $A(\psi)$ be the union of the oriented atoms of $\psi$.

\begin{proposition}\label{prop:generic}
Under (H$_{\mathrm{M0}}$):
\begin{enumerate}[leftmargin=*,label=(\alph*)]
\item some $\psi^*\in\mathcal N$ has an atom set of least measure, and $A(\psi^*)$ is contained, modulo null sets, in $A(\psi')$ for every $\psi'\in\mathcal N$;
\item for $\kappa\in(0,\infty)^3$, $\kappa\odot\psi\in\mathcal N$, and $E(\kappa\odot\psi)$ and the target weighting are the images of $E(\psi)$ and $p^*$ under one projective bijection of $\Delta$ fixing the vertices; so ``$p^*\in E$'' is invariant under rescaling;
\item if the ranking covers $T_1$, then every $\psi\in\mathcal N$ has $E(\psi)\ne\varnothing$, hence an atom chord, and $\mu(A(\psi))>0$. Genericity cannot remove the chords.
\end{enumerate}
\end{proposition}
\status{(a): proved (\src{N25} Thm.~6.2); confirmed with repair (\src{N25\_referee} \S1): normalise the weights, $\|\nu_k\|\le1$, rather than $\int|c_k|\le1$, so that the countable mixtures $\sum_kw_k\nu_k$ converge and have countably additive part $\psi^w$. (b): proved (\src{N25} Prop.~6.4(a)); confirmed. (c): proved in a referee report (\src{N25\_referee} \S3), one reader, from Proposition~\ref{prop:M0} and Theorem~\ref{thm:Estructure}; it needs no (E$_\psi$).}

\paragraph{The hypothesis (E$_\psi$).} Say that \textbf{(E$_\psi$)} holds if $\psi$ itself endorses the relevant comparisons: $B\in\C_L$, $X$ regular and $B\preceq_iX$ imply $\psi_i(B)\le\psi_i(X)$. It holds when the endorsing weights are countably additive (then $\psi_i=\nu_i\circ E_i$). It is not supplied by \src{N22} or \src{N24} in general; with purely finitely additive parts, endorsement gives only $\psi_i(C)+\nu^{\mathrm s}_i(E_iC)\le\psi_i(Y)+\nu^{\mathrm s}_i(E_iY)$ with uncontrolled singular parts $\nu^{\mathrm s}_i$ (\src{N25} N25.12, \src{N25\_referee} \S3). Call a selected $C_p$ \emph{$j$-improvable} if some regular $Z$ has $Z\prec_jC_p$ and $\psi_i(Z)\le\psi_i(C_p)$ for both $i\ne j$. On a chord $\ell$, let $\ell^\circ$ be its points on no other chord, so that $Z_t=G\cup G'$ for $t\in\ell^\circ$ ($G,G'$ the two oppositely oriented atoms on $\ell$); a \emph{split} is a measurable $S\subseteq G\cup G'$.

\begin{proposition}\label{prop:Epsi}
Assume (H$_{\mathrm{M0}}$) and (E$_\psi$).
\begin{enumerate}[leftmargin=*,label=(\alph*)]
\item For $t\in\Di$, every point of $\mathfrak F_t$ is the profile of a selected $t$-Bayes model; every $t$-Bayes $Y$ is weakly dominated by a selected $C$ with $\psi(C)=\psi(Y)$, strictly at some world if $Y$ is incoherent.
\item $q\mapsto\psi(C_q)$ is continuous at every $p$ whose selected credal set lies in $\Di$.
\item For each world $j$, the projected profiles $(\psi_i(C_p))_{i\ne j}$ of $j$-improvable selected models lie in countably many graphs of $1$-Lipschitz functions over the line $\one^\perp\subseteq\R^2$; in particular they form a null set.
\item Hence at a two-dimensional face (which by Theorem~\ref{thm:Estructure} occurs only where atom chords cross), for area-a.e.\ tied profile no selected realiser is $j$-improvable for any $j$ (``ideal tie-breaking''); on an atom face, all but countably many tied profiles have no selected realiser improvable at the odd world.
\item Along a chord, (IB) makes the comparisons among the models $B_t\cup S$ the same at every $t\in\ell^\circ$; the undominated splits do not depend on $t$, and for each undominated split $S$ the model $B_t\cup S$ is coherent at every $t\in\ell^\circ$.
\end{enumerate}
\end{proposition}
\status{proved (\src{N25} Lemma 1.2, Lemmas 7.1--7.4, Thm.~7.5, Cor.~7.6, Lemma 8.1, Thm.~8.2); confirmed by one referee (\src{N25\_referee} \S\S1, 4). (E$_\psi$) is genuinely needed: with charges, the tail estimate behind (b) fails (the obstruction of \src{N22} Remarks 3.3). (c) is the analogue, without lexicographic structure, of \src{R2} Thm.~7; it costs one dimension.}

\subsection{Residual configurations of (M0)}\label{sec:M0-residual}
Under (H$_{\mathrm{M0}}$), without (E$_\psi$), a cover of $T_1$ forces, for every $\psi\in\mathcal N$: $E(\psi)$ is a nonempty countable union of open atom chords, there are no fan points, and each chord cuts off the vertex of its odd world. With $\psi^*$ as in Proposition~\ref{prop:generic}(a), one of the following holds (\src{N25\_referee} \S7, superseding \src{N25} \S10).
\begin{itemize}[leftmargin=*]
\item[\textbf{R-a$'$}] $p^*\in E(\psi^*)$: the target weighting lies on an atom chord. Sub-cases: (i) a single oriented atom at $p^*$; (ii) both oriented atoms of the chord non-null; (iii) $p^*$ at a crossing of two or more chords, with a two-dimensional face. \emph{Status:} open. The chain of \S\ref{sec:main} cannot start (Theorem~\ref{thm:pstarE}). In (i), $T_1$ is the accepting extreme and $\sigma$ a corner pattern. Under (E$_\psi$), Proposition~\ref{prop:Epsi} applies.
\item[\textbf{R-b$'$}] $p^*\notin E(\psi^*)$, and $E(\psi^*)$ is a nonempty countable union of atom chords, possibly crossing and possibly dense. \emph{Status:} open; it needs a version of \S\ref{sec:main} that tolerates chords.
\end{itemize}
Closed: every configuration with fan points (\src{N25}'s residual R-c and the fan half of R-a), by Prop.~F; and \src{N25} Conj.~N25.L, which is vacuous because a countable $E(\psi)$ is empty.

\emph{Localisation.} The proof of Theorem~\ref{thm:main} as written uses Hausdorff continuity of $K$ on sets that together cover $\Di$ (\src{N25} Prop.~9.3, CHECKED, confirmed as a reading of the proof). The edge box (Lemma~\ref{lem:edgebox}) holds for every strictly signed integral field whose Bayes models are coherent at a.e.\ weighting (\src{N25} Prop.~9.2, proved, confirmed; also \src{N26} \S4.3). Lemma~\ref{lem:chords} and Lemma~\ref{lem:levels} survive off a null $E$ (\src{N26\_referee} \S4.2, proved there for the class $\Ltwo$; on our reading the argument uses only continuity of $K$ on $\Di\setminus E$ and nullity of $E$, and \src{N25\_referee} \S5 makes the same observation for straight levels). What genuinely crosses an atom chord is the clopen argument of Lemma~\ref{lem:PL2} on $\operatorname{int}\co(G\cup V_l)$: every chord disconnects $\Di$ and separates $p^*$ from some vertex. The remaining failure point is evaluation at $p^*$ when $p^*\in E$.

\subsection{Resolving exceptional ties: two worlds, and the lexicographic class}\label{sec:M0-lex}

Let $\Ltwo$ be the class of \emph{two-level lexicographic rankings with additive levels} on $\M_U$: world $i$ ranks $X$ by the lexicographic order of $(\psi_i(X),S_i(X))$, where the \emph{primary} $\psi_i(X)=\int_{E_i(X)}\rho_i\,d\mu$ has $\rho_i>0$ a.e., and the \emph{secondary} $S_i(X)=\int\sigma_i(g,a_X(g))\,d\mu(g)$ has an integrable envelope and is otherwise arbitrary. The primary endorses \emph{every} comparison at world $i$, so every world adds up, and (E$_\psi$) holds with $\psi$ the primary. A model is \emph{improvable at $i$} if some $Z$ has $\psi_i(Z)=\psi_i(X)$ and $S_i(Z)<S_i(X)$.

\begin{theorem}\label{thm:N26.1}
Let $n=2$ and let $\M_U$ be the $U$-domain with $K'=\{\tfrac12\}$. There is an explicit ranking in $\Ltwo$ that satisfies strict (G), (IB), (LA), (S) (with an explicit selection and closed sections for every model), (C), and (LF) at the target $T\coloneqq D_K$, where $K$ is the set of probabilities whose world-$1$ coordinate lies in $[0.25,0.625]$, such that:
\begin{enumerate}[leftmargin=*,label=(\alph*)]
\item the primary has exactly one exceptional interior weighting $t_0=0.3$, whose tie set is a non-null cone of mixed problems;
\item $T$ is a coherent primary-minimiser at $t_0$ that is neither one-sided limit of the Bayes models (it is intermediate);
\item every incoherent $t_0$-minimiser is strictly dominated, at both worlds, by the coherent model with the same primary profile;
\item the undominated models form a continuous monotone path of coherent models, and the selection index is not the primary's Bayes weighting.
\end{enumerate}
\end{theorem}
\begin{proof}[Proof sketch]
For a mixed pair put $r\coloneqq-g_2/(g_1-g_2)$, which is the same for $g$ and $-g$. Pairs with $r<\tfrac12$ have their $U$-orientation in $H$; pairs with $r>\tfrac12$ have $U$-orientation $-g$, with negative first coordinate, and evaluable models may accept $-g$ there. Take $\rho_1=1-\theta(r)$, $\rho_2=\theta(r)$ on $U$, with $\theta$ nondecreasing on $r<\tfrac12$ and equal to $t_0$ on $[0.2,0.3]$, and strictly increasing on $r>\tfrac12$. The secondaries are built from the path $u\mapsto C_u$ of undominated models so that each $C_u$ is the a.e.-unique pointwise minimiser of $S_1+u\psi_1$ and of $S_2+(3-u)\psi_2$. Then no selected model is improvable at either world, the sections of (S) are primary sublevel sets of a continuous monotone function of $u$, and at $t_0$ a bathtub comparison shows that every incoherent tied model loses strictly at both worlds to the coherent model of equal primary mass.
\end{proof}
\status{proved (\src{N26} \S1, N26.1); confirmed by one referee (\src{N26\_referee} \S1: every axiom re-derived by hand, with independent numerics). It is a genuine M0 mechanism: a ranking that adds up, satisfies every axiom, and resolves an exceptional additive tie that its supporting assessment cannot. In two worlds every target is already covered (\src{Full Coverage \dots v3}), so it adds no coverage. Referee repair to an unused remark (CHECKED): other assessments that endorse only the relevant comparisons can break the tie.}

\begin{theorem}\label{thm:Ltwo}
Let $n=3$, let a ranking in $\Ltwo$ on $\M_U$ satisfy (LA), (S), (C), and let $\psi$ be its primary.
\begin{enumerate}[leftmargin=*,label=(\alph*)]
\item \emph{Multiplier normal form.} At a world $i$ and a level $x$ interior to the range of $\psi_i$, put $g_i(x)\coloneqq\min\{S_i(Z):\psi_i(Z)=x\}$. Every undominated model at level $x$ that is not improvable at $i$ minimises $S_i+\lambda\psi_i$ pointwise for every $\lambda\in-\partial g_i(x)$. If two such models at one level differ on a non-null set, the multiplier at that level is unique; the set of multiplier \emph{values} arising at such levels is countable.
\item \emph{Splice lemma.} If two distinct undominated $t_0$-Bayes models ($t_0\in\Di$) share a multiplier at every world, a bounded splice of them is undominated and incoherent, so (LA), (S), (C) fail.
\item Every interior face of the primary has dimension at most $1$. Hence $E(\psi)$ is a countable union of non-crossing chords; the tie set of a chord $\ell$ is the union of two oppositely oriented atoms, of sign patterns $\sigma$ and $-\sigma$, with chart line $\ell$.
\item \emph{Rigidity.} No chart line of a gamble outside the tie set of $\ell$ crosses $\ell$ inside $\Delta$; the faces at the points of $\ell$ are translates of one fixed segment; $\ell$ cuts off a vertex, and the credal sets of the face models meet both edges at its endpoints.
\item If $T=D_K\in\C_L$ with $K\subseteq\Di$, every interior weighting $p$ at which $T$ is $\psi$-Bayes lies outside $E(\psi)$, and $\operatorname{dist}(p,E(\psi))>0$ whenever $K_q\to K$ as $q\to p$.
\end{enumerate}
Consequently, a ranking in $\Ltwo$ on $\M_U$ that covers $T_1$ has $E(\psi)$ a nonempty countable union of rigid, vertex-cutting, non-crossing chords at positive distance from $p^*$.
\end{theorem}
\status{proved (\src{N26} Lemmas 2.1--2.3, Thm.~3.1, Cor.~3.2, Thm.~3.3, Cor.~3.4, Thm.~4.1, Cor.~4.2, i.e.\ N26.2--N26.7); confirmed with repairs by one referee (\src{N26\_referee} \S\S2--3). Repairs incorporated: in (a), what is countable is the set of multiplier values, not the set of levels carrying a value; in (c), the tie set has the two patterns $\pm\sigma$, not one; in (d), the face is a translate of a fixed segment in direction $d$, not $[y^-,y^-+D]$. \emph{Scope.} These results concern the class $\Ltwo$ with $\psi$ its primary. They do \emph{not} transfer to general M0 rankings: the proofs use the lexicographic structure at four points (undominated implies weakly $\psi$-Pareto; \src{R2} Thm.~7; a multiplier from the convex value function of an additive secondary; the splice is undominated because it is a pointwise Lagrangian minimiser) (\src{N26\_referee} \S3.2, \src{N25\_referee} \S6). For general M0 rankings, ``no isolated exceptional points'' holds by Prop.~F, but ``no two-dimensional faces'' and ``$p^*\notin E$'' are not established. Theorem~\ref{thm:N26.1} shows that (e) is genuinely a three-world fact. Within $\Ltwo$ one may take $\psi$ to be the primary, which lies in $\mathcal N$ and satisfies (E$_\psi$); then R-a$'$ is excluded and only R-b$'$-type configurations remain (\src{N26\_referee} \S3.3).}

\begin{proposition}\label{prop:R1N26}
Let a ranking in $\Ltwo$ on $\M_U$ satisfy (LA), (S), (C), with $T_1\in\C_L$ primary-Bayes at $p^*$. Suppose that the cluster-line conclusion of Theorem~\ref{thm:pivot} holds for the primary's field on the middle range $[\tau_i,\tau_j]$ of each facet: every $h(\tau)$ has a unique oriented cluster line $\lambda_\tau\ni p^*$, $\tau\mapsto\lambda_\tau$ is continuous, and $\lambda_{\tau_i}=\lin(p^*,e_i)$, $\lambda_{\tau_j}=\lin(p^*,e_j)$. Then $E(\psi)=\varnothing$; with Theorem~\ref{thm:main}, no ranking in $\Ltwo$ on $\M_U$ covers $T_1$.
\end{proposition}
\begin{proof}[Proof sketch]
The lines $\lambda_\tau$ sweep every direction through $p^*$, so one of them crosses a given chord transversally; gambles near that $h(\tau)$ have chart lines crossing the chord, against Theorem~\ref{thm:Ltwo}(d).
\end{proof}
\status{proved in a referee report (\src{N26\_referee} \S4.3, Prop.~R1), one reader only.}

\begin{conjecture}\label{conj:N26R}
In $\Ltwo$ on $\M_U$, the cluster-line conclusion of Theorem~\ref{thm:pivot} holds for the primary off a countable union of rigid chords at positive distance from $p^*$. Sharper: a cover of $T_1$ in $\Ltwo$ would need rigid chords whose tie sets accumulate in measure at every point of some crossing ray of $\partial D_{T_1}$ (every ball around such a point meets some chord's tie set in a non-null set).
\end{conjecture}
\status{open (\src{N26} Conj.~N26.R, sharpened in \src{N26\_referee} \S4). By Proposition~\ref{prop:R1N26} the first sentence implies that no $\Ltwo$ ranking on $\M_U$ covers $T_1$; Theorem~\ref{thm:nosliding}, Corollary~\ref{cor:PCNS} and Proposition~\ref{prop:boxbridge} need no re-check. The second sentence is the referee's sharpened residual N26.R$'$: by Remark~R2 of that report (CHECKED, sketch), Lemma~\ref{lem:PL2} and Lemma~\ref{lem:nopin}, and with them Theorem~\ref{thm:pivot}, survive whenever no chord's tie set meets a small ball around a point of a crossing ray.}

\section{What is excluded and what is open}\label{sec:open}

\subsection{Two concrete open questions}\label{sec:two-questions}
An ordinary additive ranking cannot vindicate $T_1$ on the $U$-domain (Corollary~\ref{cor:coverage}). A more general ranking on that domain could escape this result in either of two ways.
\begin{itemize}[leftmargin=*]
\item \emph{One} (M0): positive error weights endorse its relevant comparisons, but the supporting additive assessment loses distinctions at exceptional weightings, where it ties coherent models with incoherent ones. The original ranking would have to distinguish them in a way compatible with (IB) and closed choice.
\item \emph{Two} (M1 or M2): no positive weighting of errors endorses all the relevant comparisons. This may have a finite witness (M1), or emerge only from the whole collection of comparisons (M2).
\end{itemize}
Neither escape has been ruled out or turned into a covering construction. General two-sided domains (\S\ref{sec:general}) are a separate open route. Two questions are concrete:
\begin{enumerate}[leftmargin=*,label=\textbf{(Q\arabic*)}]
\item \emph{Can exceptional additive ties be resolved ordinally?} On $\M_U$, can a ranking with strict (G), (IB), (LA), (S), (C) that adds up at every world cover $T_1$, by resolving the exceptional ties of its supporting assessments (the atom chords of Theorem~\ref{thm:Estructure}) in configuration R-a$'$ or R-b$'$ (\S\ref{sec:M0-residual})? Two things are \emph{not} to be demanded: that the supporting assessment itself satisfy (S) or (C), and that its Bayes weightings equal the selection indices (an auxiliary numerical representation can fail to supply the obvious selection without the original ranking failing closed choice; compare the conditional-probability example of \src{Loss Dominance Generalized v2}). A negative answer would prove that closed choice and independence prevent the needed distinctions, or permit supporting weights without exceptional ties. A positive answer would exhibit a concrete mechanism by which ordinal rankings outperform additive ones. Theorem~\ref{thm:N26.1} supplies such a mechanism in two worlds; in three worlds, Theorem~\ref{thm:Ltwo} and Conjecture~\ref{conj:N26R} describe what is left for lexicographic rankings. For coherent almost-preference rankings the same question is (U) (Conjecture~\ref{conj:U}, Remark~\ref{rem:EPchords}).
\item \emph{Can additive rankings on two-sided domains cover $T_1$?} Equivalently (Proposition~\ref{prop:CF}), does CF$(T_1)$ have a solution? The open case is sector caps at the failing edges glued to a non-precise layer along the crossing chords (\S\ref{sec:general-open}, Conjecture~\ref{conj:N21}).
\end{enumerate}
The present results do not force exact coverage to be replaced by approximation.

\subsection{Weakened guidance}\label{sec:open-ties}
Two target-independent weakenings of (G) were tested, both allowing good advice and silence to tie at world $i$ on mixed problems with small world-$i$ stake, $|g_i|\le\theta\|g\|$ (the band).
\begin{itemize}[leftmargin=*]
\item \emph{Forced ties} (\src{BRIEF3}). For additive rankings on decomposable domains, ``NSI$_\theta$ covers $T_1$'' fails for every $\theta>0$, for both the sup and the Euclidean norm, every $\varepsilon\in(0,1)$, and indeed for $\varepsilon$-contamination of every interior baseline in every $n\ge3$. \emph{Status:} proved, with two independent complete proofs (\src{N1}, \src{N2}); refereed once (\src{N4}). Scope: additive rankings only; forced-tie reading only.
\item \emph{Permitted ties} (\src{BRIEF4}; ties allowed, placed by the ranking). The conditional no-go of \src{R15} Thm.~6.2 (under the first-order hypotheses (DF), (ND) and non-adjacency) transfers verbatim for every $\theta>0$. \emph{Status:} proved; refereed once (\src{N7}, on \src{N5}, \src{N6}). The unconditional chain of \S\ref{sec:main} also transfers, below an explicit threshold (Theorem~\ref{thm:permitted}).
\end{itemize}
In both readings (C) forces (BU) at every interior weighting for every $\theta$ and both norms, because \src{R1} Lemma 1.4 uses no guidance; the caveat ``$\theta<1$ (sup), $\theta<1/\sqrt3$ (Euclidean)'' of \src{BRIEF4}, \src{N5} and \src{N7} is unnecessary (\src{N27} Prop.~27.4, confirmed in \src{N27\_referee} \S3.1).

A \emph{permitted-tie realisation} of $T$ at $p^*$ on $\M_U$ is a measurable $c:U\to\R^3$ with $c_i(g)\ge0$ if $g$ is mixed with $0<g_i\le\theta\|g\|$, $c_i(g)>0$ if $g_i>\theta\|g\|$ or $g$ is unanimous, $c_i(g)<0$ if $g_i<0$, together with (BU) at every $q\in\Di$ and (TP). For a facet $F_k$ and a crossing coordinate $i\ne k$ of $F_k$ consider
\begin{itemize}[leftmargin=*]
\item[\textbf{(X$_{ik}$)}] a non-null set of $x\in U$ has $v_k\cdot x<0$ and $c_i(x)>0$;
\item[\textbf{(VP)}] for every $x\in\relint F_k$ and $l\ne k$ with $x_l>0$, $\liminf_{q\to e_l,\,q\in\Di}P(q,x)>0$ (vertex positivity),
\end{itemize}
and put $\theta_x(T)\coloneqq\min_k\min_{i}\sup_{h\in F_k\setminus0}h_i/\|h\|$, the minimum over facets $F_k$ and crossing coordinates $i\ne k$.

\begin{theorem}\label{thm:permitted}
Let $T\subseteq\Di$ be an all-facing triangle in general position, $p^*\in\Di$, $\theta>0$, and the norm sup or Euclidean. If (X$_{ik}$) holds for every facet $F_k$ and every crossing coordinate $i$ of $F_k$, or if (VP) holds, then $T$ has no permitted-tie realisation at $p^*$ on any $U$-domain $\M_U$. (X$_{ik}$) holds automatically when $\theta<\theta_x(T)$. Consequently no additive ranking with permitted ties, integrable envelopes, (LA), (S), (C) on $\M_U$ covers $T$ in that range. For $T_1$, at every $\varepsilon\in(0,1)$ and every interior target weighting:
\begin{enumerate}[leftmargin=*,label=(\alph*)]
\item sup norm: $\theta_x=1$, so the conclusion holds for every $\theta\in(0,1)$, exactly the range in which the band is a proper subset of the mixed problems;
\item Euclidean norm: $\theta_x=\sqrt{(\kappa^2+1)/(\kappa^2+2)}$ ($\approx0.837$, $0.928$, $0.999$ at $\varepsilon=0.05$, $0.3$, $0.9$), so the conclusion holds for every $\theta<\sqrt{(\kappa^2+1)/(\kappa^2+2)}$, and in particular for every $\theta<\sqrt{2/3}$ uniformly in $\varepsilon$.
\end{enumerate}
\end{theorem}
\begin{proof}[Proof sketch]
A covering ranking yields a permitted-tie realisation ((BU) from \src{R1} Lemma 1.4, $p^*$ interior by \src{N1} Prop.~2.3, which uses only strict badness). In the chain of \S\ref{sec:main} the only input sensitive to ties is Lemma~\ref{lem:UVN}(a) at relative-interior facet points: positivity near $e_l$ of a facet gamble with $x_l>0$. It is equivalent to (VP) at $x$, and \src{N5} Thm.~3.3 (facet clause) supplies it from (X$_{lk}$). Every other sign use is strict badness ($c_i<0$ where $g_i<0$, never tied) or holds with weak signs (Lemma~\ref{lem:PL2}, Step 4 of Theorem~\ref{thm:pivot}, Proposition~\ref{prop:boxbridge}(a)). The corner-ray gap found in \src{N7} is never used. The transposition of Proposition~\ref{prop:T1} keeps the band, since both norms are permutation invariant. For $T_1$ in the Euclidean norm the maximiser of $h_l/|h|$ on $F_k$ is the normalised projection of $e_l$ onto $v_k^\perp$, which lies in $\relint F_k$, and $\theta_x^2-\tfrac23=(\kappa^2-1)/(3(\kappa^2+2))>0$.
\end{proof}
\status{proved (\src{N27} Thm.~27.6 with Prop.~27.4, Lemmas 27.5, 27.5$'$); confirmed by one referee (\src{N27\_referee} \S\S3, 6; thresholds recomputed symbolically, numerics reproduced). It rests on the facet clause of \src{N5} Thm.~3.3 (refereed once, \src{N7}). The referee adds that the Euclidean maximiser maximises $x_l/|x|$ over the whole half-space $\{v_k\cdot x\le0\}$, so enlarging $U$ does not raise the threshold for $T_1$.}

\begin{corollary}\label{cor:permitted-nonadj}
For $\theta<\theta_x(T)$, no additive ranking with permitted ties covers $T$ on any decomposable domain with (LF) (arbitrary open $N$) that is non-adjacent near $T$.
\end{corollary}
\status{CHECKED (\src{N27} Cor.~27.8, through \src{N27} Lemma 27.7, proved: the set of gambles accepted at some interior weighting is a.e.\ a countable union of the cones $D_{K_q}$, so the set of acceptances given by Bayes models is conic without vertex facts); confirmed at CHECKED level, with two wording repairs, by one referee (\src{N27\_referee} \S4, which also supplied the half-plane form of non-adjacent charts for advice range $\{+,0,-\}$ under weak signs).}

\emph{What remains open.}
\begin{itemize}[leftmargin=*]
\item \emph{The Euclidean window $\theta\in[\theta_x(T_1),1)$.} There no gamble at all beyond a facet plane is off-band at the crossing coordinate, so (X$_{lk}$) holds exactly when the ranking leaves untied a non-null set of $l$-positive points of $U\cap\{v_k\cdot x<0\}$; it is not automatic, but it does not fail automatically either. The open case is the \emph{all-tied configuration}: for some crossing pair $(l,k)$, $c_l=0$ a.e.\ on $U\cap\{x_l>0,\ v_k\cdot x<0\}$. The missing input is (VP); equivalently, every Hausdorff cluster set $W$ of $K_q$ as $q\to e_l$ satisfies $W\subseteq\{m:m\cdot w_{lk}\ge0,\ m\cdot r>0\}$, with $w_{lk}$ the $l$-crossing of $F_k$ and $r$ the corner ray at the $l$-positive end of $F_k$ (CHECKED). A directional vertex limit $W'=\co(K'\cup\{e_l,v_k\})$ meets the local vertex requirements and violates (VP) on the $l$-positive arc of $F_k$ (CHECKED, local; this repairs \src{N27}'s $W=[e_l,v_k]$, which violates $K_q\supseteq K'$). \emph{Conjecture} (Euclidean norm, unsupported either way): no permitted-tie realisation of $T_1$ exists for $\theta\in[\theta_x(T_1),1)$. In the sup norm nothing is open for $\theta\in(0,1)$. The window needs $\theta>\sqrt{2/3}\approx0.816$, far from negligible stakes.
\item Not examined: decomposable domains that are not non-adjacent near the target, under permitted ties, and the charge version (Corollary~\ref{cor:charges}) under permitted ties.
\end{itemize}

\subsection{Lexicographic and other non-additive rankings}
\begin{itemize}[leftmargin=*]
\item Two-level lexicographic rankings with additive levels: interior lexicographic Bayes models are a.e.\ unique and coherent (\src{N9} N9.5, \src{R11} Prop.~LX). A face-continuum theorem kills every design in which a secondary trades off against another world's primary at a continuum of rates (\src{N9} N9.6). The strongest plateau design (lexicographic Bayes map constant on $\Di$) and the free-collar design both fail, the latter with an explicit incoherent witness (\src{N9} N9.7, N9.8). \emph{Status:} proved; refereed once (\src{N10}). Coverage of $T_1$ by general lexicographic rankings is open; the precise residual is in \src{N10} \S4 (F4). For the class $\Ltwo$ on $\M_U$ the exceptional set of the primary consists of rigid, vertex-cutting, non-crossing chords away from the target weighting, and the remaining question is Conjecture~\ref{conj:N26R} (\S\ref{sec:M0-lex}).
\item Arbitrary (IB) preorders: a long list of necessary structure is proved (vertex anchoring, interior target index, closed tie fibres, edge repulsion; \src{N8}, \src{N9}, refereed in \src{N10}), but none of it excludes coverage. By \src{N9} N9.9 any ranking-level impossibility contains the additive one, which \S\ref{sec:main} now supplies on non-adjacent domains. \emph{Rankings that add up} (M0) are analysed in \S\S\ref{sec:M0}--\ref{sec:M0-residual}: the residuals are R-a$'$ and R-b$'$, question (Q1). \emph{Status:} open.
\item \emph{Rankings that do not add up} ((M1) or (M2) of \S\ref{sec:addup-def}). On $\M_U$, a cover outside (M0) needs a world with no weight, nonnegative on the relevant comparisons, whose countably additive part is positive a.e.\ on mixed problems (condition $\neg$SP$^{\mathrm{reg}}$; \src{N17}, \src{N19}). This arises either through a finite obstruction (M1), whose certificates in the stated formulation use at least four comparisons with overlapping gains and crossing error sets, or with no finite obstruction at all (M2). An (M1) certificate need not be an executable sequence of losing trades, and the restrictions just listed do not constrain (M2). \src{N19} states the necessary condition as conditional on excluding the integral class of \src{N13} Prop.~B, i.e.\ integral fields with unique coherent Bayes models at \emph{almost every} weight. Theorem~\ref{thm:main} removes the regularity part of that condition; the ``almost every'' part is exactly the (M0) residual of \S\ref{sec:M0-residual} (as framed in \src{BRIEF8} and \src{N25} \S10), and for coherent almost-preferences it is (U) (an identification made here, not in the sources). No system is known, for any target, in which (M1) or (M2) occurs on the relevant comparisons together with (S). \emph{Status:} route open; \src{N19} Conj.~N17.B$'$.
\end{itemize}

\subsection{Other targets and more worlds}\label{sec:open-other}
\begin{itemize}[leftmargin=*]
\item \emph{Other triangles.} Theorem~\ref{thm:main} covers every all-facing triangle in general position. For the crossing classes $(1,2,3)$ and non-facing $(2,2,2)$ the analogue is \src{R15} Conjecture R15, supported by numerics on about 1340 random triangles ($\min\Lambda_{\max}=1.0045$ for $(1,2,3)$, $1.2187$ for $(2,2,2)$; corrected values in \src{R18} \S7). The first-order arguments of \S\S4.1--4.5 apply facet by facet wherever a facet has a crossing with a shared positive coordinate; the pinching algebra is specific to the all-facing class, and for $(1,2,3)$ a nondegeneracy at the one-crossing facet is needed (\src{R18} \S5). \emph{Status:} open.
\item \emph{Segments and boundary targets} ($T_3$, $T_4$): open. Several lemmas need $\dim K=2$ (\src{R12} \S1 repair; \src{R1} Lemma 3.5).
\item \emph{$n\ge4$.} Not addressed. Vertex confinement, Lemma UVN, Lemma M, \src{R10b} Thm.~PL and Cor.~PL2 hold in every $n$; Theorem~\ref{thm:ccibu}, the box/corner system and the pinching are three-world arguments. On all-model domains, v3's additive obstruction excludes full additive coverage of targets with three vertices of $\Delta$ (\src{Full Coverage \dots v3}, App.~C).
\end{itemize}

\section{Positive results and approximation}\label{sec:positive}
\begin{sloppypar}
\begin{itemize}[leftmargin=*]
\item \emph{Gauge mixtures.} For compact convex $M_i\subseteq\R^3_{\ge0}\setminus0$ with $a_ie_i\in M_i$ and $\nu(M_i)\subseteq\co(\nu(M_j)\cup\{e_i\})$ ($\nu$ = normalisation), the penalties built from $f_i(h)=\min_{m\in M_i}m\cdot h$ give an additive system on a decomposable domain whose Bayes model at every $p\in\Delta$ is unique and equals $D_{K_p}$, $K_p=\nu(\sum_ip_iM_i)$; every interior $K_{p^*}$ is covered with (LF) (\src{R10a} Thm.~M; refereed once, \src{R12}). Examples include a regular hexagon around $u$ and 9-gons converging to $T_1$ in the Hausdorff metric (\src{R10a} Prop.~M$'$; refereed once, \src{R12}). No gauge mixture covers a triangle in general position (\src{R10a} Thm.~R, ray lemma; refereed once).
\item \emph{Other covered classes.} Precise models; normalised boxes and their conditional closures (\src{Conditional Shunning of Error.tex}); every coherent target when $n=2$ (\src{Full Coverage \dots v3}); the vacuous model under its accepted exception.
\item \emph{Approximation.} For an interior triangle $L$ and its parallel bodies, $\lp_{L+rB}(g)=\lp_L(g)-r\|P_Hg\|$, so $D_{L+rB}\uparrow D_L$ as $r\downarrow0$; parallel bodies are covered (CPM manuscript v19, cone domains). This is a genuine approximation result, but the representable class is not closed in the target topology: flat facets appear only at $r=0$ (\src{CONTAMINATION\_NOGO\_SCOPE\_AND\_GEOMETRY.md} \S\S4--5).
\item \emph{Earlier no-go for continuous maps.} For continuous, positively homogeneous coherence-preserving maps with weak, generically strict coordinate signs, the three-state contamination triangle at any interior baseline has no realiser at any interior prior, by the same facet maps, corner equations and box as in Proposition~\ref{prop:T1} (\src{CONTAMINATION\_NOGO\_SCOPE\_AND\_GEOMETRY.md} \S2; earlier work, not re-refereed here). In the ranking setting, on domains non-adjacent near the target, Theorem~\ref{thm:main} needs no continuity assumption; on general decomposable domains the question is open and reduces to CF$(T_1)$ (\S\ref{sec:general-open}).
\end{itemize}
\end{sloppypar}

\begin{thebibliography}{9}
\bibitem{HU1977} F.~Hiai and H.~Umegaki, Integrals, conditional expectations, and martingales of multivalued functions, \emph{Journal of Multivariate Analysis} 7 (1977) 149--182.
\bibitem{W1991} P.~Walley, \emph{Statistical Reasoning with Imprecise Probabilities}, Chapman and Hall, London, 1991.
\bibitem{W2000} P.~Walley, Towards a unified theory of imprecise probability, \emph{International Journal of Approximate Reasoning} 24 (2000) 125--148.
\end{thebibliography}

\section*{Sources}\label{sec:sources}
\addcontentsline{toc}{section}{Sources}
Project notes, in the working folder unless stated. Short codes in the text refer to these files.
\begin{description}[leftmargin=*,font=\normalfont\ttfamily,style=nextline]
\item[BRIEF.md, BRIEF2.md, BRIEF3.md--BRIEF8.md] Setting, axioms, (LF), decomposable domains, the realisation problem, weakened guidance, (CP-a).
\item[Loss Dominance Generalized v2.tex] Rankings as preorders, competitiveness, (E), (LA), (A), (C), (S), KKM dominance theorem (Generalized dominance paper folder).
\item[Full Coverage by Homothetic Credal Families v3.tex] Binary problems, four-rank guidance order, (G), (IB); two-world coverage; all-model obstruction (same folder).
\item[Consolidated Statements - Decomposable Domains 2026-10-12.md] Status of R1--R9 (R1 table of charts for each advice range; chart theorem, continuity; R2 Lemma U; R4 Thms A, B, Prop.~G1).
\item[R1\_notes.md, R2\_notes.md, R4\_notes.md, R6\_referee.md, R8\_referee.md] Additive rankings on decomposable domains (charts, Bayes uniqueness, adjacency, ray equations); referee reports.
\item[R10a\_notes.md, R10b\_notes.md, R12\_referee.md] Gauge mixtures, ray lemma; vertex confinement, Lemma UVN, Lemma M, Thm.~PL, Cor.~PL2; referee report.
\item[R11\_notes.md, R13\_referee.md] Lemma IC, Lemma K, Thm.~W; referee report.
\item[R15\_notes.md, R18\_referee.md] Corner calculus, all-facing triangles, box and Thm.~6.2; referee report.
\item[CCIBU.tex] Closed choice from interior Bayes uniqueness in three worlds.
\item[N12\_notes.md, N14\_referee.md] Reduction to a credal family, chords, edge box, (PC)+(NS) $\Rightarrow$ contradiction, two-slope projectivity; referee report with Thm.~C2 and the non-adjacent reduction.
\item[N15\_notes.md, N18\_referee.md, N20\_referee.md] Straight levels, pinning, pivot chords, no sliding edges, slope--rotation inequalities, two-stage blow-up; two independent referee reports.
\item[N21\_notes.md, N23\_referee.md] General decomposable domains; referee report with Lemma R23.1 and Cors.~R23.2--R23.4.
\item[N11\_notes.md, N13\_referee.md, N22\_notes.md, N24\_referee.md] Coherent almost-preference rankings; referee reports with Props A, B, R1, R2.
\item[N1--N7, N8--N10, N16, N17\_notes.md, N19\_referee.md] Weakened guidance, non-additive constructions, adversarial attempts, rankings that do not add up.
\item[N25\_notes.md, N25\_referee.md] The exceptional-tie case M0: structure of $E(\psi)$, coherent envelopes, atoms, genericity, consequences of (E$_\psi$), the target at an exceptional weighting; referee report with Prop.~F (no fan points) and corrected residuals R-a$'$, R-b$'$.
\item[N26\_notes.md, N26\_referee.md] Resolving exceptional ties lexicographically: the two-world construction, the class $\Ltwo$ in three worlds; referee report with Prop.~R1 and the sharpened Conjecture N26.R.
\item[N27\_notes.md, N27\_referee.md] The compact-collar repair of $N_0$, ray transfer of charts, the main theorem under permitted ties; referee report.
\item[N28\_notes.md, N28\_referee.md] Competing acceptances on general (LF) decomposable domains: (TS) in every cover, cut apexes, vertex-free interiors, CF$(T)$, diagonal symmetries; referee report with Lemma R1, Prop.~R3 and the sector-cap example C1.
\item[START HERE -- Coverage Status.md] Glossary (terminology of \S\ref{sec:addup-def}).
\item[Auditor's note (14 Oct 2026)] Corrections adopted by the author: relevant comparisons, the content of the additive reduction, non-adjacency, ``supporting additive assessment'', and the (M1)/(M2) wording.
\item[w1991.txt] Extracts of Walley (1991) \S\S3.7--3.8 (lines 8200--8600).
\item[CONTAMINATION\_NOGO\_SCOPE\_AND\_GEOMETRY.md] Continuous CPM no-go for contamination triangles (CPMs folder; earlier work).
\end{description}

\end{document}
